% 《量子金融》精读导读（中文原创精讲，配合原书阅读；非翻译、不替代原书）
% 原书：Belal E. Baaquie, Quantum Finance: Path Integrals and Hamiltonians
%       for Options and Interest Rates, Cambridge University Press, 2004
% 编译：xelatex ×3
\documentclass[11pt,a4paper,twoside,openright]{ctexbook}
\usepackage{amsmath,amssymb}
\usepackage{mathrsfs}
\usepackage[margin=2.4cm]{geometry}
\usepackage{graphicx}
\usepackage{float}
\usepackage{caption}
\usepackage{needspace}
\usepackage{booktabs}
\usepackage{multirow}
\usepackage{longtable}
\usepackage{array}
\usepackage{enumitem}
\ctexset{
  chapter/name = {第,章},
  section/format = \Large\bfseries,
}
\allowdisplaybreaks
\raggedbottom
\clubpenalty=10000 \widowpenalty=10000 \predisplaypenalty=10000
\graphicspath{{figures/}}
\xeCJKDeclareCharClass{CJK}{"201C, "201D, "2018, "2019}

% 插入的空白衬页真正空白（无页眉页脚）
\makeatletter
\def\cleardoublepage{\clearpage\if@twoside\ifodd\c@page\else
\hbox{}\thispagestyle{empty}\newpage\if@twocolumn\hbox{}\newpage\fi\fi\fi}
\makeatother

% 指向原书图表的行内导航标记
\newcommand{\origfig}[2]{\paragraph{原书图 #1}\quad #2\par}
\newcommand{\origtab}[2]{\paragraph{原书表 #1}\quad #2\par}

% 长表（术语表）列头
\newcolumntype{L}[1]{>{\raggedright\arraybackslash}p{#1}}

\begin{document}
\frontmatter
\begin{titlepage}
\centering
{\Huge\bfseries 量子金融\\[0.8em]}
{\LARGE 期权与利率的路径积分与哈密顿量\\[0.4em]}
{\LARGE 精读导读\\[5em]}
{\Large Quantum Finance:\\
Path Integrals and Hamiltonians for Options and Interest Rates\\[3em]}
{\Large Belal E. Baaquie（剑桥大学出版社，2004）原著\\[1.5em]}
{\large 中文精读导读\\[0.5em]}
{\large 供持有原书的读者对照阅读\\[0.3em]
含全书术语中英对照表与记号总表\\[3em]}
{\large 仅供个人学习研究使用\\[2em]}
\end{titlepage}
\thispagestyle{empty}\mbox{}\clearpage

% ch00 导读说明 (原创)
\chapter*{导读说明}
\addcontentsline{toc}{chapter}{导读说明}

本册是 Belal E. Baaquie所著 \emph{Quantum Finance: Path Integrals and
Hamiltonians for Options and Interest Rates}（Cambridge University Press,
2004）的\textbf{中文精读导读}，供持有原书的读者对照阅读。它\textbf{不是译本}：
全书以讲解体重新组织——每节先说明"这一节要解决什么问题"，再按逻辑顺序
呈现模型设定与关键公式（编号与原书一致，形如 (N.M)），最后用自己的话
解释推导主线与金融/物理含义。原书的完整行文、例题展开与文献细节仍以
原书为准；本册的作用是把全书的数学骨架提出来，让读者知道每一步"为什么
这么走"，并在回头精读原书时有所依凭。

\section*{全书结构一览}

原书论证的主线是一条不断"升级物理工具"的路线：

\begin{enumerate}[itemsep=0.2em]
\item 第 1 章给出全书 synopsis：把期权定价表述为倒向随机演化的（虚时）
      量子力学问题，把利率曲线表述为一维"量子场"。
\item 第 2--3 章是金融学背景：有效市场、随机游走、无套利与鞅、远期/
      期货/期权、Black--Scholes 方程与 Merton--Garman 方程（随机波动率）。
\item 第 4--5 章把股票期权改写成量子力学语言：定价核 = 传播子，
      Black--Scholes 哈密顿量、路径积分表示、障碍期权 = 势垒问题。
\item 第 6 章把同样的机器用于即期利率模型（Vasicek 等），
      引入正演/倒向哈密顿量与 Fokker--Planck 哈密顿量。
\item 第 7--10 章是本书核心贡献：把整条远期利率曲线 $f(t,x)$ 看作
      二维量子场（$t$ 时间、$x$ 未来时间），构造场论作用量、传播子
      （刚性/刚度）、鞅条件与风险中性测度的场论形式，并用欧洲美元
      期货数据做实证检验（第 8 章），推广到国债衍生品定价与对冲
      （第 9 章）与哈密顿量形式（第 10 章）。
\item 附录 A 补数学工具：泛函高斯积分、白噪声、朗之万方程与
      金融学基本定理。
\end{enumerate}

\section*{使用约定}

\begin{itemize}[itemsep=0.2em]
\item \textbf{公式编号}与原书一致："式 (7.38)"即原书第 7 章第 38 式；
      原书无编号的推导行在本册中不加编号。个别推导步骤从略处均注明。
\item \textbf{图表}不复制：以"原书图 N.M / 原书表 N.M"指路，请对照原书。
      提取的原书插图片文件随本工程存于 \texttt{figures/} 文件夹
      （\texttt{fig\_N\_M.png}），仅供个人对照查阅。
\item \textbf{术语}首次出现括注英文；全书术语与记号见卷末
      "术语中英对照表"与"记号总表"。
\item \textbf{记号}完全沿用原书：股票价格 $S(t)$、对数价格
      $x(t)=\ln S(t)$、即期利率 $r(t)$、远期利率 $f(t,x)$、
      零息国债 $P(t,T)$、距到期时间 $\tau=T-t$、远期利率速度场
      $\mathcal{A}(t,x)$ 等；花体字母
      $\mathcal{A},\mathcal{B},\mathcal{F},\mathcal{H},\mathcal{L},
      \mathcal{V},\mathcal{O},\mathcal{P}$ 保留原书含义。
\item 虚时（Euclidean）时间与金融时间的对应是全书反复出现的主题：
      期权定价用倒向（backward）演化，国债定价用正演（forward）
      演化，两者由 Feynman--Kac 公式联系。
\end{itemize}

\section*{阅读建议}

熟悉量子力学与路径积分的读者：可直接从第 4 章进入，把第 2.5 节
（无套利与风险中性测度）当作金融侧的唯一"公理"补读；第 7 章的
场论构造是全书技术核心。熟悉金融的读者：第 2--3 章可快速通过，
注意第 4 章起"哈密顿量 $H = -\frac{\sigma^2}{2}\partial_x^2 +
\frac{\sigma^2}{2}\partial_x + r$ 与薛定谔形式"的对应，以及第 5 章
作用量 $S = \int L\,dt$ 与定价核 $p \sim e^{-S}$ 的路径积分表示。
第 8 章的实证部分（相关结构、刚性/刚度参数拟合、心理未来时间）
是理解第 7 章模型为何取该形式的关键证据。

\tableofcontents

\mainmatter

% ch01 Synopsis (书页1-6 / p019-024)
% 内容核对说明：原书第1章（书页1--4，p019--p022）为全书Synopsis，纯散文，
% 全章无编号公式、无图表（已对照PNG逐页核验，文字层与MinerU md一致）。
% 书页5（p023）为"Part I"扉页、书页6（p024）为空白页，均不属于本章正文。
% 故本章无 \tag 公式，文末 %TAGS: 清单为空。
\chapter{概要}

原书第1章是一份四页的概要（Synopsis）：不含一个公式、一幅图，任务是画出
全书的论证地图。Baaquie先申明全书的两条主线，澄清"量子"一词的准确含义；
然后按三个组成部分逐一预告第2--10章——每章要解决什么问题、动用什么数学
工具、得到哪一类结果；章末一段说明全书体例：数学细节放进各章附录，可跳过
而不伤连贯性。

对本导读的读者，第1章的正确用法是把它当作"支票清单"：这里点到的每一个
概念——鞅条件、障碍期权、心理未来时间、刚性模型、重正化——都会在后面某章
兑现为精确定义与推导。本章没有任何编号公式，导读的任务相应变为：把每句
预告"翻译"成它在前方章节中的数学落点，并注明"见第N章"。阅读本章不需要
任何先修知识；相反，它是为全书准备的先修地图。

\section{两条主线与"量子"的含义}

本节要解决的问题是全书的名分：一本讲期权与利率的书，凭什么叫"量子"？
作者的回答分两层——两条主线，与一条概念边界。

全书的两条主线。其一：在量子理论的概念与数学框架内定义并分析量化金融
（quantitative finance），特别强调其路径积分（path integral）表述；其二：
把量子场论（quantum field theory）的技术与方法论引入利率（interest rate）
研究。两条主线共用同一套工具，差别在对象的规模：前者（股票期权）是有限
自由度问题，后者（利率曲线）是无穷自由度问题——这条"自由度升级"的路线
正是原书三部分划分的依据，见下面三节的解读。

"量子"的边界。本书不打算用量子理论去改造金融学的基本原理：无套利、货币
的时间价值这些金融学公理原样保留。"量子"一词专指量子理论的抽象数学
构造——概率论、态空间（state space）、算符（operator）、哈密顿量
（Hamiltonian）、对易方程（commutation equations）、拉格朗日量
（Lagrangian）、路径积分、量子化场（quantized field）、玻色子（boson）、
费米子（fermion）等。作者的立场可以概括为：取量子理论的数学，不取其
物理。这对两类读者各是一条信息：物理读者不必寻找波函数的物理诠释；金融
读者可以放心，定价逻辑仍从金融公理出发（见第2--3章）。

为什么路径积分与哈密顿量唱主角。原书给出三层理由。第一，随机量子过程的
路径积分与哈密顿量表述，同随机微积分（stochastic calculus）既等价又相互
独立，而随机微积分正是当今数理金融的基石之一。等价，保证旧结果都能在
新语言中复现——Black--Scholes方程与Merton--Garman方程都将改写为哈密顿量
与路径积分形式（见第4、5章）；独立，保证新语言能走出旧框架——把整条远期
利率曲线当作二维量子场建模，是随机微积分做不到的事（见第7章）。第二，
路径积分为重正化（renormalization）提供了高效的实施框架，而非线性量子
场论恰恰需要重正化才有适定的量（非线性远期利率模型将用到，见第7章）。
第三层隐含在书名的定义里：量子金融（Quantum Finance）指量子理论的概念、
方法与数学同理论及应用金融学的合成——本书即以此合成物为名。
\footnote{对金融背景的读者，"等价"可以落到实处：马尔可夫过程的路径积分
表示与Kolmogorov前向、后向方程描述同一对象，费曼--卡茨（Feynman--Kac）
公式是两种语言之间的桥，第5章将用它把期权定价写成路径积分。}

\section{写作体例：先改写已知，再推导新知}

本节交代全书的教学法承诺，读者可以据此规划自己的阅读路径。

从已知到新知。为了减轻进入量子数学（尤其是路径积分）的门槛，凡引入新
思想，原书都先做"改写"再做"推导"：把金融中已经确立的模型与已知结果用量子
形式重新表述一遍——这是较容易的练习——然后再去推导新结果。作者指出这一
做法的一个意外好处：习惯于随机微积分与偏微分方程（partial differential
equation）的金融理论工作者，看到的是一套与随机微积分完全平行、互为镜像的
形式论，向量子场论数学的过渡因而是平滑的。

自足性与范例。所有重要方程都从金融学第一性原理推导，主要结果尽可能给出
完整而自足的数学处理。若干精确可解模型被反复精讲，充当演示量子金融一般
原理的范例。这本"范例清单"后文可以点名：障碍期权的哈密顿量精确解
（见第4章）、随机波动率期权的路径积分精确解（见第5章）、Vasicek模型的
路径积分解（见第6章）、"刚性"高斯场论的市场拟合（见第7、8章）。这些模型
本身就有金融学上的兴趣，为路径积分的应用提供了现实语境。

数值算法的边界（原书脚注1）。金融问题的路径积分表述为强力计算算法打开了
大门，但数值算法是专门课题，原书除第5.16节一笔带过的蒙特卡罗（Monte
Carlo）算法外不作讨论。

\section{第一部分：金融学基本概念（原书第2--3章）}

原书把全书分为三个组成部分；本节起依次对应。第一部分题为"金融学基本
概念"（Fundamental concepts of finance），回顾金融学的标准概念与期权
理论的方程，任务是为后两个部分定义框架与语境。这部分材料是标准的——
熟悉通行期权教材的读者可以快读。

第2章"金融学导论"（Introduction to finance）面向不熟悉金融的读者，引入
理解本书方程所需的概念与术语：风险与收益（risk and return）、货币的时间
价值（time value of money）、套利（arbitrage）、对冲（hedging），最后是
国债（Treasury Bond）与固定收益证券（fixed income securities）——最后
两项为第二、三部分的利率理论预作铺垫（见第2章）。

第3章"衍生证券"（Derivative securities）引入金融衍生品（financial
derivatives）概念并讨论其定价：回顾经典分析，简述随机微积分的数学，
阐明随机过程与朗之万方程（Langevin equation）的联系，并从第一性原理推导
随机波动率（stochastic volatility）股票期权的价格，即Merton--Garman方程
及其解（见第3章）。值得注意的是，这块"标准材料"的最后一项——波动率本身
的随机性——已经埋下全书的伏笔：让波动率也变成随机对象，正是后文场论思想
登场的入口之一。

\section{第二部分：有限自由度系统（原书第4--6章）}

第二部分题为"有限自由度系统"（Systems with finite number of degrees of
freedom），把哈密顿量与路径积分用于股票期权与随机利率模型。这里出现全书
第一个关键定义：自由度（degree of freedom）指每一时刻$t$的独立随机变量
（random variable）个数。随机演化的股票价格$S(t)$有一个自由度，即期利率
（spot interest rate）$r(t)$同样有一个。这个定义把"系统规模"归结为"每个
时刻随机源的维数"，立刻给出一条分界线：有限自由度系统算出的所有量完全
有限，不需要重正化即适定；只有当自由度成为连续统（第三部分的利率曲线）
时，重正化这类场论手段才成为必需。
\footnote{用随机微分方程的语言说，自由度就是驱动方程的独立布朗运动
（Brownian motion）的条数：$S(t)$或$r(t)$由一条布朗运动驱动，是单自由度
系统；第7章的远期利率曲线则由"每个未来时刻各一条"的布朗运动驱动，是无穷
自由度系统。}

第4章"哈密顿量与股票期权"（Hamiltonians and stock options）把衍生证券
定价整体改写为量子力学问题，四件事各有落点：对常数波动率与随机波动率两种
情形，推导驱动期权价格的哈密顿量（$H_{BS}$与$H_{MG}$，见第4章）；把风险
中性（risk-neutral）演化所需的鞅条件（martingale condition）重新表述为
哈密顿量的语言（见第4章）；证明哈密顿量中的势能项代表一类路径依赖期权
（path-dependent option）（见第4章）；并用相应的哈密顿量精确求解障碍
期权（barrier option）（见第4章）。

第5章"路径积分与股票期权"（Path integrals and stock options）把期权定价
表示为费曼路径积分（Feynman path integral）：第4章的哈密顿量在"期权定价
的偏微分方程"与"它的路径积分实现"之间架起桥梁（见第5章）；书中显式算出
若干路径积分，用来演示路径积分的数学；随机波动率情形被精确解出——这是个
非平凡问题，是演示路径积分精妙之处的范例（见第5章）。此外，随机波动率
路径积分表示带来的精确简化导出一个高效算法，原书借它展示路径积分的数值
方面——即上一节提到的第5.16节蒙特卡罗算法（见第5章）。

第6章"随机利率的哈密顿量与路径积分"（Stochastic interest rates'
Hamiltonians and path integrals）回顾即期与远期利率的重要现有随机模型：
对即期利率得到福克--普朗克（Fokker--Planck）哈密顿量与路径积分，并给出
Vasicek模型的路径积分解（见第6章）；再把Heath--Jarrow--Morton（HJM）
远期利率模型改写为路径积分问题，用路径积分重新推导HJM的已知结果
（见第6章）。原书明言：第6章是为全书的主攻方向——用量子场论建模利率——
所作的准备。

\section{第三部分：利率模型的量子场论（原书第7--10章）}

第三部分题为"利率模型的量子场论"（Quantum field theory of interest rates
models），是本书的核心贡献所在。量子场论是研究无穷多自由度系统的数学
结构；无穷维带来许多随机微积分之外的新特征，最重要的当属非线性场论
（nonlinear field theory）的重正化。本部分所有章节都把远期利率（forward
interest rate）当作量子场处理。

用一行记号即可把"曲线即场"说清楚。远期利率$f(t,x)$是$t$时刻约定的、未来
时刻$x>t$的瞬时贷款利率；固定$t$，曲线
\begin{equation*}
\theta=x-t,\qquad x\mapsto f(t,t+\theta),\qquad \theta\in[0,\infty),
\end{equation*}
就是以未来时间（future time）$\theta$为坐标的一条利率期限结构曲线；第7章
让它随日历时间$t$随机涨落——曲线上每个点各有一个随机源，自由度成为
连续统，"场"由此得名。（记号预告：远期利率速度量子场$\mathcal{A}(t,x)$、
传播子$D(\theta,\theta')$、刚性（rigidity）参数$\mu$与刚度（stiffness）
参数$\lambda$，定义见第7章与卷末记号总表。）

第7章"远期利率的量子场论"（Quantum field theory of forward interest
rates）把路径积分形式用于随机演化曲线：远期利率被建模为随机涨落的曲线，
自然地由量子场论描述。各种线性（高斯）模型（linear (Gaussian) models）
被逐一研究，以展示场论进路的理论灵活性；心理未来时间（psychological
future time）为高斯场论模型提供自然推广（见第7章）；高斯模型的鞅条件被
解出，并给出计价物变换（change of numeraire）的场论推导（见第7章）；
非线性场论则在为取正值的远期利率以及带随机波动率的远期利率建模时自然
出现（见第7章）。

第8章"远期利率的实证与场论模型"（Empirical forward interest rates and
field theory models）转向实证（empiricism）：讨论远期利率的实证性质，
展示如何用欧洲美元（Eurodollar）期货合约（futures contract）的市场数据
校准（calibration）并检验场论模型。本章最重要的结果：所谓"刚性的"
（stiff）高斯场论模型对远期利率的市场行为给出几乎精确的拟合；这项实证
研究为场论刻画远期利率行为的有效性提供了有说服力的证据（见第8章；刚性、
刚度参数的数学身份见第7章）。

第9章"国债衍生品与对冲的场论"（Field theory of Treasury Bonds'
derivatives and hedging）研究国债期货、债券期权（bond option）与利率上限
期权（interest rate cap）的定价；讨论场论利率模型中国债的对冲（hedging），
并证明它是更标准进路的推广；对瞬时与有限时间两种情形都导出精确的对冲
结果，并作半实证（semi-empirical）分析（见第9章）。

第10章"远期利率的场论哈密顿量"（Field theory Hamiltonian of the forward
interest rates）为线性与非线性理论推导态空间（state space）与哈密顿量：
哈密顿量表述给出非线性远期利率以及随机波动率远期利率的鞅条件的精确解；
并给出非线性理论的计价物变换、债券期权价格与远期利率定价核（pricing
kernel）的哈密顿量推导（见第10章）。至此，第三部分在第7章"路径积分侧"与
第10章"哈密顿量侧"各形成一条完整的理论闭环。

\section{体例与使用建议}

原书第1章结尾交代全书体例。各章正文聚焦量子形式应用于金融的概念与理论
方面；更数学性的材料放进各章之后的附录（Appendix）。少数几处读者会受益于
更多细节，推导直接写进正文但用小号字体排印；书末另有一个总附录，收录对
全书材料起辅助作用的数学结果。设置这些附录的用意是把所有专题处理得完整
而全面：打算动手复算的读者会用上它们；原则上，各章附录与小号字推导都可以
跳过，不损失任何连贯性。

本导读增补一条使用建议：初读者把本章当目录，按第2--10章的预告分配深浅；
要动手计算的读者盯住上文那份"范例清单"——第4章障碍期权、第5章随机波动率、
第6章Vasicek与HJM、第7--8章刚性高斯模型——它们是全书精确可解结果的完整
名单。术语与记号可随时查卷末的术语中英对照表与记号总表。

%TAGS: （原书第1章Synopsis无编号公式、无图表，tag清单为空；已对照PNG逐页核验，p019--p022为正文，p023为Part I扉页，p024空白）

% ch02 Introduction to finance (书页7-24 / p025-042)
\chapter{金融学导论}

本章是全书的“金融学地图”：它不推导任何具体的定价公式，而是把后文需要的最小编号集合一次配齐——股票价格$S(t)$与波动率$\sigma$（第3--8章的主角）、即期利率$r(t)$与远期利率$f(t,x)$、零息与附息国债价格$P(t,T)$、$\mathcal{B}(t,T)$（第9--15章的主角），以及贯穿全书的三个原理：有效市场假说（回答“模型为什么必须是随机的”）、货币的时间价值（回答“为什么要贴现”）、无套利与鞅测度（回答“为什么贴现期望就是价格”）。本章编号公式只有10个，但式(2.1)、(2.4)、(2.6)、(2.10)组成的“贴现期望定价公式族”是全书反复调用的引擎。阅读本章只需初等概率论；鞅的严格定义见原书附录A.1，完备市场中风险中性测度的存在与唯一性见原书附录A.6。

\section*{开篇：资本、金融资产与金融工具}

原书正文（书页7--8）先划定学科版图。经济学关心维系社会物质与精神生活的各种生产活动，其实物形态是实物资产（real assets）——资本品、管理、劳动力等；资本（capital）指一个社会实物资产的经济价值。在多数发达经济中，经济资产取货币化形态，资本可以被赋予货币价值即纸面形态，称为资本的货币形式（money form of capital）。金融学（finance）就是研究资本货币形式的经济学分支，不确定性与风险在其中具有根本重要性。本书聚焦金融资产（financial assets）：与实物资产相对，它们是“纸”，持有人凭其持有对一部分实物资产及其产生收入（若有）的索取权——例如股票持有人有权分得年度红利，并在公司清算时按比例分得剩余资产。金融资产又称证券（security），其具体形态（股票或债券等）称金融工具（financial instrument）；同一张纸对发行方而言同时是金融负债（financial liability），因为其利润与资产要在全体股东之间分割。

原书在此强调金融学区别于其他经济学分支的一点：它首先是实证学科（empirical discipline）。金融业每天产生海量数据，另有堆积如山的历史数据；市场强迫模型“实用、透明、可校准、可检验”。这使金融比其他经济学分支更接近自然科学——也正是本书敢于引入物理数学工具的经验论基础。

一个关键的供给结构：股票与债券处于正净供应（positive net supply）；衍生品则处于零净供应（zero net supply）——持有人数目恰等于卖出人数目，是纯粹的零和博弈，持有人的收益恒等于卖出方收益的相反数。这一点预先决定了衍生品的定价逻辑：不能靠“供求均衡”，只能靠无套利复制（见本章2.5、2.6节）。

原书把金融工具归为三种初级形态。其一，股权（equity），即普通股与股份：代表公司所有权份额，不保证回报，只有公司盈利时宣布红利的按比例份额；股价随公司经营起落，故股东暴露于公司风险之下，但只负有限的责任（limited liability）——最多损失初始投资，因此股票价值非负，最小值为零；股权是资产，因为持有人是资本的净所有者。其二，固定收益证券（fixed income securities），即债券（bond）：由公司与政府发行，承诺单笔或一列固定支付；债券是债务工具，发行人实为向买方借款，偿付安排在固定的时间区间上，该区间称为债券的到期日（maturity）。债券种类繁多：例如五年期附息美国国债承诺每半年付息、五年末还本，而零息美国国债只在到期日一次性付清。债券常被认为比股票安全——只要发行人不破产且持有人能持有到期，回报有保证；但由于利率风险、信用风险以及外币债券的汇率风险，债券组合的亏损完全可能不亚于甚至超过股票组合。其三，衍生证券（derivative securities）：收益派生自其他金融资产——例如取决于某股票或另一个衍生品的价格。

三种初级形态可近乎无穷地组合出新工具：例如优先股（preferred stock）兼有股权与债务特征——既分得发行人股权的份额，又像债券一样收取一列（固定）支付。原书随后给出理论金融的三大问题：货币的时间价值、风险与收益、证券与资产的估值；而为投资者量身创造新工具并做无套利定价，其设计、风险—收益分析与对冲构成金融工程（financial engineering）——这正是本书前半（期权与对冲）与后半（利率与固定收益）两条主线的总纲。

\section{有效市场：证券的随机演化}

本节要解决的问题是：股票价格凭什么应当用随机过程来描写？原书的回答不是“数据看起来像随机”，而是把随机性当作有效市场假说（efficient market hypothesis）的逻辑推论：一旦市场消化了全部可得信息并达到均衡，剩下的价格变动就只能是噪声。

先补齐市场词汇。买入（或约定买入）金融资产称持有多头（long position）或做多（going long）；卖出称做空（short position，shorting）；卖出自己并不实际持有的资产称卖空（short selling），回购日通常定在未来某时。市场组织有三种基本形态：直接市场（direct market，买卖双方直接搜寻成交）、经纪市场（brokered market，经纪人收费撮合）与拍卖市场（auction market，买卖双方在集中化市场中同时互动）。

有效市场（efficient market）概念与证券的“公平价格”（fair price）绑定：市场均衡时证券取得公平价格，投资者不再交易；均衡中的有效市场，其工具价格对公平价格只有小而暂时的偏离。有效市场又与市场信息紧密相关——区分不同市场参与者的正是各自可得的市场信息量。市场信息分三类：(a)金融资产价格与收益的历史数据，(b)关于金融资产各方面的公开数据，(c)少数参与者私有的信息；据此可分别定义市场效率的弱式（weak）、半强式（semi-strong）与强式（strong）形态。直观地说，有效市场意味着大多数买卖双方财富与信息对等，没有任何一方占（不公平的）优势。原书给出有效市场假说的精确表述，其内容可转述为：对处于均衡的金融市场，没有任何参与者在既有禀赋与信息下想再交易；价格揭示可得的市场信息；新信息以随机的、零碎的方式流入，由于信息本身不完备，市场参与者做出随机响应，于是各金融工具的价格随机变动。

这里有一个影响深远的推论：市场一旦均衡，所有证券的价格变化（除漂移（drift）外）都是随机的。理由是：有效市场中工具价格已经吸收了全部市场信息并形成均衡价格，此后任何对均衡的偏离都应当不确定、不可预测，且向上与向下变动的可能性相等。于是价格变化必须用随机变量（random variable）表示：设某股权在时刻$t$的价值为$S(t)$，其变化率$\mathrm{d}S/\mathrm{d}t$被建模为随机变量，这又蕴含$S(t)$本身也是随机变量，只在初始时刻$t_0$给定确定性初始条件；随机性的强度由证券的波动率（volatility）刻画，记作$\sigma_S$或简记$\sigma$。波动率越大，价格的随机涨落越大；$\sigma=0$则未来演化完全确定。\footnote{这一论证是启发式的。它的严格化就是后续章节的随机微分方程$\mathrm{d}S=\mu S\,\mathrm{d}t+\sigma S\,\mathrm{d}W$的形式：漂移$\mu$在2.3节以风险溢价的面目回归，$\sigma$从第3章起进入Black--Scholes与Merton--Garman方程。}所有投资者面对的风险，正是金融工具随机演化的反映，而后者最终反映（金融）市场吸纳底层实体经济全部相关特征的方式。

要紧的是，有效市场假说并不主张消息不能移动市场。它说的是：未预期的、出乎意料的新信息会扰乱各证券市场价格的均衡，并把它们系统地搬移到一组新均衡价格；一旦均衡达成，普通信息几乎人人可得，于是只造成所揭示价格的随机变动。

该假说可否被经验检验？原书引文献[23]指出，有效市场的存在暗中包含两个假说——效率假说本身，加上“市场处于某个特定均衡”的假说；能够被经验检验的只是二者的联合假说（joint hypothesis）。围绕这一联合检验，学术界关于市场效率的争论经久不息。

原书以一个物理类比收束本节（书页11首段）：市场均衡类似热力学系统的平衡——平衡态中单个粒子的位置与速度仍然随机，正如均衡中金融工具的价格仍然随机；市场效率则类似热机效率：没人指望真实热机达到$100\%$的效率，$60\%$--$70\%$已属常见；同样，即使金融市场并非完全有效，应用基于有效市场概念的数学模型往往仍然成立。

\section{金融市场}

本节回答两个问题：金融工具在哪里、以何种组织形式交易？市场的体量有多大？

金融市场的第一层划分是资本市场（capital markets）与货币市场（money markets）。资本市场交易初级形态的金融工具——股权、债务与衍生品；指数（indexes）本是一篮子证券的加权平均、属于资本市场，但因其重要性与深度被当作独立实体处理。货币市场严格说属于债务市场，但其工具高度流动、期限极短（现金及现金等价物、外汇交易等），故另设市场。原书给出如下分类（转述）：(1)资本市场——股权市场（普通股、优先股），债务市场（国债（政府）票据与债券、公司债与市政债、抵押贷款支持证券（mortgage-backed securities，MBS）），衍生品市场（期权、远期与期货）；(2)指数——股权指数（道琼斯与标普指数、日经指数、DAX指数、STI指数等），债务指数（债券市场指标）；(3)货币市场——现金定期存款、短期国库券（Treasury bills）、大额存单（certificates of deposit）、商业票据（commercial paper）、欧洲美元存款（Eurodollar，指存于非美银行或美国银行海外分行的美元存款）。

\paragraph{原书图2.1（1993年美国资本市场构成，总值\$11.3万亿）}图饼显示：股权仅占$36\%$，债务市场合计占$64\%$——国债$26\%$、MBS $14\%$、公司债$11\%$、市政债$11\%$、高收益债$2\%$。若再把货币市场计入，股权份额还要更低。

\paragraph{原书图2.2（1993年全球债务市场的借款人构成）}全球债务市场在正文中的数字为1993年\$14.08万亿，并称“图2.1[100]给出主要国际借款人”——实际上给出借款人构成的是图2.2，其图题写的是\$14.8万亿。正文与图题在交叉引用与数字上均不一致（原书排印疏漏，两说照录）。借款人构成：美国$46\%$、日本$18\%$、德国$10\%$、英国$8\%$、其他$8\%$、意大利$5\%$、法国$5\%$。

量级感（原书照录的数据）：美国GDP（2001年）约\$10万亿；美国信贷市场规模（2002年）约\$29万亿，其中金融部门借款占\$9万亿；对照之下，美国股票总市值（2002年）约\$12万亿——债券市场的体量明显大于股票市场。

衍生品有两条交易渠道：受监管的交易所与不受监管的场外市场（over-the-counter，OTC）。市场规模通常以合约的名义价值（notional value）度量。交易所交易衍生品（涵盖各大期权与期货交易所的全部合约）的近期估计为名义金额\$13万亿至\$14万亿；OTC衍生品为特定客户定制，2000年底未平仓合约的名义金额估计为\$95.2万亿，且专家认为此数很可能还偏低。

\section{风险与收益}

本节把投资决策的两要素——收益与风险——变成可计算的量：收益用比例变化定义，风险用收益率的（隐含）概率分布的二阶矩定义。

设在时刻$t$以价格$S(t)$买入某股票，持有时长$T$，期末价格为$S(t+T)$，期间获得红利$d$，则该时段的（比例）收益率（fractional rate of return）$R$定义为
\begin{equation*}
R=\frac{S(t+T)+d-S(t)}{S(t)},
\end{equation*}
而$R/T$为瞬时收益率。风险来自未来股价可能上涨也可能下跌这一不确定性：市场可能繁荣（股价上涨）、衰退（股价暴跌）或沉闷（股价微动）。投资者可以为每种情景（scenario）指派概率，从而把市场的不确定性变成可把握的对象。原书表2.1给出一个典型例子（证券现价\$100，年度情景）：

\begin{center}
\begin{tabular}{llccc}
\hline
$s$ & 情景 & 期末价 & 概率$p(s)$ & 年收益$R(s)$\\
\hline
1 & 沉闷（doldrum） & \$100 & 0.70 & 0.00\\
2 & 下行（downturn） & \$85 & 0.20 & $-0.15$\\
3 & 繁荣（boom） & \$105 & 0.10 & 0.05\\
\hline
\end{tabular}
\end{center}

以离散变量$s$标记情景，$p(s)$为其概率，$R(s)$为其收益。投资的期望收益（expected return）即收益的平均（均值）值：
\begin{equation*}
\bar R=\text{期望收益}=\sum_s p(s)\,R(s)\equiv E[R],
\end{equation*}
其中记号$E[X]$表示随机量$X$的期望值。获得期望收益所冒的风险，显然是收益可能偏离该期望值：概率论告诉我们“实际收益$=$期望收益$\pm$标准差（以一定概率）”。偏离的精确幅度与可能性，只有知道股价$S(t)$的概率分布$p(S)$才能算出。标准差$\sigma$是方差的平方根，方差定义为
\begin{equation*}
\sigma^{2}=\sum_s p(s)\bigl(R(s)-\bar R\bigr)^{2}\equiv E\bigl[(R-\bar R)^{2}\bigr],
\end{equation*}
投资的风险即由$\sigma$度量：风险越大$\sigma$越大，反之亦然。按表2.1的数字，年初投资\$100的投资者年末的期望现金为$\$100\times[1+\bar R]\pm6=\$97.50\pm6$。

若收益分布重尾到方差不存在——例如证券服从Lévy分布——则$\sigma$为无穷，此时更合适的风险度量是在险价值（value at risk）。

接着建立无风险基准。银行定期存款等工具被视为无风险（risk-free）的：无风险工具在时刻$t$的收益率，是存入无风险银行的瞬时存款所赚的比率，称为即期利率（spot interest rate）或隔夜拆借利率，记作$r(t)$。于是对无风险工具
\begin{equation*}
\sigma_{\text{risk-free}}=0,\qquad \text{无风险收益率}=\text{即期利率}=r.
\end{equation*}
风险中性（risk-neutral）投资者满足于等于即期利率$r$的回报。但对风险投资$\sigma>0$，要诱导投资者承担高风险，必须有相应的高回报：资本市场为高风险投资标出溢价，以促成资本流向高风险处。风险溢价（risk premium）是高风险投资的收益率高出无风险利率的数额；对平均年收益率为$\bar R$的投资，原书把比值$(\bar R-r)/\sigma$称为风险溢价，又名Sharpe ratio。\footnote{通行术语与此略有差别：risk premium通常专指分子$\bar R-r$，而比值$(\bar R-r)/\sigma$才是Sharpe ratio。原书二者混用，导读照录并提醒读者注意。}投机者（speculator）只有在分析显示潜在回报带有足够风险溢价时才投入高风险证券——这不同于没有溢价也承担高风险的赌徒（gambler）。

本节末尾，无套利原理（principle of no arbitrage）首次亮相：任何无风险金融工具的收益率都不可能高于无风险利率。换言之，天下没有免费午餐——想收获高回报，就必须承担相应的高风险。其数学含义留待2.5节展开。

\section{货币的时间价值}

本节的任务是回答：不同时刻的美元不能直接相加比较，换算的“汇率”是什么？答案是即期利率，而随机利率下的换算规则由本章第一个编号公式(2.1)给出。

先看钱的价值为何依赖时间。资本的货币形式背后是会随时间折蚀的真实生产性资产；通货膨胀、货币贬值等因素也使金融资产所代表的价值随时间变化；金融资产代表启动或促成经济活动的能力，而机会总与变动不居的环境相连。诸如此类的原因使货币的有效价值强烈依赖时间。

如何估算货币的时间价值？原书的经济学答案是：对货币时间价值的全部内生与外生影响，都汇总在一个量里——即期利率。投入其他风险工具的钱估值更复杂，因为风险溢价因投资者而异；归根结底，货币的时间估值就是把未来值贴现（discounting）成“期望的”现值（present value）。

具体机制：设$t_0$时刻持有\$1，将其投入固定存款之类的无风险工具，并把收益不断再投入同一存款（复利（compounding））。由于无风险工具的收益率为$r(t)$，到未来时刻$t_*$这笔投资的无风险价值变为
\begin{equation*}
t_0<t_*:\quad \$1\ \text{在}\ t_0\ \text{时刻}
=\$\,e^{\int_{t_0}^{t_*}dt\,r(t)}\ \text{在}\ t_*\ \text{时刻}\quad(\text{现金的增值})，
\end{equation*}
反过来，把未来值贴现回今天：
\begin{equation*}
\$1\ \text{在}\ t_*\ \text{时刻}
=\$\,e^{-\int_{t_0}^{t_*}dt\,r(t)}\ \text{在}\ t_0\ \text{时刻}\quad(\text{现金的贴现值})。
\end{equation*}
货币时间价值的实质因此是：\$1并不是衡量价值的正确单位——它的有效价值随时间不断变动；正确的单位是贴现量（discounted quantity）。未来现金流应乘以贴现因子（discount factor）$\exp\{-\int_{t_0}^{t_*}dt\,r(t)\}$折到今天，今天的现金流应乘以它的倒数推到未来。

$r(t)$本身如何决定？从第一性原理出发须研究宏观基本面、货币的供给与需求等；利率反映消费的边际效用——诱使人们放弃当下消费、把钱储蓄（投资）起来备作未来消费的比率。后文研究即期利率模型时，$r(t)$将被当作随机（random）变量处理；此时贴现要对随机函数$r(t)$的一切可能取值取平均，贴现因子成为
\begin{equation*}
t_0<t_*:\quad \$1\ \text{在}\ t_*\ \text{时刻}
=\$\,E\!\left[e^{-\int_{t_0}^{t_*}dt\,r(t)}\right]\ \text{在}\ t_0\ \text{时刻}
\tag{2.1}
\end{equation*}
（即现金的贴现值；式中前置的\$表示金额以美元计，$E$为期望。）

推导主线只有一步，却极为要紧：确定性情形的贴现式$e^{-\int r}$——承认$r(t)$是随机过程——“对贴现因子整体取期望”得到式(2.1)。注意期望作用在整个指数上，而不是“先对$r$取平均再代入指数”。这一差别（Jensen不等式的效应）意味着随机利率下的现值一般不等于用平均利率算出的现值；更重要的是，式(2.1)已经埋下伏笔：期望总要对着某个概率测度取，而“选哪个测度”正是2.5节的主题。式(2.1)也是本章后文零息国债定价公式(2.6)的种子。

\section{无套利、鞅与风险中性测度}

本节是全章的理论核心，要解决的问题是：给衍生品定价时，怎样保证定出的价格本身不制造套利机会？答案可以压缩成一句话：在某个（风险中性）测度下让贴现价格成为鞅，然后“价格$=$贴现后的条件期望”。

套利（arbitrage）的定义：通过同时进入两笔或多笔金融交易——无论在同一市场还是跨市场——获取无风险（有保证）的利润。既然经济中存在现金存款之类的无风险工具，套利就特指获得高于货币市场无风险收益的保证收益。原书举了一个教科书式例子：同一公司股票在纽约交易所值US\$1，在新加坡值S\$1.8，而汇率为US\$1$=$S\$1.7；经纪人同时在纽约买进100股、在新加坡卖出100股，即可锁定S\$10的无风险利润。交易成本会消灭小散户的套利机会，但对交易成本近乎为零的大型券商，套利是重要的利润来源；同时套利者的抛售又把新加坡股价压向S\$1.7——套利活动本身就是价格回归均衡的动力。

有效市场与无套利的关系要说清楚：有效市场中不存在套利机会；套利是资本市场在实践中得以作为有效市场运转、并确定金融工具均衡（“正确”）价格的机制之一。有效市场的存在是无套利原理成立的充分条件而非必要条件；均衡中无套利机会成立。无套利是一个稳健（robust）的概念，因为它只表达“所有投资者宁要多财富、不要少财富”这一偏好；而大多数市场均衡模型对投资者行为施加了更苛刻的假设。

理论金融的一个重要结果（原书引[23,40,100]）是：要使金融工具的价格杜绝任何套利可能，就必须让该工具的贴现价值按鞅过程（martingale process）演化。要注意：证券（如股票）的真实市场演化并不是鞅过程——若是，持有该证券便没有风险溢价。鞅演化是给衍生品定价而设的理论构造：先在“风险中性世界”里定价，使价格免于套利。概率论中的鞅（原书附录A.1讨论）是“公平博弈”（fair game）概念的数学化；有效市场中证券的无风险演化等价于其演化满足鞅条件（martingale condition）。由于真实投资者并非风险中性、要求风险溢价，他们的（真实）演化需要从风险中性测度出发作变换测度（change of measure）。

鞅的定义。设随机过程由$N+1$个随机变量$X_i$（$1\le i\le N+1$）给出，联合概率分布函数为$p(x_1,x_2,\ldots,x_{N+1})$（原书附录式(A.2)），则鞅过程由如下条件概率定义：
\begin{equation*}
E\!\left[X_{n+1}\mid x_1,x_2,\ldots,x_n\right]=x_n
\qquad:\ \text{鞅}
\tag{2.2}
\end{equation*}
左端是随机变量$X_{n+1}$在$X_1,X_2,\ldots,X_n$取值为$x_1,x_2,\ldots,x_n$条件下的条件期望（conditional expectation）。在金融应用中，这些随机变量取某股票在时刻$t_1,t_2,\ldots,t_{N+1}$的未来价格$S_1,S_2,\ldots,S_{N+1}$。

要把鞅条件用于股价演化，必须先贴现：被比较的是两个不同时刻的价格，按2.4节的教训须先把金额换到同一时刻的时间可比单位——贴现由无风险利率完成。原书附录A.6证明：对完备市场（complete market），存在唯一（unique）的风险中性测度（risk-neutral measure），使得该资产的一切衍生品的贴现演化都满足鞅条件（引[40,42]）。\footnote{现代教科书的“资产定价基本定理”表述为两条：无套利$\Leftrightarrow$存在等价鞅测度（第一定理）；市场完备$\Leftrightarrow$测度唯一（第二定理）。原书的表述与之一致，参见Harrison--Kreps与Harrison--Pliska的经典谱系及原书附录A.6；对照阅读可参考Shreve《Stochastic Calculus for Finance》。}

定价公式族的出发点：设某股权在未来时刻$t$的价值为$S(t)$，并假设存在贴现股价的无风险演化
\begin{equation*}
e^{-\int_0^t r(t')\,dt'}\,S(t)
\tag{2.3}
\end{equation*}
使它按鞅过程演化（引[42]）。由鞅的条件期望表述（原书附录式(A.40)），股权贴现值在时刻$t$的条件期望就是其现值$S(0)$，换言之
\begin{equation*}
S(0)=E\!\left[\left.e^{-\int_0^t r(t')\,dt'}\,S(t)\,\right|S(0)\right].
\tag{2.4}
\end{equation*}

推导主线值得逐步拆开（这一步没有任何高深数学，力量全部来自测度的选择）：记$Y_t=e^{-\int_0^t r(t')dt'}S(t)$，即把$t$时刻的股价用无风险利率贴现回零时刻；(2.2)中取$n=0$，“以今日全部信息为条件，$Y_t$的期望等于$Y_0$”，而$Y_0=S(0)$，这正是式(2.4)。换个角度读：式(2.4)说“贴现后的未来价格的期望$=$今天的价格”。为什么这等价于无套利？因为若期望价格偏离$S(0)$，买低卖高并持有到期即可获得超出无风险利率的保证收益——恰是套利的定义；鞅条件堵死了这条路。换测度不改变路径的实现，只改变路径的概率权重，从而把真实世界漂移（含风险溢价）换成无风险漂移，这就是“风险中性定价”的机制。

式(2.4)的通用性极强——它对任何证券都成立，是全书引用频率最高的公式之一：给衍生品定价时，把$S(t)$换成衍生品价格、在风险中性测度下取贴现期望即可；本章的式(2.6)与(2.10)都是它的直接特例。有作者（引[42]）把鞅测度的存在性称为金融学基本定理（fundamental theorem of finance），原书附录A.6有简述。总结成一对双向命题：若存在测度使贴现工具的演化满足鞅条件，则该工具一切衍生品的价格免于套利机会；反之，若贴现价格免于套利机会，则存在鞅测度。本书所分析的大多数模型都在鞅测度下演化，从而保证一切（衍生）金融工具的价格无套利。

\section{对冲}

本节处理一个自然的问题：金融工具的演化既然随机，能否用若干风险资产拼出一个无风险的投资组合（portfolio）？换言之，能否用一个工具的随机涨落抵消另一个工具的随机涨落，且抵消精确到组合整体成为无风险证券？对冲还有一层比较功能：收益相同的两个组合，经过对冲的那个风险更低，因而一般是更优的组合。

对冲（hedging）是把某金融工具放入与其他相关工具构成的投资组合、以削减其随机涨落的程序的统称。完全对冲（perfectly hedged）的组合免除一切随机涨落——被对冲工具价格的随机涨落恰被组合内其他工具的反向涨落精确抵消；实践中通常最好的结果也只是部分对冲（partially hedged）。

对冲类似买保险：其成本是买卖所需证券的交易成本，是降险必须支付的价格；交易成本高使对冲更贵，但它对抗风险依然有效。对冲常有的副作用是压低未来收益，原书举国债对冲为例：期货合约入场零成本，可用期货空头对冲债券——利率上升时对冲起效（期货盈利抵补债券跌价）；但利率下降时债券涨价，期货却亏损，反而拉低净利润。消灭“坏”波动的同时也消灭了“好”波动。期权入场虽非零成本，却常能让投资者管理风险而不必接受未来收益缩水。总之，对冲策略取决于投资者的目标。

并非一切涨落皆可对冲。完备市场中原则上存在可用于对冲某一特定工具每种风险的资产；实践中能否完全对冲取决于市场上实际可得的其他工具——衍生工具发展的主要动力，正来自对常用金融工具的对冲需求。

对冲的结构条件：至少需要第二种工具，才能尝试两种工具涨落之间的抵消；这第二种工具必须依赖被对冲的工具，否则二者的随机涨落谈不上关联。对冲初级工具时常需要其衍生工具，反之亦然：衍生工具与标的由同一随机过程驱动，其演化与标的完全相关（perfectly correlated），从而允许完全对冲。

看具体构造。设某证券（如普通股）由随机变量$S(t)$代表，要削弱持有$S(t)$的风险，需要持有第二工具——$S(t)$的一个衍生品，记作$D(S)$——使组合整体的涨落变小。设$S(t)$可以被完全对冲，对冲组合记作$\Pi(S,t)$；组合的具体形态可以是：做多（买入）一份衍生品$D(S)$，同时卖空价值$\Delta(S)$的标的股票$S$。于是在时刻$t$
\begin{equation*}
\Pi(S,t)=D(S)-\Delta(S)\,S .
\end{equation*}
由于证券$S(t)$在时刻$t$的值是已知的，若组合$\Pi(S,t)$的时间演化不含任何随机涨落——实为确定性函数——则组合被完全对冲：
\begin{equation*}
\frac{d\Pi(S,t)}{dt}:\ \text{无随机性}\ \Longrightarrow\ \text{完全对冲}.
\end{equation*}
$\mathrm{d}\Pi(S,t)/\mathrm{d}t$既无随机涨落，它就是无风险证券；无套利原理随即要求完全对冲组合的收益率等于无风险即期利率$r(t)$：
\begin{equation*}
\frac{d\Pi(S,t)}{dt}=r(t)\,\Pi(S,t).
\end{equation*}

推导主线（这是Black--Scholes--Merton论证的骨架，值得记牢）：选对冲参数$\Delta(S)$使得组合瞬时无风险——无套利强制确定性组合只赚无风险利率——由此反解出衍生品$D(S)$必须满足的方程。第3章将用伊藤微积分展开$\mathrm{d}\Pi/\mathrm{d}t$、定出$\Delta$并导出Black--Scholes方程（含随机波动率时为Merton--Garman方程）。\footnote{$\Delta(S)$就是期权交易员口中的Delta；原书此处把它抽象为“卖空价值$\Delta(S)$的股票”，第3章给出其显式表达式。}

实践中对冲可行还须满足若干条件（原书列了三条）：其一，市场必须交易衍生品$D(S)$，否则无法构造对冲组合——许多金融工具因缺乏适当的衍生工具而无法对冲，例如证券的波动率本身。其二，须查明对冲参数（hedging parameter）$\Delta(S)$是否存在、其对股价$S$的函数依赖为何——这要求知道$D(S)$与$S(t)$的精确关系，以及对$S(t)$随机动力学的完整刻画。其三，组合$\Pi(S,t)$依赖时间，对冲必须连续进行——这要求市场有足够流动性，流动性又决定对冲的交易成本。股权对冲的概念在第3章（衍生品）展开，国债对冲在第9章详细讨论。

\section{远期利率与固定收益证券}

本节把视野从股权（一维随机对象$S(t)$）切换到债务市场：那里的随机对象是整条利率曲线。远期利率与固定收益证券是债务市场的基本构件；原书把它们写成“用远期利率逐段贴现\$1”，由此得到三个恒等式(2.5)、(2.8)、(2.9)与一个鞅定价式(2.10)——这正是后半书把远期利率升级为量子场的出发点。

即期利率再回顾：在时刻$t$的瞬时贷款，借款成本是即期利率$r(t)$，通常按年百分比报价；即期利率的典型范围是每年$0.1\%$至$20\%$。

远期利率（forward interest rate，forward rate）的定义：借款人常要在未来特定时刻借钱——例如为一年后的一笔商品买卖——希望现在就锁定届时利率。资本市场为此提供的工具就是远期利率，记作$f(t,x)$：在较早时刻$t<x$以合约形式约定的、未来时段$x$与$x+dx$之间借款的瞬时利率。从数学观点看，即期利率与远期利率都是借钱——借一个无穷小时间$\epsilon$——的瞬时成本。全体远期利率构成利率的期限结构（term structure），并与收益率曲线（yield curve）相关。由定义，时刻$t$的隔夜贷款的即期利率为
\begin{equation*}
r(t)=f(t,t).
\tag{2.5}
\end{equation*}

债券。债券是政府与公司为从资本市场融资而发行的债务工具，发行人承担支付一列预定固定现金流的义务，故各类债券通称固定收益证券（fixed income securities）。国债（Treasury Bond）是无违约风险的工具；公司债、市政债及某些国家（如俄罗斯、阿根廷）的主权债原则上存在违约风险。美债的无风险性质，使美国政府能以尽可能低的利率进行大规模国际举债。零息国债（zero coupon Treasury Bond）是无风险金融工具，只有单一现金流：在未来时刻$T$支付\$1的固定收益；其在$t<T$时刻的价格记作$P(t,T)$，且$P(T,T)=1$。

定价路线一：用即期利率。由货币的时间价值，到期时刻为$T$的债券在到期前的价值，是把$P(T,T)=1$用即期利率贴现到$t$；利率被视为随机变量的一般情形下，按式(2.1)取期望得
\begin{align*}
P(t,T)&=E\!\left[e^{-\int_t^T dt'\,r(t')}\,P(T,T)\right]\\
&=E\!\left[e^{-\int_t^T dt'\,r(t')}\right]
\tag{2.6}
\end{align*}
（即$P(T,T)=1$的贴现值），其中期望值对$r(t)$所服从的随机过程取。式(2.6)表明：国债价格只是即期利率初值$r(t)$的泛函。

定价路线一（续）：附息债券化归零息债券。附息国债$\mathcal{B}(t,T)$有一列预定现金流：在递增时刻$T_i$支付价值$c_i$的息票（coupon），到期时刻$T$偿还价值$L$的本金（principal）。用“现金流相同的两种金融工具是同一的”这一原则，可以证明（引[58]）$\mathcal{B}(t,T)$可用零息债券表出：
\begin{equation*}
\mathcal{B}(t,T)=\sum_{i=1}^{K}c_i\,P(t,T_i)+L\,P(t,T).
\tag{2.7}
\end{equation*}
可见附息债券等价于一个零息债券的组合：任何零息债券模型自动给出附息债券模型。市政债、公司债与高收益债因税收规则、流动性等因素而更难建模，且有有限的违约概率，带有国债所没有的风险成分，因此须支付高于国债的风险溢价。

定价路线二：用远期利率（推导重点）。零息债券是发行人借入的一笔有限期贷款，而远期利率只对瞬时未来借款有定义，故须用远期利率逐段贴现、迭代得到现值。到期时$P(T,T)=\$1$；把\$1从未来时刻$T$逐段贴现到当前时刻$t$。做法：把时间切成间隔$\epsilon$的格点，远期利率$f(t,x_n)$定义在未来时刻$x_n=t+n\epsilon$（$n=0,1,\ldots,[(T-t)/\epsilon]$）。从未来时刻$x_n$到$x_{n-\epsilon}$的瞬时贷款的贴现因子是$e^{-\epsilon f(t,x_n)}$。对时刻$T$的确定性收益\$1逐段贴现到$t$，得
\begin{equation*}
P(t,T)=e^{-\epsilon f(t,x_0)}\,e^{-\epsilon f(t,x_1)}\cdots e^{-\epsilon f(t,x_{N-1})}\,e^{-\epsilon f(t,x_N)}\cdot 1 .
\end{equation*}
取$\epsilon\to 0$的极限：
\begin{align*}
P(t,T)&=e^{-\int_t^T dx\,f(t,x)}
\tag{2.8}\\[2pt]
\Rightarrow\ -\frac{\partial}{\partial T}\ln P(t,T)&=f(t,T)
\tag{2.9}
\end{align*}
式(2.8)用远期利率表出国债价格，是恒等式，可反过来取作远期利率的定义；由式(2.8)与(2.9)，$P(t,T)$已知则$f(t,T)$亦已知，反之亦然。

鞅条件用于远期利率。把式(2.4)的精神用于零息债券的演化：设到期时刻$T$的零息国债在$t_0<T$的价格为$P(t_0,T)$，在$t_0<t_*<T$的价格为$P(t_*,T)$。满足鞅条件意味着：$t_*$时刻的债券价格沿时间倒推至$t_0$、并以无风险即期利率$r(t)$连续贴现后，其（条件）期望必须等于$t_0$时刻的债券价格。由式(2.4)，零息债券的鞅条件有如下非常一般、与模型无关（model-independent）的表达：
\begin{equation*}
P(t_0,T)=E_{[t_0,t_*]}\!\left[\left.e^{-\int_{t_0}^{t_*}r(t)\,dt}\,P(t_*,T)\,\right|P(t_0,T)\right],
\tag{2.10}
\end{equation*}
其中记号$E_{[t_0,t_*]}[X]$表示随机变量$X$在时段$(t_0,t_*]$上的平均。

两种定义的对峙是本节的收束，也是后半书的路线选择。零息债券有两个定义：用即期利率的式(2.6)与用远期利率的式(2.8)，二者很不一样。可以取即期利率为基本量（远期利率为导出量），如式(2.6)：优点是简单，适于处理只与即期利率行为直接相关的问题；缺点是即期利率必须与观察到的各期限债券价格全部相容。而式(2.8)把$P(t,T)$写成远期利率，是更一般的表达式：远期利率原则上可直接由市场确定；对全部期限的远期利率建模时，以整条初始期限结构作为输入，因此模型不会基于观测到的债券价格生成套利机会。代价是$f(t,x)$依赖两个变量$t,x$，而$r(t)$只依赖时间$t$。尽管更复杂，把远期利率视为基本量、把即期利率仅仅当作远期利率曲线上的一点，仍是远期利率建模的一条高产路线。\footnote{这正是HJM模型（Heath--Jarrow--Morton框架）的哲学：直接为整条瞬时远期利率曲线指定随机动力学，并以无套利约束漂移。本书后半书更进一步，把$f(t,x)$写成由量子场$\mathcal{A}(t,x)$驱动的场论模型。}

\section{小结}

原书以几段短论收束全章，每段都是全书路线图的一条支线，导读整理如下。第一，金融工具的随机演化是用随机变量为金融建模的全部根基；正是金融工具的随机时间演化，提供了与量子理论及路径积分的关键连接（第4章起全面展开）。第二，风险概念居于对冲概念的中心；把风险概念推广到无穷维情形，将在为国债对冲建立量子场论（第9章）时起关键作用。第三，货币的时间价值自然引出鞅测度——其定义式的要求是“贴现后的未来值等于现值”（即式(2.4)的形状）；为非线性远期利率找到恰当的鞅测度，将耗费全书大量笔墨。第四，对冲与期权构成金融衍生品理论的基石，期权定价是本书前半的焦点。第五，本书后半处理远期利率与固定收益证券；对远期利率的研究自然引向把它表述为量子场——这是全书的主旨。

读法建议（导读）：把本章10个编号公式分成三组记忆——定价公式族(2.1)、(2.4)、(2.6)、(2.10)同属“贴现期望”结构，式(2.4)是母式，其余是特例；(2.5)、(2.8)、(2.9)是$r$、$f$、$P$三者间的恒等关系，构成利率端的概念字典；(2.2)是唯一的概率论定义，(2.3)是其作用对象，(2.7)则是化归（附息$\to$零息）。带着这张公式地图进入第3章，Black--Scholes方程的每一个组成部分都已在本章就位。

%TAGS: 2.1,2.2,2.3,2.4,2.5,2.6,2.7,2.8,2.9,2.10

% ch03 Derivative securities (书页25-44 / p043-062)
% 覆盖说明：原书第3章正文为书页25--42（p043--p060），含3.1--3.9节（3.9为章末附录）。
% 书页43（p061）为"Part II"扉页、书页44（p062）空白，不属于本章正文。
% 编号公式核对：原书本章tag为(3.2)--(3.10)与(3.12)--(3.35)，共33条；
% 原书无(3.1)与(3.11)——编号断档系原书自身如此（已对照PNG逐页核验，非漏转）。
% 原书已知排印疑点两处，见式(3.25)脚注与3.9节文字说明。
\chapter{衍生证券}

本章是第一部分"金融学基本概念"的技术核心：第2章给出了套利、对冲、货币的
时间价值这些定性观念，本章把它们锻造成两个可以实际求解的定价方程——常数
波动率下的Black--Scholes方程与随机波动率下的Merton--Garman方程。论证工具
也随之升级：先引进随机微分方程（stochastic differential equation）给股票
价格建模，再为它配备一套新的微积分——伊藤微积分（Ito calculus）；随后用
"对冲组合"（hedged portfolio）这一全书反复出现的论证模式导出定价方程。
先修内容只有第2章的无套利原理；本章结论在后面的去向是明确的：第4章把两个
方程改写为哈密顿量语言，第5章用路径积分重新求解并推广，第9章把"对冲组合"
的思路移植到国债并把"残余风险"做成场论中的可计算量。

按内容本章分四步走。第一（3.1--3.2节）按现金流特征给衍生证券分类：远期
（forward）、期货（futures）与期权（option），并定义收益函数（payoff
function）与路径无关/路径依赖（path-independent/path-dependent）之分；
第二（3.3--3.4节）建立数学引擎：股票价格的朗之万型随机微分方程与伊藤
微积分的两条新规则；第三（3.5节）用对冲组合导出Black--Scholes方程，并用
风险中性估值解出欧式看涨期权的显式价格；第四（3.6--3.7节）让波动率本身
也随机化，重复"对冲组合"论证得到Merton--Garman方程。3.9节附录处理
$\rho=0$情形的随机波动率解，它把3.4节算出的几何平均方差"三分之一因子"
直接搬来使用。

\section{远期与期货合约}

本节解决的问题是：如何在今天就锁定一笔未来交易的价格。原书用钢铁进口商
作引子：公司一年后（记$T$）要买钢材，担心钢价上涨，希望按今天（记$t$）
的市场条件把价格定死。两类合约都能实现这一点，差别在现金流的安排方式。

\paragraph{模型设定与记号}
设钢材在$t$时刻的每吨价格为$S(t)$。远期合约（forward contract）是买方与
卖方之间的双边协议：买方持有多头头寸（long position），卖方持有空头头寸
（short position）；卖方承诺在$T$时刻按远期价格（forward price）
$F(t,T)$交货，这个价格反映签约当时的市场价格与利率水平。合约在$t$时刻
签订，签约时没有任何现金易手——合约条款被选成使初始价值为零；此后合约
价值随行就市地波动，直到$T$时刻结算。到期时合约对双方的价值为
\begin{equation*}
S(T)-F(t,T)\quad\text{（多头）},\qquad F(t,T)-S(T)\quad\text{（空头）}.
\end{equation*}
即整个合约只在到期日产生一次现金流，其价值取决于届时现货价与$F(t,T)$之
差，正负皆可。

期货合约（futures contract，记$\mathcal{F}(t,T)$）同样是"今天签约、未来
$T$时刻交割特定数量资产"的协议，但有三点制度化差异：格式标准化、受严格
监管；总有一家第三方——通常是清算所（clearing house）——居中担保；以及
由此产生的保证金制度。初始保证金（initial margin）在签约时缴入；此后若
价格偏离前值，则按日缴存或划出维持保证金（maintenance margin），这一手续
称为逐日盯市（marking-to-market）。标的资产价格随机波动，期货合约的价值
因此也随机波动。
\footnote{原书脚注提醒：对一般金融工具无须严格区分"价值"与"价格"两个词，
全书将混用；唯独期货合约需要把两者分开——签约时双方都不付钱，故合约的
\emph{价值}为零，但市场会给它指派一个名义上的公允\emph{价格}
$\mathcal{F}(t_0,T)\neq 0$。}

\paragraph{关键公式}
到期日期货价格必须收敛于现货价格：
\begin{equation*}
\mathcal{F}(T,T)=S(T).
\end{equation*}
这不是假设而是套利逼出来的结论：若$\mathcal{F}(T,T)>S(T)$，卖方可做空
期货再按市价$S(T)$买入交割，赚取无风险差价；若$\mathcal{F}(T,T)<S(T)$，
想买钢的公司做多头期货即可低于市价拿货。两股套利压力把到期期货价钉在
现货价上。

\origfig{3.1}{可能的期货价格路径（书页27）。两条实线为期货价
$\mathcal{F}(t,T)$的可能轨迹，虚线为现货价$S(t)$：合约存续期内期货价可在
现货价上下两侧游走，只在到期时刻$T$收敛于现货价。}

\paragraph{解读}
远期与期货都服务于同一种需求（锁价），选择标准是流动性与信用：远期合约
按买卖双方需要定制，一笔现金流、双边信用风险；期货标准化、在交易所交易、
流动性更好，但逐日盯市把一次到期结算拆成一串现金流。这串现金流使期货
价格与远期价格在利率随机时可以不同——期货定价是利率理论的试金石，原书
把它留到第9章的国债期货再处理；本节只需记住$\mathcal{F}(T,T)=S(T)$这条
边界条件。本节出现的"多头/空头""盯市""保证金"等词汇也是后文对冲论证
的日常语言。

\section{期权}

本节引入第三类、也是本书真正的主角——期权。远期与期货的卖方被合约
\emph{强制}交收资产；期权则把权利与义务分开：持有人付一笔期权费
（premium）买来的只是\emph{权利}而非义务。期权可以写在任何证券（乃至其他
衍生品）之上，市场上大量期权在场内衍生品市场交易。

\paragraph{模型设定与记号}
期权$C$是一份买或卖的合约（分别称看涨期权call与看跌期权put）。以欧式
（European）看涨期权为例：卖方有义务在预定价格$K$与预定未来时刻交割公司
股票$S$，买方则有权选择行权（exercise）或放弃行权。到期时若股价低于$K$，
看涨期权持有人自然放弃；高于$K$则行权获利。看跌期权（put）的持有人对称地
拥有按预定价格\emph{卖出}的权利。一般地，欧式期权是固定到期日的合约，
买方可在到期日按某个预定（但不必然固定）的执行价格（strike price）买入
或卖出证券；执行价格的精确形式称为期权的收益函数（payoff function）。
期权按收益是否依赖路径分两大类：路径无关与路径依赖。

\subsection{路径无关期权}

路径无关期权的收益函数只依赖标的证券在到期时刻的取值，与证券
\emph{如何}走到终点无关。最重要的两类是欧式看涨与看跌期权。设标的证券为
$S$，其上的欧式看涨期权价格记$C(t)=C(t,S(t))$，赋予持有人在$T>t$以执行
价格$K$买入证券的权利。到期日$t=T$时其价值显然为
\begin{equation*}
C(T,S(T))=
\begin{cases}
S(T)-K, & S(T)>K,\\[2pt]
0, & S(T)<K,
\end{cases}
\;=\;g(S),
\end{equation*}
其中$g(S)$即收益函数。

\origfig{3.2}{看涨期权的收益（书页28）。实线为到期收益$g(S)$：$S<K$段
贴地为零，$K$处以斜率1折起；虚线为到期前$C(t,S(t))$的可能取值，作为
$S$的函数是一条高于收益折线的光滑曲线——这段高出折线的"溢价"正是
期权费的时间价值。}

欧式看跌期权价格记$P(t)$，定义同上，只是持有人有权以$K$\emph{卖出}证券。
设即期利率（spot interest rate）$r$为常数，一个简单的无套利论证给出
\emph{看跌--看涨平价}（put--call parity）：
\begin{equation*}
C(t)+Ke^{-r(T-t)}=P(t)+S(t);\qquad t\leq T.
\end{equation*}
读法：左边"看涨期权加一笔到期本息恰为$K$的现金存款"，与右边"看跌期权
加一股股票"在到期日的价值同为$\max(S(T),K)$——两个组合今日若不等价，
买卖之间即有套利可图。欧式看涨、看跌期权由此确属路径无关：收益只看终价。

\origfig{3.3}{看跌期权的收益（书页29）。实线为到期收益$g(S)$：$S>K$段
贴地，$K$处以斜率$-1$折下；虚线为到期前期权值的可能轨迹。}

\subsection{路径依赖期权}

路径依赖期权的收益函数依赖证券在到期前的整条路径。最有名的一类是美式
（American）期权：与欧式唯一的不同是可以提前行权。它显然路径依赖——
"现在行不行权"取决于到期前全程的价格轨迹。无套利论证给出两条经典结论：
不付红利的股票上，美式看涨与欧式看涨同价；而美式看跌一般贵于欧式看跌
（提前行权锁定利息的能力有价值）。

亚式期权（Asian option）的收益取决于证券在存续期（从签约的$t$到到期的
$T$）内的\emph{平均}价格。障碍期权（barrier option）为股价预设一条屏障：
敲出（knock-out）障碍欧式看涨期权与普通欧式看涨相同，只是一旦股价$S(t)$
触及屏障，期权立即作废；对称地有敲入（knock-in）期权，以及敲入/敲出的
各种组合。更奇特的品种——回望期权（look-back）、quanto期权、篮子期权
（basket）、混合期权（hybrid）、双执行价期权（dual-strike）等——为投资者
的特殊需求定制；还有场外（over-the-counter，OTC）期权市场按需开制合约。

\paragraph{解读}
本节分类的关键在"终值问题"这四个字：路径无关期权的收益只给定$T$时刻的
终值，定价因此成为"由终值反推现值"的问题——这正是下一节随机微分方程与
3.5节Black--Scholes方程的出发点。而路径依赖期权（美式、亚式、障碍）需要
的数学更多：障碍期权将在第4章用哈密顿量精确求解，亚式期权的几何平均
恰在本章3.4节作为一个可精确计算的例子出现，随机波动率期权的路径积分
（第5章）则同时覆盖两类。

\section{随机微分方程}

本节把定价问题写成数学形式并给出股价的动力学假设。期权定价的基本问题是：
已知到期日$T$的收益函数$g(S)$，问当证券现价为$S(t)$时期权在更早的
$t<T$时刻该值多少？显然$C=C(t,S(t))$，且终值$C(T,S(T))=g(S(T))$。对路径
无关期权，这是一个\emph{终值问题}（final value problem）：终值已知，要往
回解出更早时刻的值。今天的期权价取决于证券未来的取值，所以必须先规定
$S(t)$如何演化到$S(T)$——理论金融的标准做法是把股价建模为随机过程
（stochastic process），由随机微分方程驱动。
\footnote{原书脚注指出这套方程的名号对照：同一件东西在物理文献中叫朗之万
方程（Langevin equation），在概率论文献中叫伊藤--维纳过程
（Ito--Wiener process）；白噪声的统计性质见原书附录A.4，狄拉克
$\delta$函数见附录A.2，朗之万方程的简介见附录A.5。}

\paragraph{模型设定与关键公式}
股价$S(t)$的演化方程取
\begin{equation*}
\frac{dS(t)}{dt}=\phi S(t)+\sigma S R(t)
\tag{3.2}
\end{equation*}
其中$\phi$是证券的期望收益（expected return），$\sigma$是其波动率
（volatility），$R$是零均值的高斯白噪声（Gaussian white noise）。沿用
Black--Scholes的分析，本节先取$\sigma$为常数。$\sigma$度量股价演化中
随机性的强度；特殊情形$\sigma=0$时股价确定性增长，$S(t)=e^{\phi t}S(0)$。

白噪声在任意时刻相互独立，其关联函数由狄拉克$\delta$函数给出：
\begin{equation*}
E[R(t)]=0;\qquad E[R(t)R(t')]\equiv\langle R(t)R(t')\rangle=\delta(t-t').
\tag{3.3}
\end{equation*}
随机变量$X$的期望值本书并用$E[X]$与$\langle X\rangle$两种记号。

把时间离散化为$t=n\epsilon$、$R(t)\to R_n$后，白噪声的概率分布函数为
\begin{equation*}
P(R_n)=\sqrt{\frac{\epsilon}{2\pi}}\,e^{-\frac{\epsilon}{2}R_n^{2}}.
\tag{3.4}
\end{equation*}
据此（原书附录A.4给出证明）可得白噪声最奇特的性质：
\begin{equation*}
R_n^{2}=\frac{1}{\epsilon}+\text{量级为}0(1)\text{的随机项}.
\tag{3.5}
\end{equation*}

\paragraph{推导主线与解读}
式（3.5）是本章一切"新数学"的源头，值得讲透。$R_n$是量级$1/\sqrt{\epsilon}$
的随机变量（密度里配了$\sqrt{\epsilon}$的归一因子），平方后领先项
$1/\epsilon$却是\emph{确定性}的——随机性只体现在次领头阶。换言之，到
$\epsilon$的领先阶，白噪声的平方不再是随机量。这个看似无害的事实就是
概率论中的伊藤微积分：任何"白噪声的函数"（如$S(t)$、$C(t)$）的泰勒展开
都会多出经典微积分没有的项，下一节系统地展开这一点。

\section{伊藤微积分}

本节解决的问题是：白噪声的奇异性如何改写泰勒展开与乘积求导法则。设$f$是
白噪声$R(t)$的任意函数（例如$f=C(t,S(t))$）。按导数定义
\begin{equation*}
\frac{df}{dt}=\lim_{\epsilon\to 0}\frac{f(t+\epsilon,S(t+\epsilon))-f(t,S(t))}{\epsilon},
\end{equation*}
做泰勒展开到二阶，得
\begin{equation*}
\frac{df}{dt}=\frac{\partial f}{\partial t}
+\frac{\partial f}{\partial S}\frac{dS}{dt}
+\frac{\epsilon}{2}\frac{\partial^{2}f}{\partial S^{2}}
\Bigl[\frac{dS}{dt}\Bigr]^{2}+0(\epsilon^{1/2}).
\tag{3.6}
\end{equation*}
对光滑函数，末项是$\epsilon$阶、取极限后消失。但式（3.2）中的$dS/dt$
含白噪声，不光滑：
\begin{align*}
\Bigl[\frac{dS}{dt}\Bigr]^{2}
&=\sigma^{2}S^{2}R^{2}+0(1)\\
&=\frac{1}{\epsilon}\,\sigma^{2}S^{2}+0(1).
\tag{3.7}
\end{align*}
$\epsilon\times[\,dS/dt\,]^2$因而不趋于零，而趋于$\sigma^{2}S^{2}$——
二阶项\emph{存活}下来。把式（3.2）、（3.6）、（3.7）合并，取
$\epsilon\to 0$得
\begin{align*}
\frac{df}{dt}
&=\frac{\partial f}{\partial t}
+\frac{\partial f}{\partial S}\frac{dS}{dt}
+\frac{\sigma^{2}}{2}S^{2}\frac{\partial^{2}f}{\partial S^{2}}\\
&=\frac{\partial f}{\partial t}
+\frac{1}{2}\sigma^{2}S^{2}\frac{\partial^{2}f}{\partial S^{2}}
+\phi S\frac{\partial f}{\partial S}
+\sigma S\frac{\partial f}{\partial S}R.
\tag{3.8}
\end{align*}
第一行是普通链式法则多出的伊藤修正项$\frac{\sigma^{2}}{2}S^{2}
\partial^{2}f/\partial S^{2}$；第二行把$dS/dt$按式（3.2）展开，确定性
部分与随机部分（正比于$R$的末项）分明。

乘积的求导法则同样被改写。设$g(t,R(t))\equiv g_t$是另一个白噪声函数，
记$\delta g_t\equiv g_{t+\epsilon}-g_t$，则
\begin{align*}
\frac{d(fg)}{dt}
&=\lim_{\epsilon\to 0}\frac{1}{\epsilon}\bigl[f_{t+\epsilon}g_{t+\epsilon}-f_tg_t\bigr]\\
&=\lim_{\epsilon\to 0}\frac{1}{\epsilon}\bigl[\delta f_t\,g_t+f_t\,\delta g_t+\delta f_t\,\delta g_t\bigr],
\end{align*}
通常末项$\delta f_t\delta g_t$是$\epsilon^{2}$阶而消失；白噪声的奇异性
却让它存活，于是
\begin{equation*}
\frac{d(fg)}{dt}=\frac{df}{dt}\,g+f\,\frac{dg}{dt}
+\frac{df}{\sqrt{dt}}\,\frac{dg}{\sqrt{dt}}
\;:\;\text{伊藤链式法则（Ito's chain rule）}.
\tag{3.9}
\end{equation*}
比经典乘积法则多出的第三项，与式（3.8）的二阶项同源。

\paragraph{微分形式的复推}
式（3.8）是衍生证券理论的中心结果，原书又用微分形式把它重推一遍。把
式（3.2）写成微分形式
\begin{equation*}
dS=\phi S\,dt+\sigma S\,dz;\qquad dz=R\,dt,
\tag{3.10}
\end{equation*}
其中$dz$是维纳过程（Wiener process）。由式（3.5），$R^{2}(t)=1/dt$，故
\begin{equation*}
(dz)^{2}=R_t^{2}(dt)^{2}=dt+0(dt^{3/2}),
\end{equation*}
从而
\begin{equation*}
(dS)^{2}=\sigma^{2}S^{2}\,dt+0(dt^{3/2}).
\end{equation*}
于是$f$的微分展开为
\begin{align*}
df&=\frac{\partial f}{\partial t}\,dt+\frac{\partial f}{\partial S}\,dS
+\frac{1}{2}\frac{\partial^{2}f}{\partial S^{2}}(dS)^{2}+0(dt^{3/2})\\
&=\Bigl(\frac{\partial f}{\partial t}
+\frac{1}{2}\sigma^{2}S^{2}\frac{\partial^{2}f}{\partial S^{2}}\Bigr)dt
+\sigma S\frac{\partial f}{\partial S}\,dz
+\phi S\frac{\partial f}{\partial S}\,dt,
\end{align*}
用$dz/dt=R$即回到式（3.8）。微分形式下的链式法则与式（3.9）对应：
\begin{equation*}
d(fg)=df\,g+f\,dg+df\,dg.
\end{equation*}
\footnote{这条"$dS$的平方等于$d t$"规则在通行教材（如Shreve）中表述为
"布朗运动的二次变差（quadratic variation）为$t$"或"伊藤修正项"；物理
读者可把它读作高斯白噪声关联$\langle R R\rangle=\delta$的直接推论。
读推导时抓住一句口诀即可：凡$(d\cdots)^{2}$出现，保留$d t$阶，其余丢弃。}

\subsection*{股票价格：对数正态随机变量}

为了演示随机微积分的用法（也为3.5.2节做准备），原书把式（3.2）积分出来。
取对数变量并求导——注意这里\emph{必须}用伊藤版链式法则（3.8），多出的
$-\sigma^{2}/2$正是伊藤修正项：
\begin{equation*}
x(t)=\ln[S(t)];\qquad\Rightarrow\quad
\frac{dx}{dt}=\phi-\frac{\sigma^{2}}{2}+\sigma R(t).
\tag{3.12}
\end{equation*}
关键在于：$R$的系数是常数，$x$方程是\emph{线性}的，逐点积分即得
\begin{equation*}
\Rightarrow\quad
x(T)=x(t)+\Bigl(\phi-\frac{\sigma^{2}}{2}\Bigr)(T-t)
+\sigma\int_{t}^{T}dt'\,R(t').
\tag{3.13}
\end{equation*}
随机变量$\int_t^T dt' R(t')$是一列正态随机变量之和，由原书式（A.29）知它
是$N(0,\sqrt{T-t})$随机变量，于是
\begin{equation*}
S(T)=S(t)\,e^{(\phi-\frac{\sigma^{2}}{2})(T-t)+(\sigma\sqrt{T-t})Z},
\qquad Z=N(0,1).
\tag{3.14}
\end{equation*}
股价从给定的$S(t)$出发，到期日散布成一族可能的$S(T)$；指数上是正态
（高斯）随机变量，故$S(T)$是对数正态（lognormal）随机变量。把证券价格
建模为对数正态变量的经验效度，原书引Campbell等的实证研究讨论。

\subsection*{股票价格的几何平均}

路径依赖量并非都不可解：股价的几何平均（geometric mean）的概率分布可以
\emph{精确}算出——这个例子同时为3.9节附录预置了工具。记$\tau=T-t$、
$m=x(t)+\frac{1}{2}(\phi-\frac{\sigma^{2}}{2})\tau$，由式（3.13）
\begin{align*}
&S_{\text{geometric mean}}=e^{G},\\
&G\equiv\frac{1}{\tau}\int_{t}^{T}dt'\,x(t')
=m+\frac{\sigma}{\tau}\int_{t}^{T}dt'\int_{t}^{t'}dt''\,R(t'')\\
&\phantom{G\equiv}=m+\frac{\sigma}{\tau}\int_{t}^{T}dt'\,(T-t')\,R(t').
\end{align*}
（末步交换积分次序：$\int_t^T dt'\int_t^{t'}dt''$换成权重$T-t'$的单一
积分。）白噪声的积分是高斯随机变量，由均值与方差完全决定。均值
$E[G]=m$；方差用式（3.3）的$\delta$关联逐点落网：
\begin{align*}
E[(G-m)^{2}]
&=\Bigl(\frac{\sigma}{\tau}\Bigr)^{2}\int_{t}^{T}dt'\,(T-t')
\int_{t}^{T}dt''\,(T-t'')\,E[R(t')R(t'')]\\
&=\Bigl(\frac{\sigma}{\tau}\Bigr)^{2}\int_{t}^{T}dt'\,(T-t')^{2}
=\frac{\sigma^{2}\tau}{3}.
\end{align*}
因此
\begin{equation*}
G=N\Bigl(m,\;\frac{\sigma^{2}\tau}{3}\Bigr).
\tag{3.15}
\end{equation*}
\paragraph{解读}
几何平均是对数正态的，均值与股价本身相同，方差却是股价的$1/3$——因子
$1/3$来自梯形权重$\int_t^T (T-t')^{2}dt'=\tau^{3}/3$：越靠近起点的噪声
被平均掉的路径越多。3.9节附录里"平均波动率的分布
$P_M=N(V_0,\xi^{2}\tau/3)$"正是同一计算换个马甲。

\section{Black--Scholes方程：对冲组合}

本节回答全书的核心问题之一：期权价格$C$到底由什么方程决定？约束只有
无套利一条，单凭它推不出$C$。Black与Scholes的根本洞见是：若能把期权
\emph{完全对冲}（perfectly hedge），就能为它定价——完全对冲的组合没有
不确定性，其收益率必须等于即期利率$r$给出的无风险收益，否则即可套利。
要构造这样的组合，必须先分析期权随时间的演化。

\paragraph{推导主线}
Black--Scholes的基本想法：构造一个组合，使其价值的瞬时变化与白噪声$R$
无关——没有随机性，即为完全对冲。取组合
\begin{equation*}
\Pi=C-\frac{\partial C}{\partial S}\,S,
\tag{3.16}
\end{equation*}
即持有一份期权$C$，同时卖空$\partial C/\partial S$份证券$S$。对$\Pi$用
式（3.8）（对$C$）与式（3.2）（对$S$）：
\begin{align*}
\frac{d\Pi}{dt}
&=\frac{dC}{dt}-\frac{\partial C}{\partial S}\frac{dS}{dt}\\
&=\frac{\partial C}{\partial t}
+\frac{1}{2}\sigma^{2}S^{2}\frac{\partial^{2}C}{\partial S^{2}}.
\tag{3.17}
\end{align*}
选$\partial C/\partial S$份$S$做对冲的用意正在于此：式（3.8）末尾的随机项
$\sigma S(\partial C/\partial S)R$与$dS/dt$里的随机项$\sigma SR$系数相同，
相减即消。$t$时刻$S(t)$已知，$C(t,S(t))$是确定性量，故$d\Pi/dt$中不再
含$R$——用一种证券的随机涨落抵消另一种证券的随机涨落，就是对冲（delta
对冲）的实质。
\footnote{原书脚注声明：混合二阶导项$(\partial^{2}C/\partial S\partial t)S$
被当作可忽略；这在$\sigma$为常数的模型里是自洽的近似。}

组合$\Pi$的收益率现在是确定性的，按无套利原理必须等于短期无风险利率$r$：
\begin{equation*}
\frac{d\Pi}{dt}=r\Pi.
\end{equation*}
把式（3.17）代入，得到著名的Black--Scholes方程：
\begin{equation*}
\frac{\partial C}{\partial t}
+rS\frac{\partial C}{\partial S}
+\frac{1}{2}\sigma^{2}S^{2}\frac{\partial^{2}C}{\partial S^{2}}
=rC.
\tag{3.18}
\end{equation*}

\paragraph{解读}
两点最值得咀嚼。其一，式（3.2）的期望收益参数$\phi$在式（3.18）中
\emph{消失}了：风险中性组合的价值与投资者的期望（$\phi$所体现者）无关，
等价地说，衍生品定价基于一个与投资者风险偏好无关的无风险过程——这就是
风险中性（risk-neutral）定价的萌芽。对冲实际上意味着：对被对冲的组合，
存在证券$S(t)$的一种风险中性演化，也称风险中性测度（risk-neutral
measure）或无风险测度。其二，"完全对冲"在实践中做不到：Black--Scholes
框架悬于这个理想化概念之上。原书随即点出后文伏笔——Bouchaud与Potters用
\emph{不完美}对冲组合为期权定价，残余风险（residual risk）恒为有限；
第9章的国债对冲场论同样留下一个消不掉的有限风险，称残余方差（residual
variance）。第4、5章将分别用哈密顿量与路径积分重解式（3.18）。

\subsection{Black--Scholes推导中的假设}

推导实际动用了六条假设，逐条列出（这是日后一切推广的突破口）：
\begin{itemize}
\item 组合满足无套利条件；
\item 股票无限可分割、可以卖空（否则式（3.16）的非整数头寸无从谈起）；
\item 股价连续演化——若允许不连续跳跃，组合无法完全对冲，Black--Scholes
分析失效；
\item 即期利率$r$为常数（可推广为随机即期利率）；
\item 组合$\Pi$可以连续再平衡（对冲比$\partial C/\partial S$须随时调整）；
\item 无交易成本。
\end{itemize}
这些条件在真实市场中并不完全成立（交易成本尤其显著），但市场仍以
Black--Scholes定价为行业标准，并将其作为更复杂期权定价的基础。

\subsection{Black--Scholes方程的风险中性解}

解式（3.18）的路子有两条：直接当偏微分方程解；或用哈密顿量与路径积分
（原书4.6.1节与5.2节）。本节演示最优雅的一条——风险中性估值
（risk-neutral valuation）原理，这条原理在第9章定价国债期权时还要大规模
使用。记剩余时间$\tau=T-t$。式（3.14）现在读作
\begin{equation*}
S(T)=e^{x(T)};\qquad
x(T)=N\Bigl(\ln S(t)+\Bigl(\phi-\frac{\sigma^{2}}{2}\Bigr)\tau,\;
\sigma\sqrt{\tau}\Bigr).
\tag{3.19}
\end{equation*}
风险中性估值原理说：欧式看涨期权的现值等于其到期期望值
$E[\max(S-K,0)]$按无风险利率贴现。这个"风险中性概率分布"又称鞅测度
（martingale measure），它须满足鞅条件（原书式A.40）：
\begin{equation*}
S(t)=E\bigl[e^{-r\tau}S(T)\,\big|\,S(t)\bigr]
=\int_{-\infty}^{+\infty}e^{x}\,P_m(x)\,dx.
\tag{3.20}
\end{equation*}
（含义：贴现后的股价在风险中性测度下是鞅——未来贴现值的条件期望等于
现值。）用式（3.19）的对数正态分布解式（3.20），得$\phi=r$：期望收益被
无风险利率替换。于是鞅概率分布为（记$\ln S(t)=x(t)$、$x(T)=x$）
\begin{equation*}
P_m(x)=\frac{1}{\sqrt{2\pi\sigma^{2}\tau}}\,
e^{-\frac{1}{2\sigma^{2}\tau}\bigl[x-x(t)-(r-\frac{\sigma^{2}}{2})\tau\bigr]^{2}}.
\tag{3.21}
\end{equation*}
与式（3.19）对照，唯一的改动是$\phi$换成了$r$——这正是风险中性估值的
操作内容。期权现价即到期收益在鞅测度下的期望经贴现：
\begin{equation*}
C=e^{-r\tau}E[\max(S-K,0)]
=e^{-r\tau}\int_{\ln K}^{+\infty}(e^{x}-K)\,P_m(x)\,dx.
\tag{3.22}
\end{equation*}
积分（3.22）的显式值在下面式（3.25）算出，欧式看涨期权价格为
\begin{equation*}
C(\tau,S,K,r)=S\,N(d_{+})-Ke^{-r\tau}N(d_{-}),
\tag{3.23}
\end{equation*}
其中$N(x)$是正态随机变量的累积分布（原书式A.22）：
\begin{equation*}
N(x)=\frac{1}{\sqrt{2\pi}}\int_{-\infty}^{x}e^{-\frac{1}{2}z^{2}}\,dz;
\qquad
d_{\pm}=\frac{\ln\bigl(\frac{S}{K}\bigr)+\bigl(r\pm\frac{\sigma^{2}}{2}\bigr)\tau}
{\sigma\sqrt{\tau}}.
\tag{3.24}
\end{equation*}

\paragraph{解读：公式与市场的两处裂缝}
式（3.23）与真实报价吻合得并不好。其一，原则上$\sigma$应与执行价格$K$
无关，数据却非如此。于是从业者反其道而行：对每个执行价$K$，用观测到的
期权价反解出使式（3.23）成立的波动率，得隐含波动率（implied volatility）
$\sigma(K)$——它不是常数，而恰是交易员的报价语言与策略路标。其二，高斯
分布（式3.21）下股价对数偏离超过$\sigma\sqrt{T-t}$的概率急剧衰减，实际
观测的偏离却大得多，即肥尾（fat tail）现象；文献中已提出列维（Levy）
分布等非高斯分布来解释。这两条裂缝是后文引入更复杂模型的实证动机。

\subsection*{欧式看涨期权的Black--Scholes价格}

因为式（3.23）用途极广，原书给出积分（3.22）的显式推导。记
$S(T)=e^{x}$、$x_0=\ln S+\tau(r-\frac{\sigma^{2}}{2})$。由式（3.22）与
式（3.21），
\begin{align*}
C&=e^{-r\tau}\int_{-\infty}^{+\infty}\frac{dx}{\sqrt{2\pi\tau\sigma^{2}}}
\,(e^{x}-K)_{+}\,e^{-\frac{1}{2\tau\sigma^{2}}(x-x_{0})^{2}}\\
&=e^{-r\tau}\int_{-\infty}^{+\infty}\frac{dx}{\sqrt{2\pi\tau\sigma^{2}}}
\,(e^{x+x_{0}}-K)_{+}\,e^{-\frac{1}{2\tau\sigma^{2}}x^{2}}\\
&=e^{-r\tau}\int_{\ln K-x_{0}}^{+\infty}\frac{dx}{\sqrt{2\pi\tau\sigma^{2}}}
\,(e^{x+x_{0}}-K)\,e^{-\frac{1}{2\tau\sigma^{2}}x^{2}}\\
&=S\left[\int_{\ln K-x_{0}}^{+\infty}\frac{dx}{\sqrt{2\pi\tau\sigma^{2}}}
\,e^{-\frac{1}{2\tau\sigma^{2}}(x+\tau\sigma^{2})^{2}}\right]
-e^{-r\tau}K\,N(d_{-})\\
&=S\,N(d_{+})-e^{-r\tau}K\,N(d_{-}).
\tag{3.25}
\end{align*}
\footnote{排印勘误：第4行方括号内指数原书印作$(x+\tau\sigma^{2})^{2}$；
按配方方向应为$(x-\tau\sigma^{2})^{2}$（把$e^{x-\frac{x^{2}}{2\tau\sigma^{2}}}
\cdot e^{-\tau\sigma^{2}/2}$配方，再作代换$u=x-\tau\sigma^{2}$），否则末行
的$SN(d_{+})$无法得出。导读按原书照录，在此注明。}
推导的骨架值得复述：第一步写下对数正态期望（$(\cdot)_{+}$表示取正部）；
第二步平移积分变量消去$x_0$；第三步把取正部换成积分下限$\ln K-x_0$；
第四步把$e^{x}$项配方化成高斯（$e^{-r\tau}e^{x_0}=Se^{-\tau\sigma^{2}/2}$
给出前因子$S$），第二项认出累积分布$N(d_-)$；第五步配方后的高斯积分
正是$N(d_+)$。

\section{带随机波动率的股票价格}

本节放松Black--Scholes的核心假设：波动率不再是常数，而是一个独立随机的
对象。想法上仍是老一套——构造（无风险的）对冲组合为期权定价；障碍在于
波动率本身不是资本市场上的交易品种，可用的金融工具不够把波动率的风险
完全对冲掉。市场因此是\emph{不完备}的（incomplete）：演化随机波动率的
风险中性测度不再唯一，其演化依赖于投资者的风险偏好。这是全章论证链上
第一个真正的裂缝，也是第9、10章场论对冲与非线性模型登场的伏笔。

\paragraph{文献中的波动率过程}
记$V\equiv\sigma^{2}$。文献中提出过多种方差$V$的随机过程。Hull--White、
Heston等取
\begin{equation*}
\frac{dV}{dt}=a+bV+\xi V^{1/2}Q;\qquad \sigma^{2}\equiv V,
\end{equation*}
其中$Q$是白噪声；Baaquie、Hull--White等取
\begin{equation*}
\frac{dV}{dt}=\mu V+\xi VQ;
\end{equation*}
Stein--Stein则取（对$\sigma$本身）
\begin{equation*}
d\sigma=-\delta(\sigma-\theta)\,dt+k\,dz,
\tag{3.26}
\end{equation*}
$\delta$与$\theta$分别是均值回复（mean reversion）强度与波动率的均值。
除式（3.26）外，上述各过程都是如下一般形式的特例
\footnote{原书脚注指出：式（3.26）也能纳入，只要在漂移项中添一项
$\gamma V^{1/2}$。}：
\begin{equation*}
\frac{dV}{dt}=\lambda+\mu V+\xi V^{\alpha}Q.
\tag{3.27}
\end{equation*}
$\lambda$与$\mu$的选取受$V>0$的约束。

\paragraph{完整耦合过程与广义伊藤公式}
股价与方差由两条白噪声驱动，构成完备的耦合系统：
\begin{equation*}
\frac{dS}{dt}=\phi S\,dt+S\sqrt{V}\,R_{1},
\tag{3.28}
\end{equation*}
\begin{equation*}
\frac{dV}{dt}=\lambda+\mu V+\xi V^{\alpha}R_{2};\qquad V\equiv\sigma^{2}.
\tag{3.29}
\end{equation*}
注意原书式（3.28）左端写作$dS/dt$而右端保留$dt$，是记号上的将就，读作
微分式即可。$R_1$、$R_2$是相关系数为$\rho$（$-1\leq\rho\leq 1$）的高斯
白噪声：
\begin{equation*}
\langle R_{1}(t)R_{1}(t')\rangle=\langle R_{2}(t)R_{2}(t')\rangle
=\delta(t-t')=\frac{1}{\rho}\,\langle R_{1}(t)R_{2}(t')\rangle,
\tag{3.30}
\end{equation*}
即两噪声的互关联为$\rho\,\delta(t-t')$；$\phi,\lambda,\mu,\xi$取常数。
离散化到$\epsilon$的领头阶，白噪声的平方规则带上相关项：
\begin{equation*}
R_{1}^{2}(t)=\frac{1}{\epsilon}=R_{2}^{2}(t);\qquad
R_{1}R_{2}(t)=\frac{\rho}{\epsilon}.
\tag{3.31}
\end{equation*}
把式（3.8）推广到双噪声情形——对$R_1,R_2$的任意函数$f$，二阶项中除两个
各自的平方外还多出互相关项——得
\begin{align*}
\frac{df}{dt}
&=\frac{\partial f}{\partial t}
+\phi S\frac{\partial f}{\partial S}
+(\lambda+\mu V)\frac{\partial f}{\partial V}
+\frac{\sigma^{2}S^{2}}{2}\frac{\partial^{2}f}{\partial S^{2}}
+\rho V^{1/2+\alpha}\xi\,\frac{\partial^{2}f}{\partial S\,\partial V}\\
&\quad+\frac{\xi^{2}V^{2\alpha}}{2}\frac{\partial^{2}f}{\partial V^{2}}
+\sigma S\frac{\partial f}{\partial S}R_{1}
+\xi V^{\alpha}\frac{\partial f}{\partial V}R_{2}\\
&=\Theta+\Xi R_{1}+\Psi R_{2},
\tag{3.32}
\end{align*}
末行把确定性部分记$\Theta$、两个随机系数分别记$\Xi$与$\Psi$，为的是把
随机与非随机项分开，供下一步对冲消元。

\section{Merton--Garman方程}

现在只有一个随机方程（3.32）、两个随机源，一份期权只提供一个对冲工具，
消不掉两个随机项——这正是"市场不完备"在代数上的样子。Merton与Garman的
办法：同时取\emph{两种}期权$C_1$与$C_2$（同一标的证券，执行价与到期日
分别为$K_1,K_2$与$T_1,T_2$），加上股票本身，凑出足够多的对冲工具。取组合
\begin{equation*}
\Pi=C_{1}+\Gamma_{1}C_{2}+\Gamma_{2}S,
\end{equation*}
对每项用式（3.32）：
\begin{equation*}
\frac{d\Pi}{dt}
=\Theta_{1}+\Gamma_{1}\Theta_{2}+\Gamma_{2}\phi
+(\Xi_{1}+\Gamma_{1}\Xi_{2}+\Gamma_{2}\sigma S)R_{1}
+(\Psi_{1}+\Gamma_{1}\Psi_{2})R_{2}.
\end{equation*}
与常数波动率情形一样，完全对冲要求组合价值的变化不含随机项。两个随机
系数置零，$\Gamma_1$、$\Gamma_2$两个未知数恰可解出：
\begin{align*}
&\Xi_{1}+\Gamma_{1}\Xi_{2}+\Gamma_{2}\sigma S=0,\\
&\Psi_{1}+\Gamma_{1}\Psi_{2}=0,
\end{align*}
即
\begin{align*}
\Gamma_{1}&=-\frac{\Psi_{1}}{\Psi_{2}}
=-\frac{\partial C_{1}/\partial V}{\partial C_{2}/\partial V},\\
\Gamma_{2}&=\frac{\Psi_{1}}{\Psi_{2}}\frac{\partial C_{2}}{\partial S}
-\frac{\partial C_{1}}{\partial S}
=\frac{\partial C_{1}/\partial V}{\partial C_{2}/\partial V}
\,\frac{\partial C_{2}}{\partial S}
-\frac{\partial C_{1}}{\partial S}.
\end{align*}
组合既已无风险，无套利原理要求它按无风险利率增值：
$\frac{d\Pi}{dt}=r\Pi$。代入$\Theta,\Xi,\Psi$的表达式并整理——关键是
做\emph{变量分离}：左边只含$K_1,T_1$，右边只含$K_2,T_2$，两边欲恒等必须
等于一个与$K_1,K_2,T_1,T_2$都无关的量$\beta(S,V,t,r)$：
\begin{align*}
&\frac{1}{\partial C_{1}/\partial V}\biggl(
\frac{\partial C_{1}}{\partial t}
+(\lambda+\mu V)\frac{\partial C_{1}}{\partial V}
+rS\frac{\partial C_{1}}{\partial S}
+\frac{VS^{2}}{2}\frac{\partial^{2}C_{1}}{\partial S^{2}}\\
&\qquad\qquad+\rho V^{1/2+\alpha}\xi\,\frac{\partial^{2}C_{1}}{\partial S\,\partial V}
+\frac{\xi^{2}V^{2\alpha}}{2}\frac{\partial^{2}C_{1}}{\partial V^{2}}
-rC_{1}\biggr)\\
&=\frac{1}{\partial C_{2}/\partial V}\biggl(
\frac{\partial C_{2}}{\partial t}
+(\lambda+\mu V)\frac{\partial C_{2}}{\partial V}
+rS\frac{\partial C_{2}}{\partial S}
+\frac{VS^{2}}{2}\frac{\partial^{2}C_{2}}{\partial S^{2}}\\
&\qquad\qquad+\rho V^{1/2+\alpha}\xi\,\frac{\partial^{2}C_{2}}{\partial S\,\partial V}
+\frac{\xi^{2}V^{2\alpha}}{2}\frac{\partial^{2}C_{2}}{\partial V^{2}}
-rC_{2}\biggr)\\
&\equiv\beta(S,V,t,r).
\tag{3.33}
\end{align*}
$\beta$称波动率风险的市场价格（market price of volatility risk）：
$\beta$越高，投资者对波动率风险的厌恶越强。它之所以在随机波动率下
必须出现、而在Black--Scholes定价中不需要，根源仍是"波动率不可交易"：
没有足够的工具把波动率风险对冲干净，投资者的风险偏好就必须进场——
风险中性估值对波动率不再直接可用，或者说不再唯一。$\beta$难以实证估计，
但有证据表明它非零（原书引文献）；在Cox--Ingersoll--Ross模型中，消费
增长与标的资产收益有常数相关，产生正比于波动率的风险溢价——采用该模型
的好处是它只是把上式中的$\lambda$重新定义了一遍。此后原书约定：风险的
市场价格已被吸收进$\lambda$，Merton--Garman方程即
\begin{align*}
\frac{\partial C}{\partial t}
+rS\frac{\partial C}{\partial S}
+(\lambda+\mu V)\frac{\partial C}{\partial V}
+\frac{1}{2}VS^{2}\frac{\partial^{2}C}{\partial S^{2}}
+\rho\xi V^{1/2+\alpha}S\frac{\partial^{2}C}{\partial S\,\partial V}
+\xi^{2}V^{2\alpha}\frac{\partial^{2}C}{\partial V^{2}}
=rC.
\tag{3.34}
\end{align*}

\paragraph{解读}
与式（3.18）对照：Black--Scholes方程是"时间一维+股价一维"的二阶抛物型
方程；Merton--Garman方程把变量扩成$(S,V)$二维、多出$\partial_V$一阶项、
$S$--$V$交叉项与$\partial_V^2$项，且漂移$\lambda+\mu V$里藏着待定的风险
偏好。第4章将把式（3.18）与（3.34）分别改写成哈密顿量$H_{BS}$与$H_{MG}$
的作用，第5章用路径积分解$H_{MG}$版本——那是原书方法论的第一次真正亮相。

\section{小结}

本章沿"工具\,--\,方程"两条线推进。工具线：远期/期货/期权的分类与收益
函数（3.1--3.2节）；股价的朗之万型随机微分方程（3.3节）；伊藤微积分的
两条新规则——泰勒展开中的二阶存活项与乘积法则中的$df\,dg$项（3.4节）。
方程线：用对冲组合+无套利导出Black--Scholes方程（3.5节）并经风险中性
估值（鞅条件$\phi=r$）解出欧式看涨期权价格；让波动率随机化后，市场不完备
迫使引入波动率风险的市场价格$\beta$，得到Merton--Garman方程（3.6--3.7
节）。原书把本章定位为"用朗之万形式"推导——即随机微分方程语言；同样的
内容从第4章起将换用哈密顿量与路径积分语言重演，届时式（3.18）与式
（3.34）分别对应$H_{BS}$与$H_{MG}$，鞅条件将改写为哈密顿量的本征值
条件（第4章），随机波动率的解将以路径积分形式独立重推（第5章式5.115）。

\section{附录：$\rho=0$时随机波动率的解}

本附录处理式（3.28）--（3.29）体系的一个可精确求解的角落。虽然股价过程
（3.28）依赖方差$V$，但当$\rho=0$时方差过程（3.29）不依赖股价$S$——
波动率的演化可以与股价\emph{分开}处理。

\paragraph{Merton定理}
Merton的一个定理说：若波动率独立于$S$地遵循某个（可以随机的）过程
$V(t)$，则期权价格就是Black--Scholes价格、但其中的波动率换成
\emph{平均波动率}，再对该平均值的分布求积。记$\tau=T-t$、
\begin{equation*}
\bar{V}=\frac{1}{\tau}\int_{t}^{T}dt'\,V(t'),
\end{equation*}
期权价格由式（3.23）推广为
\begin{equation*}
C=\int_{0}^{\infty}\bigl[S\,N(d_{+}(\bar{V}))
-Ke^{-r\tau}N(d_{-}(\bar{V}))\bigr]\,
P_{M}(\bar{V})\,\frac{d\bar{V}}{\bar{V}},
\tag{3.35}
\end{equation*}
其中$P_M$是平均波动率$\bar{V}$的概率分布函数，$d_{\pm}(\bar{V})$仿照
式（3.24）为
\begin{equation*}
d_{\pm}=\frac{\ln(S/K)+\tau\,(r\pm\frac{1}{2}\bar{V})}{\sqrt{\bar{V}\tau}},
\end{equation*}
$K$为执行价格。式（3.35）的独立推导（连同$P_M$的形式表达式）将在第5章
用随机波动率期权定价的路径积分表述给出（原书式5.115）。
\footnote{积分测度$d\bar{V}/\bar{V}$值得注意：被积的是\emph{对数}平均
波动率的 Black--Scholes 核，与股价的对数正态性一脉相承；这层结构在第5章
的路径积分推导中会看得很清楚。}

\paragraph{两个例子}
原书举两个例子演示式（3.35）的用法（都取$\tau=T$）。其一是确定性过程：
\begin{equation*}
V=V_{0}e^{\mu t},\qquad 0\leq t\leq T.
\end{equation*}
此时平均波动率的分布退化为狄拉克$\delta$函数：
\begin{equation*}
V_{M}=\delta\Bigl(V-V_{0}\,\frac{e^{\mu\tau}-1}{\mu\tau}\Bigr),
\end{equation*}
即Black--Scholes结果中把$\sigma$换成
$\sqrt{V_0\,(e^{\mu\tau}-1)/(\mu\tau)}$——时间平均的确定性方差。

其二是随机波动率：在式（3.29）中取$\lambda=\mu=\alpha=0$，
\begin{equation*}
\frac{dV}{dt}=\xi R(t),\qquad V(0)=V_{0},\qquad 0\leq t\leq T,
\end{equation*}
$R(t)$为白噪声。
\footnote{原书脚注坦言这不是现实的过程：高斯分布给出$P(V<0)>0$，而$V$
显然非负；但只要时段较短、$P(V<0)$可忽略，它不失为合理近似。$V>0$的
约束正是3.6节说$\lambda,\mu$受限的原因，也是第7章非线性（对数）模型
的动机之一。}
方差$V$的时段均值$\bar V$在$(0,T)$内的分布直接引用式（3.15）的三分之一
因子：
\begin{equation*}
P_{M}=N\Bigl(V_{0},\;\frac{\xi^{2}\tau}{3}\Bigr).
\end{equation*}
代入式（3.35），期权价格为
\begin{equation*}
f=\sqrt{\frac{3}{2\pi\xi^{2}\tau}}
\int_{0}^{\infty}\bigl[S\,N(d_{+}(V))-Ke^{-r\tau}N(d_{-}(V))\bigr]\,
\exp\Bigl(-\frac{3(V-V_{0})^{2}}{2\xi^{2}\tau}\Bigr)\,\frac{dV}{V}.
\end{equation*}

\paragraph{解读}
附录的价值在于示范"可解结构"的叠加方式：$\rho=0$让$V$与$S$解耦，
Merton定理把随机波动率问题化归为"对平均波动率求平均的Black--Scholes"，
而平均值的分布又恰好是3.4节算过的几何平均型高斯（三分之一因子）。三个
毫不相干的部件咬合成一个显式积分。注意最终价格记$f$而非$C$——它是"对
$\bar V$平均后的期权价格"。第5章的路径积分将不借助Merton定理直接推出
式（5.115），两相对照可以看出路径积分方法的威力。

%TAGS: 3.2,3.3,3.4,3.5,3.6,3.7,3.8,3.9,3.10,3.12,3.13,3.14,3.15,3.16,3.17,3.18,3.19,3.20,3.21,3.22,3.23,3.24,3.25,3.26,3.27,3.28,3.29,3.30,3.31,3.32,3.33,3.34,3.35（原书无(3.1)与(3.11)，编号断档系原书如此；已对照PNG逐页核验，p043--p060为正文，p061为Part II扉页，p062空白）


% ch04a Hamiltonians and stock options 4.1-4.8 (书页45-62 / p063-080)
% 覆盖说明：本块覆盖原书第4章4.1--4.8节。4.8节终止于书页62上部（p080）；
% 4.9节"哈密顿量与障碍期权"自书页62中部（p080下半）起，归下一分块（ch04b）。
% 编号公式核对：本块tag为(4.2)--(4.24)、(4.26)--(4.32)与(4.34)--(4.48)，共45条；
% 原书无(4.1)、(4.25)与(4.33)——编号断档系原书自身如此（已对照PNG逐页核验，非漏转）。
\chapter{哈密顿量与股票期权}

本章是全书方法论的转折点。第3章把期权定价归结为两个抛物型偏微分方程
（Black--Scholes方程与Merton--Garman方程），并按传统路线"解方程"；本章
换一套语言：把期权价格看成某个线性空间——态空间（state space）——里的
态矢，把定价方程逐项改写成薛定谔方程（Schr\"odinger equation）的样子
$\partial C/\partial t=HC$，其中算符$H$称为期权定价的哈密顿量
（Hamiltonian）。这不是外衣式的改写：哈密顿量形式带来两件传统方法不易
做到的事——其一，本征函数（eigenfunction）：求出$H$的全部本征函数，就能
对一大类路径无关与路径依赖期权（path-dependent option）给出显式解，
障碍期权（barrier option）不过是给本征函数加约束；其二，势
（potential）：在哈密顿量上添加势函数，可以定义新的金融工具并建模
路径依赖期权。前者在本章后半（4.6节）兑现于
Black--Scholes定价核，后者在4.8节立起框架、由4.9节（障碍期权）与第5章
（路径积分）展开。

论证的顺序是：4.1--4.3节快节奏补齐量子力学清单——态矢、态空间与完备性
方程、算符与本征值问题（重点是\emph{非厄米}情形，金融哈密顿量的常态）；
4.4节把Black--Scholes与Merton--Garman方程翻译成哈密顿量语言，是全章的
核心一步；4.5--4.6节引入定价核（pricing kernel）并证明它是演化算符
$e^{-\tau H}$的矩阵元，随后用动量基把Black--Scholes定价核解到显式；
4.7节把鞅条件（martingale condition）压缩成一条
算符方程$H|S\rangle=0$；4.8节讨论给哈密顿量加势。先修内容只有第3章的
两个定价方程；本章多处指向原书章末附录4.11--4.13（两态系统与非厄米
哈密顿量的具体演算），那里是本节抽象记号的"落地样本"。

\section{量子力学要点}

本节对应原书4.1节，任务是把后文要用的量子力学词汇清点一遍。先说清楚
为什么可以这么做：本章将逐条核对地展示，期权定价与单自由度量子系统的
数学描述完全同一——第3章的Black--Scholes方程在本章会原封不动地变成
薛定谔方程（Schr\"odinger equation）的形状。方程同构，量子力学百年积累
的线性代数机器——态空间、完备性、谱分解——就可以整体搬进金融。

\paragraph{模型设定与记号}
经典力学里，粒子在$t$时刻的位置$x_t$是$t$的确定性函数，由牛顿运动定律
给出；这与零波动率（$\sigma=0$）的股票价格演化完全同型——价格轨道没有
悬念。量子力学里粒子演化是随机的，对应非零波动率（$\sigma\neq0$）的
股票。记号上，量子力学把每个时刻的位置记作$x_t$，称之为一个自由度
（degree of freedom）：它就是"系统在该时刻所取的那个随机变量"，时间
$t$只是给这些随机变量编号的参数。把所有时刻的自由度收拢成序列
$\{x_t\}$，就是概率论里的随机过程（stochastic process）。于是"单自由度
量子系统"翻译成金融语言恰好是"一维随机过程"——一只股票。

固定时刻$t$，粒子位置$x$的概率分布由
$|\psi(t,x)|^{2}\equiv\psi^{*}(t,x)\psi(t,x)$给出（星号表示复共轭）。
函数$\psi$称为系统的态矢量（state vector），它是态空间
（state space）$\mathcal{V}$的元素——$\mathcal{V}$由允许位形上的全体
函数构成，是一个线性向量空间\footnote{原书脚注给出口径：粒子在实直线
$\Re$上运动时，$\mathcal{V}$由满足$\int_{\Re}dx\,|\psi(x)|^{2}=1$的全体
$\psi(x)$组成。}。$\mathcal{V}$的对偶态空间（dual state space）
$\mathcal{V}_{\mathrm{dual}}$由从$\mathcal{V}$到复数的全体线性映射构成，
同样是线性向量空间。

记号采用Dirac的括号记号：$\mathcal{V}$的元素记右矢（ket）
$|g\rangle$，$\mathcal{V}_{\mathrm{dual}}$的元素记左矢（bra）
$\langle p|$；左矢作用在右矢上得到标量积
$\langle p|g\rangle=\langle g|p\rangle^{*}$，是个复数。物理上可观测的量
（能量、位置等）用厄米算符（Hermitian operator）表示——即把线性向量
空间映到自身的线性映射；对厄米算符而言$\mathcal{V}$与
$\mathcal{V}_{\mathrm{dual}}$同构，左右矢的区分可以忽略。

\paragraph{两条独立要素}
量子力学的数学地基由两个互相独立的部件搭成：其一，线性向量态空间及其
对偶；其二，作用在态空间上的（线性）算符。据此，原书总结出两条原则：
量子系统由态矢量$|\psi\rangle$描述，它是态空间$\mathcal{V}$的元素；系统
的一切性质都由作用在$|\psi\rangle$上的线性算符表示。

\paragraph{解读}
本节真正的载荷只有一句话：把"期权价格"看成向量空间里的向量、把"定价
方程"看成算符演化方程，有数学结构上的依据，而非修辞。传统教材（Hull、
Shreve一系）把期权价格当作待求的二元函数直接解偏微分方程；量子口径先
把"所有可能的价格函数"收进一个线性空间，再把单个方程升格为算符演化
——这步升格的回报要到4.6节求解定价核时才真正兑现。原书还预告了附录
4.11的两态系统：只有两个可能状态的最小量子系统，完备性方程（见4.2节）
在那里有看得见摸得着的矩阵版本。

\section{态空间与完备性方程}

本节对应原书4.2节。核心问题：怎样把"$\mathcal{V}$中任何矢量都能用一组
基展开"写成一条可操作的公式？答案就是完备性方程（completeness
equation）。原书开篇即声明：这一节看似形式化，实为后文反复调用的引擎
——4.6节推定价核、4.7节推鞅条件，用的都是"插入完备性"这一招。

\paragraph{模型设定与记号}
先离散后连续。设粒子（保留量子力学的说法，读者不妨直接把它读作"对数
股价"）只能在一维格点上跳动，格距为$a$，格点$x=na$。基矢（basis
states）$|n\rangle$取"第$n$个位置为1、其余分量为0"的无穷列向量，
$n=0,\pm1,\pm2,\ldots,\pm\infty$；其对偶是行向量
$\langle n|=[\ldots\,0\;1\;0\,\ldots]$。正交归一与完备性一并写出：
\begin{align*}
&|n\rangle=
\begin{pmatrix}
\vdots\\
0\\
1\\
0\\
\vdots
\end{pmatrix}
\;:\;\text{第$n$个位置};\qquad
\langle n|=[\ldots\,0\;1\;0\,\ldots]\\
&\langle m|n\rangle=\delta_{n-m}\equiv
\begin{cases}
1, & n=m,\\
0, & n\neq m,
\end{cases}\\
&\sum_{n=-\infty}^{+\infty}|n\rangle\langle n|=\mathcal{I}
\quad\text{（完备性方程）}
\end{align*}
$\mathcal{I}$是无穷维单位矩阵。完备性又称恒等算符的分解（resolution of
the identity）：只有\emph{完备}的基矢集合才能合力拼出态空间上的恒等
算符——反之，谁能拼出$\mathcal{I}$，谁就是完备基。这条判据后文会反复
使用。

\paragraph{连续极限}
允许位置其实是实直线$\Re$上的所有点，故取$a\to0$。位置基矢与标量积
（原书引其附录式(A.13)）定义为
\begin{equation*}
|x\rangle=\lim_{a\to0}\frac{1}{\sqrt{a}}\,|n\rangle;\qquad
-\infty\leq x\leq\infty
\end{equation*}
\begin{equation*}
\langle x|x'\rangle=\delta(x-x')\equiv
\begin{cases}
\infty, & x=x',\\
0, & x\neq x',
\end{cases}
\end{equation*}
即标量积由克罗内克$\delta$升级为狄拉克$\delta$函数（Dirac delta
function）。把格点完备性取极限（求和的每项配上前因子$a$）：
\begin{equation*}
\sum_{n=-\infty}^{+\infty}|n\rangle\langle n|
\;\to\;
a\sum_{n=-\infty}^{+\infty}|x\rangle\langle x|
\quad\Longrightarrow\quad
\int_{-\infty}^{\infty}dx\,|x\rangle\langle x|=\mathcal{I}
\quad\text{（完备性方程）}
\end{equation*}
其中$\mathcal{I}$现在是（函数）态空间上的恒等算符。

\paragraph{关键公式}
再给一个不借助格点的直接验证：考察两个函数的标量积，
\begin{equation*}
\langle\psi|g\rangle\equiv\int dx\,\psi^{*}(x)g(x)
=\Bigl\langle\psi\Bigl|
\Bigl\{\int_{-\infty}^{\infty}dx\,|x\rangle\langle x|\Bigr\}
\Bigr|g\Bigr\rangle
\end{equation*}
花括号里的算符夹在$\langle\psi|$与$|g\rangle$之间确实起恒等算符的
作用。于是得到本章第一条编号公式：
\begin{equation*}
\mathcal{I}=\int_{-\infty}^{\infty}dx\,|x\rangle\langle x|
\tag{4.2}
\end{equation*}
两个粒子（位置$x$与$y$）的情形只差把积分与基矢都翻倍：
\begin{equation*}
\mathcal{I}=\int_{-\infty}^{\infty}dx\,dy\,|x,y\rangle\langle x,y|
\tag{4.3}
\end{equation*}
其中$|x,y\rangle\equiv|x\rangle\otimes|y\rangle$；更多粒子依此类推。

态矢在位置基下的"分量"即波函数：$\psi(x)=\langle x|\psi\rangle$，其对偶
分量为$\psi^{*}(x)=\langle\psi|x\rangle$。两者给出矢量的模长——把(4.2)
插进中间即得
\begin{equation*}
\langle\psi|\psi\rangle
=\Bigl\langle\psi\Bigl|\int_{-\infty}^{\infty}dx\,|x\rangle\langle x|
\Bigr|\psi\Bigr\rangle
=\int_{-\infty}^{\infty}\psi(x)^{*}\psi(x)\geq0.
\end{equation*}

\paragraph{金融转译}
把$x$读作对数股价$\ln S$，本节的抽象物都有了具体身份：$|x\rangle$是
"股价钉死在$x$"的理想化状态——4.5节的数字期权收益恰是$\delta(x-x_0)$，
正是基矢$|x_0\rangle$的金融化身；完备性方程则是"对一切可能价格求和"的
恒等式。日后每见"在被积函数里插一个$\int dx\,|x\rangle\langle x|$"，读者
应当条件反射：这是在用完备基分解矢量或算符。还应指出，(4.2)的积分限是
整条实直线；对数变量$x$天然取遍$\Re$，与"粒子在直线上运动"同框，这是
本书坚持用$x=\ln S$工作（变量替换见4.4节）的深层原因之一。

\section{算符与哈密顿量}

本节对应原书4.3节。核心问题：算符如何用"矩阵元"定义？厄米共轭如何
操作？以及——本章真正的主角——\emph{非厄米}哈密顿量的本征值问题长什么
样子？

\paragraph{模型设定与记号}
算符（operator）定义为态空间$\mathcal{V}$到自身的线性映射，是张量积
空间$\mathcal{V}\otimes\mathcal{V}_{\mathrm{dual}}$的元素；附录4.11的两态
系统里，算符就是$2\times2$矩阵。由单变量$x$的全体函数构成的态空间
$\mathcal{V}=\{\psi(x)\,|\,x\in\Re\}$（约定$\langle x|\psi\rangle=\psi(x)$）
上，算符是$N\times N$矩阵在$N\to\infty$下的推广。两个基本构件：坐标
算符$\hat{x}$（把$\psi(x)$乘上$x$）与微分算符$\partial/\partial x$（把
$\psi$映为偏导）；本节出现的所有算符都是二者的函数。

\paragraph{关键公式}
与矩阵由矩阵元$M_{ij}$完全确定同理，算符由其矩阵元确定。在Dirac记号
下，
\begin{equation*}
\hat{x}\psi(x)=x\psi(x)
\quad\Longrightarrow\quad
\langle x|\hat{x}|\psi\rangle=x\,\langle x|\psi\rangle=x\psi(x),
\end{equation*}
\begin{equation*}
\Bigl\langle x\Bigl|\frac{\partial}{\partial x}\Bigr|\psi\Bigr\rangle
=\frac{\partial\psi(x)}{\partial x};
\end{equation*}
取$|\psi\rangle=|x'\rangle$立得
\begin{equation*}
\langle x|\hat{x}|x'\rangle=x\,\langle x|x'\rangle=x\,\delta(x-x').
\end{equation*}
模仿矩阵的共轭$M^{\dagger}_{ij}=M^{*}_{ji}$，任意算符$\mathcal{O}$的厄米
共轭（Hermitian conjugation）定义为
\begin{equation*}
\langle f|\mathcal{O}^{\dagger}|g\rangle\equiv
\langle g|\mathcal{O}|f\rangle^{*}
\tag{4.4}
\end{equation*}
且与矩阵同律：$(A+B+\ldots)^{\dagger}=A^{\dagger}+B^{\dagger}+\ldots$，
$(AB\ldots)^{\dagger}=\ldots B^{\dagger}A^{\dagger}$。

\paragraph{推导主线：两个算符的厄米性}
用完备性(4.2)插入与定义(4.4)，坐标算符满足
\begin{equation*}
\langle f|\hat{x}^{\dagger}|g\rangle
=\Bigl[\int_{-\infty}^{\infty}dx\,x\,g^{*}(x)f(x)\Bigr]^{*}
=\int_{-\infty}^{\infty}dx\,x\,g(x)f^{*}(x)
=\langle f|\hat{x}|g\rangle
\quad\Longrightarrow\quad
\hat{x}^{\dagger}=\hat{x}
\quad\text{（厄米）}.
\end{equation*}
微分算符对(4.4)做一次分部积分（边界项消失）：
\begin{align*}
\Bigl\langle f\Bigl|\Bigl(\frac{\partial}{\partial x}\Bigr)^{\dagger}
\Bigr|g\Bigr\rangle
&\equiv\Bigl[\int_{-\infty}^{\infty}dx\,g^{*}(x)
\frac{\partial f(x)}{\partial x}\Bigr]^{*}
=-\int_{-\infty}^{\infty}dx\,\frac{\partial g(x)}{\partial x}f^{*}(x)\\
&=-\Bigl\langle f\Bigl|\frac{\partial}{\partial x}\Bigr|g\Bigr\rangle
\quad\Longrightarrow\quad
\Bigl(\frac{\partial}{\partial x}\Bigr)^{\dagger}
=-\frac{\partial}{\partial x}
\quad\text{（反厄米，anti-Hermitian）}
\tag{4.5}
\end{align*}
一阶导数是反厄米的；乘上$i$（即$i\,\partial/\partial x$）即成厄米算符。
而二阶导数自轭：
\begin{equation*}
\Bigl(\frac{\partial^{2}}{\partial x^{2}}\Bigr)^{\dagger}
=\frac{\partial^{2}}{\partial x^{2}}
\quad\text{（厄米）}
\tag{4.6}
\end{equation*}
(4.5)与(4.6)两条小结论是4.4节检验哈密顿量厄米性的全部装备。

\paragraph{一处必须记住的细节}
矩阵元$\langle x|\partial/\partial x|f\rangle$里的算符作用在
\emph{左边}——即作用在对偶空间上。厄米算符左右不分家，但金融里的
哈密顿量一般非厄米，"向左作用还是向右作用"结果不同（原书为此专设
脚注提醒）。这枚地雷将在4.6.1小节真实引爆。

\paragraph{哈密顿量与本征值问题}
哈密顿量（Hamiltonian）$H$是驱动系统时间演化的算符，因而也是期权定价
中最重要的算符；金融中的$H$一般非厄米。先看本征态（eigenstate）的
定义性范例：坐标算符的本征方程为
\begin{equation*}
\hat{x}\,|x\rangle=x\,|x\rangle
\tag{4.7}
\end{equation*}
$|x\rangle$在$\hat{x}$作用下只被乘上实数$x$，故称$\hat{x}$的本征态，$x$
即本征值（eigenvalue）——实数性是$\hat{x}$厄米的直接红利。

非厄米哈密顿量的本征值问题（把(4.7)与附录4.11的式(4.59)合并推广）：
本征值$E$是\emph{复}数，右、左本征函数必须分开陈述，
\begin{equation*}
H|\psi_{E}\rangle=E|\psi_{E}\rangle
\quad(E\ \text{复}),
\end{equation*}
\begin{equation*}
\langle\tilde{\psi}_{E}|H=E\langle\tilde{\psi}_{E}|
\quad\Longrightarrow\quad
H^{\dagger}|\tilde{\psi}_{E}\rangle=E^{*}|\tilde{\psi}_{E}\rangle,
\end{equation*}
\begin{equation*}
\langle\tilde{\psi}_{E}|\psi_{E'}\rangle
=\frac{1}{\mu(E)}\,\delta(E-E'),
\end{equation*}
其中波浪号标记左本征函数，$\mu(E)$是本征值$E$处的态密度（density of
states），定义为
\begin{equation*}
\mu(E)=\mathrm{trace}\,\delta(H-E)
\equiv\int_{-\infty}^{+\infty}dx\,\langle x|\delta(H-E)|x\rangle
\tag{4.8}
\end{equation*}
$E$既然是复数，左右本征函数就不再互为对偶：
$\langle\tilde{\psi}_{E}|\neq\langle\psi_{E}|$。要拼出完备性，须从本征
能量中挑出一个子集$\mathcal{D}$，使其对应的$|\psi_{E}\rangle$构成态空间
的完备基；完备性方程为
\begin{equation*}
\int_{\mathcal{D}}dE\,\mu(E)\,|\psi_{E}\rangle\langle\tilde{\psi}_{E}|
=\mathcal{I}
\tag{4.9}
\end{equation*}
写成分量形式（两端各夹位置基矢）：
\begin{equation*}
\int_{\mathcal{D}}dE\,\mu(E)\,\psi_{E}(x)\tilde{\psi}_{E}(x')
=\langle x|\mathcal{I}|x'\rangle=\delta(x-x')
\tag{4.10}
\end{equation*}

对一般的（微分）哈密顿量$H=H(x,\partial/\partial x)$，本征函数方程就是
把矩阵元读成微积分：
\begin{equation*}
\Bigl\langle x\Bigl|H\Bigl(x,\frac{\partial}{\partial x}\Bigr)
\Bigr|\psi_{E}\Bigr\rangle
=H\Bigl(x,\frac{\partial}{\partial x}\Bigr)\psi_{E}(x)=E\psi_{E}(x),
\end{equation*}
且同样谨记"向左作用"：计算
$\langle\psi|H(x,\partial/\partial x)|x\rangle$时，实际在算
$\langle x|H^{\dagger}(x,\partial/\partial x)|\psi\rangle^{*}$。

\paragraph{解读}
这份清单为两件事预备。其一，(4.9)+(4.10)的谱完备性是4.6节定价核本征
展开（式(4.31)）的原型——届时只需把$e^{-\tau H}$插进(4.9)。其二，非厄米
不是技术细节：4.4节将看到，Black--Scholes哈密顿量恰因漂移项而非厄米，
而漂移正是金融的常态；附录4.13将用非厄米哈密顿量实际解一个障碍期权
（原书脚注4的预告）。\footnote{原书脚注2值得记录：量子力学只允许模长为
1（$\langle\psi|\psi\rangle=1$）的态矢，构成希尔伯特空间；金融做不到
——股价$e^{x}$这类态矢长度无穷，故本书的态空间必须取得比希尔伯特空间
大。态矢不可归一化这一点在4.7节还会重现：$|S\rangle$将是零能量本征态，
却同样不可归一。}

\section{Black--Scholes与Merton--Garman哈密顿量}

本节对应原书4.4节，是全章的核心：把第3章的两个定价方程逐项翻译成
$\partial C/\partial t=HC$，读出$H$。翻译完成后，"解期权"就等于
"解一个（非厄米的）量子系统"。

\paragraph{第一步：常数波动率}
Black--Scholes方程（原书式(3.18)）为
\begin{equation*}
\frac{\partial C}{\partial t}
=-\frac{1}{2}\sigma^{2}S^{2}\frac{\partial^{2}C}{\partial S^{2}}
-rS\frac{\partial C}{\partial S}+rC
\tag{4.11}
\end{equation*}
做对数变量替换
\begin{equation*}
S=e^{x};\qquad -\infty\leq x\leq\infty
\end{equation*}
$S>0$自动满足，且$S^{2}\partial^{2}/\partial S^{2}$与$S\,\partial/\partial S$
在$x$变量下都化为常系数。方程获得薛定谔形状——原书称之为
Black--Scholes--Schr\"odinger方程：
\begin{equation*}
\frac{\partial C}{\partial t}=H_{BS}C
\tag{4.12}
\end{equation*}
\begin{equation*}
H_{BS}=-\frac{\sigma^{2}}{2}\frac{\partial^{2}}{\partial x^{2}}
+\Bigl(\frac{1}{2}\sigma^{2}-r\Bigr)\frac{\partial}{\partial x}+r
\tag{4.13}
\end{equation*}

\paragraph{物理词典}
把(4.13)当作一个量子系统来读：自由度是$x=\ln S$；波动率扮演"质量的
倒数"——动能项$-(\sigma^{2}/2)\partial^{2}/\partial x^{2}$对应量子力学的
$-(\hbar^{2}/2m)\partial^{2}/\partial x^{2}$；漂移项对应一个依赖"速度"
的势；期权价格$C$扮演薛定谔波函数。特别地，若$r=\sigma^{2}/2$（漂移
恰好消失），$H_{BS}$退化为"自由粒子加常数"，一个纯动能系统。

\paragraph{厄米性检查}
用(4.5)与(4.6)逐项检查(4.13)：
\begin{equation*}
H_{BS}^{\dagger}
=-\frac{\sigma^{2}}{2}\frac{\partial^{2}}{\partial x^{2}}
-\Bigl(\frac{1}{2}\sigma^{2}-r\Bigr)\frac{\partial}{\partial x}+r
\neq H_{BS}
\tag{4.14}
\end{equation*}
二阶导项自轭、常数项自轭，唯独漂移项反号——Black--Scholes哈密顿量的
非厄米性全部来自漂移。这与4.3节"反厄米的一阶导"遥相呼应：漂移在物理
词典里是实的速度依赖势，而实势必破厄米（厄米要求势取虚值）。

\paragraph{第二步：Dirac记号}
把期权价格升格为态矢$|C\rangle$，其在$x$基下的表示就是函数
$C(x)=\langle x|C\rangle$。Black--Scholes方程写成
\begin{equation*}
\langle x|H_{BS}|C\rangle
=H_{BS}\Bigl(\frac{\partial}{\partial x}\Bigr)\langle x|C\rangle
=H_{BS}\Bigl(\frac{\partial}{\partial x}\Bigr)C(x)
\tag{4.15}
\end{equation*}
左端是纯记号——算符夹在基矢与态矢之间；右端是微积分——$H_{BS}$对波
函数逐项求导。(4.15)看似只是改写，实为授权：既然
$\langle x|H_{BS}|C\rangle$有意义，$|C\rangle$就是合格的态矢，4.6节可以
放手对它做指数演化$e^{-\tau H}$。

\paragraph{第三步：随机波动率}
Merton--Garman方程（原书式(3.34)）以股价$S$与方差$V$为变量：
\begin{equation*}
\frac{\partial C}{\partial t}
+rS\frac{\partial C}{\partial S}
+(\lambda+\mu V)\frac{\partial C}{\partial V}
+\frac{1}{2}VS^{2}\frac{\partial^{2}C}{\partial S^{2}}
+\rho\xi V^{1/2+\alpha}S\frac{\partial^{2}C}{\partial S\,\partial V}
+\xi^{2}V^{2\alpha}\frac{\partial^{2}C}{\partial V^{2}}
=rC
\end{equation*}
$S$与$V$都是正随机变量，各取对数（双变量替换）\footnote{$y$取为
$\ln\sigma^{2}$而非$\ln\sigma$，是迁就Merton--Garman方程以方差$V$为变量
的原口径；如此$(x,y)$对称地取遍实直线，正与4.2节"粒子在直线上运动"的
舞台吻合。}：
\begin{equation*}
S=e^{x},\quad -\infty<x<\infty;\qquad
\sigma^{2}=V=e^{y},\quad -\infty<y<\infty
\end{equation*}
Merton--Garman方程的$(x,y)$形式为（原书引[5,70]）：
\begin{align*}
\frac{\partial C}{\partial t}
+\Bigl(r-\frac{e^{y}}{2}\Bigr)\frac{\partial C}{\partial x}
+\Bigl(\lambda e^{-y}+\mu-\frac{\xi^{2}}{2}e^{2y(\alpha-1)}\Bigr)
\frac{\partial C}{\partial y}
+\frac{e^{y}}{2}\frac{\partial^{2}C}{\partial x^{2}}\notag\\
+\rho\xi e^{y(\alpha-1/2)}\frac{\partial^{2}C}{\partial x\,\partial y}
+\xi^{2}e^{2y(\alpha-1)}\frac{\partial^{2}C}{\partial y^{2}}
=rC
\tag{4.16}
\end{align*}
即Merton--Garman--Schr\"odinger方程：
\begin{equation*}
\frac{\partial C}{\partial t}=H_{MG}C
\tag{4.17}
\end{equation*}
由(4.16)逐项读出（各项变号搬场，$rC$并入哈密顿量）：
\begin{align*}
H_{MG}=&-\frac{e^{y}}{2}\frac{\partial^{2}}{\partial x^{2}}
-\Bigl(r-\frac{e^{y}}{2}\Bigr)\frac{\partial}{\partial x}
-\Bigl(\lambda e^{-y}+\mu-\frac{\xi^{2}}{2}e^{2y(\alpha-1)}\Bigr)
\frac{\partial}{\partial y}\\
&-\rho\xi e^{y(\alpha-1/2)}\frac{\partial^{2}}{\partial x\,\partial y}
-\frac{\xi^{2}e^{2y(\alpha-1)}}{2}\frac{\partial^{2}}{\partial y^{2}}
+r
\tag{4.18}
\end{align*}

\paragraph{解读}
$H_{MG}$是双自由度系统，且系数全是$e^{y}$的幂，强非线性——原书形容它
"按任何标准都令人生畏"。一般$\alpha$只能数值求解；$\alpha=1/2$可用偏
微分方程技术精确解出（原书文献[45]），$\alpha=1$则可用路径积分方法
（原书文献[5]）——后者正是第5章的节目。本章余下的解析工作只对
$H_{BS}$展开，但框架（态矢、演化、完备性、谱分解）全部通用，向多变量
推广"直截了当"（原书语）。另须注意(4.16)与(4.18)的互逆结构：从偏微分
方程读出哈密顿量时，所有$\partial C$项变号搬家、$rC$原样保留——这个
"变号"规则是4.8节构造$H_{V}$的手法来源。

\section{期权的定价核}

本节对应原书4.5节。核心问题：从定价公式里把"给定现值、到期如何分布"
的那个条件概率剥出来——定价核（pricing kernel）。它携带着为任何路径
无关期权定价所需的全部信息；下一节将证明，它就是哈密顿量演化算符的
一个矩阵元。

\paragraph{模型设定与记号}
就一般性而言，取随机波动率的股票与路径无关期权：到期时刻$T$，收益
函数$g(x,y)$，求$t<T$时刻的价格。记$p(x,y,T-t;x',y')$为（风险中性的）
条件概率：给定$t$时刻的价格$x$与波动率$y$，$T$时刻取值$x'$与$y'$的
概率密度。"尚未演化"就是终值条件：
\begin{equation*}
p(x,y,0;x',y')=\delta(x'-x)\,\delta(y'-y)
\tag{4.19}
\end{equation*}
定价公式由费曼--卡茨公式（Feynman--Kac formula，原书引[31]）给出
（$\tau=T-t$）：
\begin{equation*}
C(\tau;x,y)=\int_{-\infty}^{+\infty}dx'\,dy'\,
p(x,y,\tau;x',y')\,g(x';y')
\tag{4.20}
\end{equation*}
$p(x,y,\tau;x',y')$之所以称为定价核：它是把收益函数$g(x',y')$向过去
演化回$t$时刻的积分变换的\emph{核}——"核"字即由此而来。由(4.19)立即
得到定价必须满足的终值条件：
\begin{equation*}
C(0,x,y)=g(x,y)
\tag{4.21}
\end{equation*}

\paragraph{关键公式：从双变量到单变量}
股票类期权的收益通常只依赖股价：$g(x',y')=g(x')$，与到期波动率无关。
此时末态波动率$y'$可以积掉：
\begin{equation*}
C(\tau;x,y)=\int_{-\infty}^{+\infty}dx'\,
p_{MG}(x,y,\tau;x')\,g(x')
\tag{4.22}
\end{equation*}
其中Merton--Garman定价核定义为
\begin{equation*}
p_{MG}(x,y,\tau;x')=\int_{-\infty}^{+\infty}dy'\,
p(x,y,\tau;x',y')
\tag{4.23}
\end{equation*}
更简单的Black--Scholes情形里只有股价在演化，记号相应简化：
\begin{equation*}
C(\tau;x)=\int_{-\infty}^{+\infty}dx'\,
p_{BS}(x,\tau;x')\,g(x')
\tag{4.24}
\end{equation*}

\paragraph{数字期权：定价核的金融化身}
定价核还有一个漂亮的工具性解读。数字期权（digital option）的收益函数
只在股价终值落入$x_0$附近的窄条时非零：
\begin{equation*}
D(x-x_{0})=\lim_{\epsilon\to0}
\begin{cases}
\frac{1}{\epsilon}, & -\epsilon/2\leq(x-x_{0})\leq\epsilon/2,\\
0, & \text{其他},
\end{cases}
=\delta(x-x_{0})
\end{equation*}
把$g=\delta$代入(4.24)：
\begin{equation*}
C(t;x)=\int_{-\infty}^{+\infty}dx'\,
p_{BS}(x,\tau;x')\,\delta(x'-x_{0})
=p_{BS}(x,\tau;x_{0})
\tag{4.26}
\end{equation*}

\paragraph{解读}
(4.26)说明定价核本身就是数字期权的价格：$p_{BS}(x,\tau;x_{0})$是"到期
钉住$x_{0}$邻域一格"的数字期权的价格，一般的定价核则由连续排布的数字
期权拼装而成。回头看4.2节，$\delta(x'-x_{0})$恰是基矢$|x_{0}\rangle$的
表示——所以(4.26)已经在暗示$p_{BS}$与$\langle x|\cdots|x_{0}\rangle$型
矩阵元相等，那正是下一节式(4.30)的内容。也不妨与标准教材对读：金融
教科书把(4.24)读作"风险中性期望乘贴现因子"$e^{-r\tau}E^{Q}[g(S_{T})]$，
定价核就是其中的（贴现后的）风险中性密度；本书口径里贴现因子已并入核
本身（见(4.39)的显式式），两种读法给出同一个数——但"核"的读法在漂移、
势与随机波动率在场时依然成立，"期望乘贴现"的直观却只对常数$r$好用。

\section{定价核的本征函数解}

本节对应原书4.6节。核心问题：定价核由哈密顿量决定——把它写成演化
算符$e^{-\tau H}$的矩阵元，再用$H$的本征函数把它显式展开。本节只做
Black--Scholes情形，向多变量的推广"直截了当"。

\paragraph{形式解与方向反转}
把(4.12)中的$H_{BS}$换成任意哈密顿量$H$：
\begin{equation*}
\frac{\partial C}{\partial t}=HC
\tag{4.27}
\end{equation*}
形式解是指数算符：
\begin{equation*}
C(t,x)=e^{tH}C(0,x)
\tag{4.28}
\end{equation*}
在Dirac记号下把时间依赖写明，就是三步连锁：
\begin{equation*}
\Bigl\langle x\Bigl|\frac{\partial}{\partial t}\Bigr|C,t\Bigr\rangle
=\langle x|H|C,t\rangle
\quad\Longrightarrow\quad
\frac{\partial}{\partial t}|C,t\rangle=H|C,t\rangle
\quad\Longrightarrow\quad
|C,t\rangle=e^{tH}|C,0\rangle
\end{equation*}
关键一步在于\emph{方向反转}：期权已知条件不是初值而是终值（到期收益）。
把终值条件(4.21)写成$|C,T\rangle=e^{TH}|C,0\rangle=|g\rangle$，解出
\begin{equation*}
|C,0\rangle=e^{-TH}|g\rangle
\qquad\Longrightarrow\qquad
|C,t\rangle=e^{-(T-t)H}|g\rangle
\end{equation*}
剩余时间$\tau=T-t$是\emph{倒着走}的：$\tau=0$对应$t=T$（拿到收益），
$\tau=T$对应$t=0$（今天）。于是
\begin{equation*}
C(t,x)=\langle x|C,t\rangle=\langle x|e^{-\tau H}|g\rangle
\tag{4.29}
\end{equation*}
插入完备性方程(4.2)把$|g\rangle$展开：
\begin{equation*}
C(t,x)=\int_{-\infty}^{\infty}dx'\,
\langle x|e^{-\tau H}|x'\rangle\,g(x')
\end{equation*}
定价核由此获得哈密顿量表示：
\begin{equation*}
p(x,\tau;x')=\langle x|e^{-\tau H}|x'\rangle
\tag{4.30}
\end{equation*}

\paragraph{解读：哈密顿量是贴现算符}
式(4.28)里$C$带着增长指数$e^{tH}$，看似不稳定；换用剩余时间$\tau$后
立即变成衰减指数$e^{-\tau H}$（式(4.30)）。原因不在数学而在边界条件的
位置：$C$的条件给在终端$T$，改用$\tau$表达，增长即翻转成衰减。换句话说，
哈密顿量在期权定价中的角色\emph{不是}像量子力学那样把系统推向未来，
而是把未来收益函数向过去倒演，同时完成贴现（discount）——这是本章最
重要的观念转折，"定价核=贴现后的转移概率"在这里获得算符身份。

\paragraph{谱表示}
把本征完备性(4.9)插进(4.30)的矩阵元中间：
\begin{align*}
p(x,\tau;x')
&=\langle x|e^{-\tau H}|x'\rangle\\
&=\Bigl\langle x\Bigl|e^{-\tau H}
\int_{\mathcal{D}}dE\,\mu(E)\,|\psi_{E}\rangle\langle\tilde{\psi}_{E}|
\Bigr|x'\Bigr\rangle\\
&=\int_{\mathcal{D}}dE\,\mu(E)\,e^{-\tau E}\,
\psi_{E}(x)\,\tilde{\psi}_{E}(x')
\tag{4.31}
\end{align*}
中间一步用了$e^{-\tau H}|\psi_{E}\rangle=e^{-\tau E}|\psi_{E}\rangle$——
本征函数的价值正在于此。只要能求出$H$的全部本征函数与态密度，定价核
就有显式谱表示：一篮子衰减指数$e^{-\tau E}$，本征值$E$扮演"利率"。
这个表示对研究定价核的形式性质、以及对能显式对角的哈密顿量（4.6.1
小节与原书附录4.13的障碍期权）是利器；对一般非线性哈密顿量（如(4.18)
的$H_{MG}$），精确对角化通常做不到，矩阵元$e^{-\tau H}$也就算不出来
——费曼路径积分正是为这类哈密顿量准备的解析与数值工具（第5章）。

\subsection{Black--Scholes定价核}

本小节对应原书4.6.1，任务是给常数波动率欧式期权把定价核算到底。起点
是费曼--卡茨公式
\begin{equation*}
C(t,x)=\int_{-\infty}^{\infty}dx'\,
\langle x|e^{-\tau H_{BS}}|x'\rangle\,g(x'),
\end{equation*}
其中$H_{BS}$即式(4.13)。定价核定义为
\begin{equation*}
p_{BS}(x,\tau|x')=\langle x|e^{-\tau H_{BS}}|x'\rangle
;\qquad \tau=T-t
\tag{4.32}
\end{equation*}

\paragraph{推导主线：动量基}
求$H_{BS}$本征函数的捷径是换到"动量"基——$H_{BS}$在其中是对角的。
对$|x\rangle$基做傅里叶变换（原书附录式(A.11)）：
\begin{equation*}
\langle x|x'\rangle=\delta(x-x')
=\int_{-\infty}^{\infty}\frac{dp}{2\pi}\,e^{ip(x-x')}
=\int_{-\infty}^{\infty}\frac{dp}{2\pi}\,
\langle x|p\rangle\langle p|x'\rangle
\end{equation*}
这给出动量基的完备性方程：
\begin{equation*}
\int_{-\infty}^{\infty}\frac{dp}{2\pi}\,|p\rangle\langle p|=\mathcal{I}
\tag{4.34}
\end{equation*}
以及标量积
\begin{equation*}
\langle x|p\rangle=e^{ipx};\qquad
\langle p|x\rangle=e^{-ipx}\;\cdot
\tag{4.35}
\end{equation*}
从(4.13)与(4.15)的定义直接计算矩阵元：
\begin{align*}
\langle x|H_{BS}|p\rangle
&\equiv H_{BS}\,\langle x|p\rangle=H_{BS}\,e^{ipx}\\
&=\Bigl\{\frac{1}{2}\sigma^{2}p^{2}
+i\Bigl(\frac{1}{2}\sigma^{2}-r\Bigr)p+r\Bigr\}e^{ipx}
\tag{4.36}
\end{align*}
读出结论：平面波$e^{ipx}$是$H_{BS}$的本征函数，以"动量"$p$标记，复
本征值$E(p)=\frac{1}{2}\sigma^{2}p^{2}+i\bigl(\frac{1}{2}\sigma^{2}-r\bigr)p+r$
——虚部来自漂移项，是非厄米性的直接痕迹；(4.34)保证这批本征函数完备。

左矢一侧的矩阵元供对照（原书无编号，用(4.14)）：
\begin{equation*}
\langle p|H_{BS}|x\rangle
=\langle x|H_{BS}^{\dagger}|p\rangle^{*}
=\bigl[H_{BS}^{\dagger}\,e^{ipx}\bigr]^{*}
=\Bigl\{\frac{1}{2}\sigma^{2}p^{2}
+i\Bigl(\frac{1}{2}\sigma^{2}-r\Bigr)p+r\Bigr\}e^{-ipx}
\end{equation*}
\footnote{原书脚注5的警示值得记录：有人想直接对$|x\rangle$求导来算
$\langle p|H_{BS}|x\rangle$——那是错的，因为
$\langle p|\partial/\partial x|x\rangle\neq
\partial/\partial x\,\langle p|x\rangle$。算符由其对\emph{对偶}坐标基
$\langle x|$的作用定义，而非对$|x\rangle$。厄米哈密顿量下两种做法同
答案，量子力学教科书因此对此保持沉默；非厄米情形必须分清——对
$|x\rangle$直接用$H_{BS}$作用，漂移项恰好差一个符号。4.3节"算符作用在
左边"的地雷在此引爆。}

\paragraph{谱展开与高斯积分}
在(4.32)中插入动量完备性(4.34)，再用(4.36)的本征值：
\begin{align*}
p_{BS}(x,\tau;x')
&=\langle x|e^{-\tau H_{BS}}|x'\rangle\\
&=\int_{-\infty}^{\infty}\frac{dp}{2\pi}\,
\langle x|e^{-\tau H_{BS}}|p\rangle\langle p|x'\rangle
\tag{4.37}\\
&=e^{-r\tau}\int_{-\infty}^{\infty}\frac{dp}{2\pi}\,
e^{-\frac{1}{2}\tau\sigma^{2}p^{2}}\,
e^{ip\left(x-x'+\tau\left(r-\sigma^{2}/2\right)\right)}
\tag{4.38}
\end{align*}
末行来自合并指数：$\langle x|e^{-\tau H_{BS}}|p\rangle
=e^{-\tau E(p)}e^{ipx}$，与$e^{-ipx'}$相乘后，虚部
$-i\tau\bigl(\frac{1}{2}\sigma^{2}-r\bigr)p$并入相位
$ip\bigl(x-x'+\tau(r-\sigma^{2}/2)\bigr)$。执行高斯积分（步骤从略），得
Black--Scholes定价核的显式形式：
\begin{align*}
p_{BS}(x,\tau;x')&\equiv p_{BS}(x,\tau;x';\sigma)
=\langle x|e^{-\tau H_{BS}}|x'\rangle\\
&=e^{-r\tau}\,\frac{1}{\sqrt{2\pi\tau\sigma^{2}}}\,
e^{-\frac{1}{2\tau\sigma^{2}}
\left\{x-x'+\tau\left(r-\sigma^{2}/2\right)\right\}^{2}}
\tag{4.39}
\end{align*}

\paragraph{解读}
(4.39)陈述：$x'$服从均值为$\log(S(t))+(r-\sigma^{2}/2)\tau$、方差为
$\sigma^{2}\tau$的正态分布——这正是常数波动率下应有的结果，与第3章用
鞅条件推出的定价核（原书式(3.21)）是同一对象\footnote{原书脚注6提醒
记号对应：$x'=\log(S(T))$，$x=\log(S(t))$，$\tau=T-t$；原书文献[7]
另给出拉普拉斯变换的推导路线。}。至此，哈密顿量方法与传统方法在常数
波动率这一试金石上会师，而且多交出一部可推广的机器：4.7节把同一套
语言用在鞅条件上，4.8节再给哈密顿量加势。

\section{鞅条件的哈密顿量表述}

本节对应原书4.7节。核心问题：风险中性定价的基石——鞅条件（martingale
condition）——在哈密顿量语言里是什么？答案简洁得像标语：
\begin{equation*}
H|S\rangle=0,
\end{equation*}
标的证券是哈密顿量的\emph{零能量本征态}。

\paragraph{关键公式}
出发点是(4.29)与(4.30)的合并重抄：收益$g(x)$的期权在$t<T$的价格为
\begin{equation*}
C(t,x)=\int_{-\infty}^{\infty}dx'\,
\langle x|e^{-(T-t)H}|x'\rangle\,g(x')
\tag{4.40}
\end{equation*}
按无风险（风险中性）演化，未来任一时刻$t_{*}$的股价经贴现后的期望必须
等于现价——这就是鞅条件（原书引其附录式(A.40)；显式形式即第3章的
式(3.20)）：
\begin{equation*}
S(t)=E\Bigl[e^{-(t_{*}-t)r}\,S(t_{*})\,\Big|\,S(t)\Bigr]
\tag{4.41}
\end{equation*}
$t_{*}$取任何值都必须成立，故鞅条件是瞬时条件。

\paragraph{推导主线}
在(4.40)里取收益$g(x)=S(x)$——把"期权"换成股票本身：向后演化必须还原
出$S$，这正是(4.41)的哈密顿量版本：
\begin{equation*}
S(x)=\int_{-\infty}^{\infty}dx'\,
\langle x|e^{-(t_{*}-t)H}|x'\rangle\,S(x')
\end{equation*}
写成Dirac记号：
\begin{equation*}
\langle x|S\rangle=\int_{-\infty}^{\infty}dx'\,
\langle x|e^{-(t_{*}-t)H}|x'\rangle\,\langle x'|S\rangle
\tag{4.42}
\end{equation*}
插入单证券完备性
$\mathcal{I}=\int_{-\infty}^{\infty}dx'\,|x'\rangle\langle x'|$
（即(4.2)），右端恰是同一个演化算符作用在$|S\rangle$上，于是得到
（本征态形式的）方程
\begin{equation*}
|S\rangle=e^{-(t_{*}-t)H}|S\rangle
\end{equation*}
$t_{*}$任意，取无穷小间隔即得鞅条件的瞬时表述（原书引[3]）：
\begin{equation*}
H|S\rangle=0
\tag{4.43}
\end{equation*}

\paragraph{解读}
$S=e^{x}$是$H$的\emph{零能量}本征态，这件事有两层意味。其一，股权是态
空间中不可归一化的元素——$e^{x}$在$L^{2}$意义下长度无穷，呼应4.3节
原书脚注2的提醒。其二，零本征能量意味着：在鞅演化驱动下，标的证券
"不变"——它随机漂移的部分恰好被贴现抵消，平均位置原地不动。

验证与意义同样重要。容易核验$H_{BS}$（式(4.13)）满足(4.43)：三项作用在
$e^{x}$上相消，
\begin{equation*}
-\frac{\sigma^{2}}{2}e^{x}
+\Bigl(\frac{\sigma^{2}}{2}-r\Bigr)e^{x}+r\,e^{x}=0;
\end{equation*}
$H_{MG}$（式(4.18)）同样满足：其中所有含波动率的项都带着
$\partial/\partial y$，而$S=e^{x}$不依赖$y$，自动湮灭，剩下的$x$部分与
$H_{BS}$同构，同样湮灭。这个结果的意义远大于核验本身：鞅测度的存在
等价于存在湮灭标的证券的无风险哈密顿量；而存在无风险哈密顿量，就意味
着用它定价的一切衍生品都无套利。无套利原理被压缩成一条算符方程
——这是哈密顿量形式交给金融理论的第一份结构性红利，也是下一节给哈密
顿量"加势"时唯一必须守住的约束。

\section{期权定价中的势}

本节对应原书4.8节。核心问题：量子力学哈密顿量的标准形象是"动能加势"，
\begin{equation*}
H=-\frac{1}{2m}\frac{\partial^{2}}{\partial x^{2}}+V(x)
\end{equation*}
（原书引其式(4.60)），而Black--Scholes与Merton--Garman哈密顿量
（(4.13)与(4.18)）的每一项都是对$x$、$y$的导数项——依赖"速度"，没有
乘性的$V(x)$。问：期权定价的哈密顿量能不能加势？原书的答案：能，而且
加势正是定义新金融工具、建模路径依赖期权（path-dependent option）的
机制（Linetsky，原书文献[67]）；一部分障碍期权（barrier option）就可以
表达成势问题。

\paragraph{机制：把路径依赖集中成势}
路径依赖期权仍然由Black--Scholes哈密顿量$H_{BS}$演化；把障碍等效应用
一个\emph{附加}在$H_{BS}$上的势$V(x)$来集中实现，得到有效哈密顿量
\begin{equation*}
H_{\mathrm{eff}}=H_{BS}+V
\end{equation*}
从第5章的拉格朗日量（原书式(5.21)）出发，一类路径依赖期权的哈密顿量为
\begin{equation*}
H_{\text{路径依赖期权}}=H_{BS}+igf(x)
\end{equation*}
其中函数$f(x)$编码路径依赖结构；特例是亚式期权（Asian option，原书
式(5.22)）：
\begin{equation*}
H_{\text{亚式期权}}=H_{BS}+ige^{x}
\end{equation*}
（原书文献[7]）。两式中的$g$是耦合常数，与收益函数的记号撞车，注意
区分；虚数单位$i$的存在使这类哈密顿量天然非厄米——势不必是实数，
这是金融哈密顿量与量子力学品味的又一处分岔。

\paragraph{关键公式：鞅条件约束下的势}
4.7节的(4.43)给哈密顿量立了规矩——必须湮灭$S(t)$；能加的势因此不能
乱加。满足鞅条件的Black--Scholes哈密顿量的直接推广是：
\begin{equation*}
H_{V}=-\frac{\sigma^{2}}{2}\frac{\partial^{2}}{\partial x^{2}}
+\Bigl(\frac{1}{2}\sigma^{2}-V(x)\Bigr)\frac{\partial}{\partial x}
+V(x)
\tag{4.44}
\end{equation*}
其中势$V(x)$是$x$的任意函数。核验只需一行：
\begin{equation*}
H_{V}\,e^{x}
=\Bigl[-\frac{\sigma^{2}}{2}+\Bigl(\frac{\sigma^{2}}{2}-V\Bigr)+V\Bigr]
e^{x}=0,
\end{equation*}
对任意$V$成立——漂移项里的$-V$与势里的$+V$严丝合缝地相消。(4.44)把
常数$r$替换成$V(x)$时"两处同步"（漂移与势），原因正在于此：单换一处
就会破坏鞅条件。

\paragraph{金融解读}
$V(x)$的金融身份是\emph{依赖股价路径的贴现率}（原书文献[7]）：证券的
贴现因子不再是$e^{-r(T-t)}$，而是$\exp\bigl\{-\int_{t}^{T}V(x(t'))\,dt'\bigr\}$
——随股价所走的路径而变。以
到期$T$、执行价格$K$的欧式期权为例，$t<T$时刻的价格即
\begin{equation*}
E_{[t,T]}\Bigl[e^{-\int_{t}^{T}V(x(t'))\,dt'}\,(e^{x}-K)_{+}\Bigr]
\tag{4.45}
\end{equation*}
其中演化由$-(\sigma^{2}/2)\partial^{2}/\partial x^{2}
+[(1/2)\sigma^{2}-V(x)]\,\partial/\partial x$驱动，而贴现率被选得与漂移
一致，恰是为了满足鞅条件。原书随即给出应有的审慎：通常的即期利率$r$
贴现由货币市场账户的无套利论证担保；按$V(x)$贴现能否由市场实现、是否
与金融学原理相容，需要进一步研究——换言之，(4.44)是形式上自洽的推广，
其金融地位仍是开放问题。

\paragraph{推导主线：非厄米性的相似变换}
$H_{V}$的非厄米性质特别简单：任意$V$下，它可经相似变换（similarity
transformation）化成厄米哈密顿量$H_{\mathrm{eff}}$（原书文献[7]）：
\begin{equation*}
H_{V}=e^{s}\,H_{\mathrm{eff}}\,e^{-s}
\tag{4.46}
\end{equation*}
其中
\begin{equation*}
H_{\mathrm{eff}}=-\frac{\sigma^{2}}{2}\frac{\partial^{2}}{\partial x^{2}}
+\frac{1}{2}\frac{\partial V}{\partial x}
+\frac{1}{2\sigma^{2}}\Bigl(V+\frac{1}{2}\sigma^{2}\Bigr)^{2}
;\qquad
s=\frac{1}{2}x-\frac{1}{\sigma^{2}}\int_{0}^{x}dy\,V(y)
\end{equation*}
$H_{\mathrm{eff}}$厄米（(4.6)保证二阶导项自轭，后两项是实函数乘性算
符），故其本征函数构成完备基；由(4.46)，$H_{V}$也能用$H_{\mathrm{eff}}$
的本征函数对角化：
\begin{equation*}
H_{\mathrm{eff}}|\phi_{n}\rangle=E_{n}|\phi_{n}\rangle
\quad\Longrightarrow\quad
H_{V}|\psi_{n}\rangle=E_{n}|\psi_{n}\rangle,
\end{equation*}
\begin{equation*}
|\psi_{n}\rangle=e^{s}|\phi_{n}\rangle
;\qquad
\langle\tilde{\psi}_{n}|=e^{-s}\langle\phi_{n}|\neq\langle\psi_{n}|
\end{equation*}
左右本征函数依旧分离——相似变换不消灭非厄米性，只把它集中到左右矢的
$e^{\pm s}$权重上；能谱$E_{n}$则是实的（与$H_{\mathrm{eff}}$共用）。
\footnote{原书脚注7提醒：对更复杂的哈密顿量，如Merton--Garman的
$H_{MG}$，等价的厄米$H_{\mathrm{eff}}$远非显然——相似变换的运气不会
次次都有。}

\paragraph{Black--Scholes特例}
取$V(x)=r$，$H_{BS}$自身就吃进这套变换：
\begin{equation*}
H_{BS}=e^{s}\,H_{\mathrm{eff}}\,e^{-s}
=e^{\alpha x}\Bigl[-\frac{\sigma^{2}}{2}
\frac{\partial^{2}}{\partial x^{2}}+\gamma\Bigr]e^{-\alpha x}
\tag{4.47}
\end{equation*}
其中
\begin{equation*}
\gamma=\frac{1}{2\sigma^{2}}\Bigl(r+\frac{1}{2}\sigma^{2}\Bigr)^{2}
;\qquad
\alpha=\frac{1}{\sigma^{2}}\Bigl(\frac{1}{2}\sigma^{2}-r\Bigr)
\tag{4.48}
\end{equation*}
即Black--Scholes哈密顿量与"自由粒子加常数能移$\gamma$"只差一个可显式
写出的相似变换——4.4节"漂移是实的速度依赖势"的判断在此拿到代数证书。

\paragraph{解读与交棒}
本节把哈密顿量形式的伸缩性展示完毕：势可以编码路径依赖（亚式的
$ige^{x}$、障碍的$V$），而鞅条件(4.43)是唯一不可让渡的约束。厄米化的
有效Black--Scholes哈密顿量（(4.47)--(4.48)）马上派上大用场：原书4.9节
将用它求解双重势垒问题，障碍期权由此变成"给本征函数加边界条件"的
习题。本块到此收束；障碍期权、双重障碍与本征函数边界条件的全部内容，
由下一分块自原书4.9节（书页62中部）接续。

% JOIN：本块（ch04a）止于原书4.8节末——书页62上部（p080），最后一个编号
% 公式为(4.48)；下一分块（ch04b）自原书4.9节"哈密顿量与障碍期权"
% （书页62中部，p080下半）起。原书图4.1出现在4.9节（p081附近），归
% 下一分块指引。

%TAGS: 4.2,4.3,4.4,4.5,4.6,4.7,4.8,4.9,4.10,4.11,4.12,4.13,4.14,4.15,4.16,4.17,4.18,4.19,4.20,4.21,4.22,4.23,4.24,4.26,4.27,4.28,4.29,4.30,4.31,4.32,4.34,4.35,4.36,4.37,4.38,4.39,4.40,4.41,4.42,4.43,4.44,4.45,4.46,4.47,4.48（原书无(4.1)、(4.25)与(4.33)，编号断档系原书如此；已对照PNG逐页核验，p063--p079为本块主输入，p080上部为4.8节末尾）

% ch05a Path integrals and stock options 5.1-5.7 (书页78-95 / p096-113)
% 覆盖说明：原书第5章 5.1--5.7 节为书页78--95（p096--p113）；5.7 节末尾的
% 未编号公式 F' 与收尾段落在书页96（p114）顶部、5.8 总结之前，为保本节完整
% 一并收入。5.8 总结与附录5.9--5.13（书页96--116 / p114--p134）由后半文件续写。
% 编号公式核对：原书本半 tag 为 (5.1)--(5.62)，共62条，连续无断档（已对照
% PNG 逐页核验）。后半（附录）另有 (5.63)--(5.94)。
% 原书已知排印疑点：(5.47)(5.48) 印作 O(e)（应为 O(ε)，已按上下文订正）；
% 式(5.42)第一行积分测度印作 dp_y dp_y、指数印作 H(x,y,p_y,p_y)（均应为
% p_x，已订正并加脚注）；5.5 节转移振幅一式测度印作 ∫_{BS}DX（BS 下标系
% 沿用笔误）；式(5.36) 分子印作 cosh(π−|θ|)a（按恒等式读作 cosh((π−|θ|)a)）；
% 式(5.32) 第二行求和下标印作 Σ_{i=1}^∞ a_n（哑标笔误，应为 n）。
% 以上均见相应脚注或正文说明。
\chapter{路径积分与股票期权}

本章把第4章的算符语言换成一套新的数学机器。哈密顿量$\mathcal{H}$固然把
期权定价组织成了"传播算符$e^{-\tau H}$的矩阵元"，但能直接用算符方法解出
的定价问题寥寥无几：除了自由Black--Scholes情形，一般的非线性系统（例如
Merton--Garman哈密顿量$\mathcal{H}_{MG}$）用微分方程手段很难解析求解。
Feynman路径积分（path integral）提供了换一个视角的入口：不再盯着哈密顿
算符，而是研究拉格朗日量（Lagrangian）$L$与作用量（action）$S$，把"求
定价核"改写成"对所有可能路径求和"。许多算符框架下无从下手的计算，在
路径积分里不过是一串高斯积分；本章的三个可解模型——Black--Scholes核、
简谐振子、随机波动率股票——分别演示了三种通用技术：积分变量替换、简正
模式展开（傅里叶展开）与微扰展开。这三种技术是第7--10章远期利率量子场
论的全部数学预备，其中简谐振子传播子$D(t,t')$将在第7章原封不动地复活为
远期利率场的传播子。

先修内容：第4章的态空间、完备性方程与定价核$p(x,\tau;x')=\langle x\vert
e^{-\tau H}\vert x'\rangle$。本章前半（5.1--5.7节，书页78--95）建立路径
积分表述并解出三个可解模型；后半（5.8节总结与附录5.9--5.13，书页96起）
补路径积分量子力学的物理背景、金融版海森伯不确定性原理，以及任意
$\alpha$的Merton--Garman路径积分等技术细节。

\section{定价核的拉格朗日量与作用量}

本节要解决的问题：定价核是矩阵元$\langle x\vert e^{-\tau H}\vert
x'\rangle$，可$e^{-\tau H}$的矩阵元一般算不出来——哈密顿量里动能项$T$
是微分算符、势能项$V$是乘法算符，两者不对易$[T,V]\neq0$，故$e^{-\tau H}$
指数不能拆开。Feynman的观察是：虽然有限时间段$\tau$的演化算符难算，
无穷小时间步$\epsilon$的演化算符却好算——$\epsilon\to0$时不对易造成的
误差可以忽略，$e^{-\epsilon H}\simeq e^{-\epsilon T}e^{-\epsilon V}$。
把$\tau$切成$N$小段、每段用一次这个近似，再把中间所有时刻的态插入完备
性方程求和，就得到路径积分。\footnote{严格说法是Campbell--Baker公式
（原书脚注1）：对不对易算符$A$、$B$有$\exp(A)\exp(B)=\exp(A+B+[A,B]/2
+\cdots)$，$e^{-\epsilon H}\simeq e^{-\epsilon T}e^{-\epsilon V}$正是取
其领先阶。}

\paragraph{模型设定与记号}
单一证券、距到期时间$\tau=T-t$，定价核由第4章式(4.30)给出：
$p(x,\tau;x')=\langle x\vert e^{-\tau H}\vert x'\rangle$。把到期前的
$\tau$均分为$N$步，每步$\epsilon=\tau/N$；连续时间变量$x(t)$离散化为
$x_i$，$t=i\epsilon$，$0\leq i\leq N$。于是矩阵元可以写成$N$个短时间
演化算符的乘积的极限：
\begin{align*}
p(x,\tau;x')&=\lim_{N\to\infty}\langle x\vert\bigl[e^{-\epsilon H}\bigr]^{N}\vert x'\rangle\\
&=\lim_{N\to\infty}\langle x\vert e^{-\epsilon H}\cdots e^{-\epsilon H}\vert x'\rangle
\tag{5.1}
\end{align*}
多数计算最终都会取$N\to\infty$极限。位置基矢的完备性方程（即第4章式
(4.3)，本节重新编号）是
\begin{equation*}
\mathcal{I}=\int_{-\infty}^{\infty}dx\vert x\rangle\langle x\vert
\tag{5.2}
\end{equation*}

\paragraph{关键公式}
在式(5.1)中插入$(N-1)$次完备性方程(5.2)，矩阵元便按时间切片分解：
\begin{equation*}
p(x;\tau\vert x')=\Biggl(\prod_{i=1}^{N-1}\int dx_{i}\Biggr)\prod_{i=1}^{N}
\langle x_{i}\vert e^{-\epsilon H}\vert x_{i-1}\rangle
\tag{5.3}
\end{equation*}
边界条件固定路径的两端：
\begin{equation*}
x_{N}=x,\quad x_{0}=x'
\tag{5.4}
\end{equation*}
拉格朗日量由Feynman公式定义——单步矩阵元的指数上永远可以提出一个
$\epsilon L$：
\begin{equation*}
\langle x_{i}\vert e^{-\epsilon H}\vert x_{i-1}\rangle\equiv
\mathcal{N}_{i}(\epsilon)\,e^{\epsilon L(x_{i};x_{i-1};\epsilon)}
\tag{5.5}
\end{equation*}
其中$\mathcal{N}_{i}(\epsilon)$是归一因子。把式(5.5)代回式(5.3)，得
离散时间Feynman路径积分：
\begin{equation*}
p(x;\tau,x')\equiv\int\mathcal{D}X\,e^{S}
\tag{5.6}
\end{equation*}
作用量（action）是拉格朗日量沿时间格点的求和：
\begin{equation*}
S=\epsilon\sum_{i=1}^{N}L(x_{i};x_{i-1};\epsilon)
\tag{5.7}
\end{equation*}
路径积分测度（path-integration measure）由完备性插入产生：
\begin{equation*}
\int\mathcal{D}X=\mathcal{N}_{N}(\epsilon)\prod_{i=1}^{N-1}\int
\mathcal{N}_{i}(\epsilon)\,dx_{i}
\end{equation*}
（原书脚注提醒：归一因子有$N$个，而$x_i$积分只有$N-1$个——两端由边界
条件(5.4)钉住。）当$\mathcal{N}_{i}$不依赖积分变量$x_{i}$时，归一因子
可并入常数、$N\to\infty$极限存在。对边界条件$x(\tau)=x$、$x(0)=x'$，
取极限得连续时间Feynman路径积分：
\begin{align*}
p(x;\tau,x')&=\prod_{t=0}^{\tau}\int_{-\infty}^{+\infty}dx(t)\,e^{S}\\
\Rightarrow\ \langle x\vert e^{-\tau H}\vert x'\rangle&\equiv\int\mathcal{D}X\,e^{S}
\tag{5.8}
\end{align*}
其中连续时间作用量是熟悉的形式：
\begin{equation*}
S=\int L\Bigl(x,\frac{dx}{dt}\Bigr)dt
\end{equation*}

\paragraph{推导主线}
整节只有三步。第一步，把$N$步演化算符连乘（式5.1）——这是恒等式，因为
$e^{-\tau H}=(e^{-\tau H/N})^{N}$。第二步，在相邻两个短演化算符之间插入
完备性方程(5.2)：每插一次就多一个对中间位置的积分$\int dx_i$，插$N-1$
次后矩阵元变成式(5.3)——"先固定两端、中间所有时刻的位置都自由"。
第三步是定义性的：式(5.5)把每个短时间矩阵元的对数按$\epsilon$展开，
其$O(\epsilon)$系数就定义为拉格朗日量。这一步看似任意，实则是全部魔力
所在：只要能算出单步矩阵元（而这往往只是高斯积分），就自动得到整个
理论的作用量。

\paragraph{解读与联系}
从式(5.7)看，每个时刻$t=i\epsilon$系统只有一个自由度$x_i$；式(5.8)则
表明路径积分是无穷多次积分——对$[0,\tau]$中每个时刻各积一次$x_t$。
概率论里这个积分叫维纳积分（Wiener integral），即"对布朗路径求和"的
严格数学对象。多个独立变量（如下文的随机波动率，股价加波动率两个自由
度）时只需把测度换成$\mathcal{D}X\,\mathcal{D}Y$，框架不变。路径积分
量子力学的物理背景（虚时间、欧几里得演化）原书放在附录5.9，见后半。

\section{Black--Scholes拉格朗日量}

本节把5.1节的抽象方案落到第一个具体模型：风险中性演化下的股票价格，
即Black--Scholes方程描述的常数波动率情形。此时只有一个独立变量——
股价$S=e^{x}$。单步矩阵元按式(5.5)的方式写成
$\langle x_{i}\vert e^{-\epsilon H_{BS}}\vert x_{i-1}\rangle
=N_{BS}(\epsilon)\,e^{\epsilon L_{BS}}$，其中$N_{BS}(\epsilon)$为归一
常数，函数$L_{BS}$就是Black--Scholes拉格朗日量。由于Black--Scholes
定价核是少数有精确闭式的情形（第4章式(4.39)），令$\tau=\epsilon$即可
读出$L_{BS}$——这正是路径积分方法的方便之处：知道短时间核，拉格朗日
量就是它的对数。

\paragraph{模型设定与记号}
常数波动率$\sigma$、瞬时利率$r$。第4章式(4.39)的精确定价核在
$\tau=\epsilon$时给出
\begin{align*}
\langle x_{i}\vert e^{-\epsilon H_{BS}}\vert x_{i-1}\rangle
&=\frac{e^{-r\epsilon}}{\sqrt{2\pi\epsilon\sigma^{2}}}\,
e^{-\frac{1}{2\epsilon\sigma^{2}}\{x_{i}-x_{i-1}+\epsilon(r-\sigma^{2}/2)\}^{2}}\\
&=\frac{e^{-r\epsilon}}{\sqrt{2\pi\epsilon\sigma^{2}}}\,
e^{-\frac{\epsilon}{2\sigma^{2}}\{\frac{\delta x_{i}}{\epsilon}
+(r-\sigma^{2}/2)\}^{2}}\\
&=N_{BS}(\epsilon)\,e^{\epsilon L_{BS}}
\end{align*}
其中$\delta x_{i}=x_{i}-x_{i-1}$。于是
\begin{align*}
L_{BS}(i)&=-\frac{1}{2\sigma^{2}}\Bigl(\frac{\delta x_{i}}{\epsilon}
+r-\frac{\sigma^{2}}{2}\Bigr)^{2}-r
\tag{5.9}\\
S_{BS}&=\epsilon\sum_{i=1}^{N}L_{BS}(i)
\tag{5.10}\\
N_{BS}(\epsilon)&=\frac{1}{\sqrt{2\pi\epsilon\sigma^{2}}}
\end{align*}
下标$i$标记剩余时间，$N\epsilon=\tau=T-t$（$t$是真实时间），边界条件
与式(5.4)相同：
\begin{equation*}
x_{N}=x,\quad x_{0}=x'
\tag{5.11}
\end{equation*}
完备性方程给出路径积分测度：
\begin{equation*}
\int_{BS}\mathcal{D}X=\Bigl[\frac{1}{2\pi\epsilon\sigma^{2}}\Bigr]^{N/2}
\prod_{i=1}^{N-1}\int_{-\infty}^{+\infty}dx_{i}
\tag{5.12}
\end{equation*}
Black--Scholes定价核的路径积分表示为
\begin{equation*}
p_{BS}(x,\tau;x')=\int_{BS}\mathcal{D}X\,e^{S_{BS}}
\tag{5.13}
\end{equation*}

\paragraph{解读与联系}
式(5.9)的结构值得记住：拉格朗日量是"速度"$\delta x_{i}/\epsilon$的
二次型减去常数$-r$。把$\delta x_{i}/\epsilon$看成速度，$L_{BS}$正是
质量为$1/\sigma^{2}$的自由粒子拉格朗日量——波动率倒数充当质量，波动
率越大质量越小、路径越"软"；漂移项$r-\sigma^{2}/2$只是一个速度平移，
它是风险中性漂移$r$减去对数坐标的伊藤修正$\sigma^{2}/2$。此外整个
定价核多出因子$e^{-r\tau}$——贴现不进入作用量，体现"期权费是贴现后
的期望"。

\subsection{Black--Scholes路径积分}

本小节把式(5.13)一路积到底：它形式上是一个$N$重耦合高斯积分，但经
一次变量替换即可完全解耦，最终还原式(4.39)。这是全书第一个完整演算
的路径积分，值得逐步吃透。

\paragraph{关键公式}
由式(5.10)，作用量的显式形式为
\begin{equation*}
S_{BS}=-\frac{1}{2\epsilon\sigma^{2}}\sum_{i=1}^{N}
\Bigl(x_{i}-x_{i-1}+\epsilon(r-\frac{\sigma^{2}}{2})\Bigr)^{2}-r\epsilon N
\end{equation*}
做变量替换，把每一步的"速度平移"吸收进新变量：
\begin{equation*}
\begin{aligned}
\zeta_{i}&=x_{i}-x_{i-1}+\epsilon\Bigl(r-\frac{\sigma^{2}}{2}\Bigr)\\
d\zeta_{i}&=dx_{i}\ ;\quad i=1,\ldots,N
\end{aligned}
\tag{5.14}
\end{equation*}
原来的$x_{i}$只有$N-1$个积分变量，而$\zeta_{i}$却定义了$N$个独立
变量，多出的一个自由度必须用边界条件(5.11)约束掉。$x_{0}$端的边界
条件自动满足；$x_{N}$端则要求
\begin{equation*}
\begin{aligned}
x_{N}&=x_{0}-N\epsilon\Bigl(r-\frac{\sigma^{2}}{2}\Bigr)
+\sum_{i=1}^{N}\zeta_{i}\\
\Rightarrow\ m-\sum_{i=1}^{N}\zeta_{i}=0\ ;\quad
m&\equiv x-x'+\tau\Bigl(r-\frac{\sigma^{2}}{2}\Bigr)
\end{aligned}
\tag{5.15}
\end{equation*}
把这条约束以$\delta$函数形式放进路径积分，式(5.13)变成
\begin{align*}
p_{BS}(x,\tau;x')&=\langle x\vert e^{-\tau H_{BS}}\vert x'\rangle
=\int_{BS}\mathcal{D}X\,e^{S_{BS}}\\
&=e^{-r\tau}\Bigl[\frac{1}{2\pi\epsilon\sigma^{2}}\Bigr]^{N/2}
\prod_{i=1}^{N}\int_{-\infty}^{+\infty}d\zeta_{i}\,
e^{-\frac{1}{2\epsilon\sigma^{2}}\sum_{i=1}^{N}\zeta_{i}^{2}}\,
\delta\Bigl(m-\sum_{i=1}^{N}\zeta_{i}\Bigr)
\end{align*}
用$\delta$函数的傅里叶表示（附录A式(A.11)），$\zeta_{i}$积分立即
分解为彼此无关的单变量高斯积分：
\begin{align*}
p_{BS}(x,\tau;x')&=e^{-r\tau}\Bigl[\frac{1}{2\pi\epsilon\sigma^{2}}\Bigr]^{N/2}\times\\
&\quad\int_{-\infty}^{+\infty}\frac{dp}{2\pi}\prod_{i=1}^{N}\int_{-\infty}^{+\infty}
d\zeta_{i}\,e^{-\frac{1}{2\epsilon\sigma^{2}}\sum_{i=1}^{N}\zeta_{i}^{2}}\,
e^{ip[m-\sum_{i=1}^{N}\zeta_{i}]}\\
&=e^{-r\tau}\Bigl[\frac{1}{2\pi\epsilon\sigma^{2}}\Bigr]^{N/2}
\int_{-\infty}^{+\infty}\frac{dp}{2\pi}\,e^{ipm}
\prod_{i=1}^{N}\int_{-\infty}^{+\infty}d\zeta_{i}\,
e^{-\frac{1}{2\epsilon\sigma^{2}}\zeta_{i}^{2}+ip\zeta_{i}}\\
&=e^{-r\tau}\int_{-\infty}^{+\infty}\frac{dp}{2\pi}\,
e^{ipm}\,e^{-\frac{N\epsilon\sigma^{2}}{2}p^{2}}\\
&=e^{-r\tau}\,\frac{1}{\sqrt{2\pi\tau\sigma^{2}}}\,
e^{-\frac{1}{2\tau\sigma^{2}}\{x-x'+\tau(r-\sigma^{2}/2)\}^{2}}
\tag{5.16}
\end{align*}
最后一步用$N\epsilon=\tau$，正是第4章式(4.39)的Black--Scholes定价核。

\paragraph{推导主线}
三处技术要点。其一，式(5.14)的替换把二次型对角化：作用量本是相邻
$x_{i}$的耦合平方和，换成$\zeta_{i}$后成为$\sum_{i}\zeta_{i}^{2}$——
这与统计物理里把耦合振子换成简正坐标是同一招。其二，$N$个新变量对
$N-1$个旧变量的超额自由度不能含糊：式(5.15)用边界条件生成一条线性
约束，并明写成$\delta$函数放进积分，否则会丢掉一个归一约束、终式差
一个因子。其三，$\delta$函数的傅里叶表示把"变量之和固定"的约束积分
换成对共轭动量$p$的单积分，各$\zeta_{i}$随之彻底解耦；每个单变量
高斯积分贡献$\exp(-\epsilon\sigma^{2}p^{2}/2)$，$N$个连乘后剩下的
$p$积分又是一个高斯积分，收出终式。整条推导只用了高斯积分，别无
其他工具。

\subsection*{Black--Scholes速度关联函数}

定价核只用到$e^{S_{BS}}$的积分，路径积分框架真正的增量在于：它还能
回答"路径涨落"层面的问题——股价速度的平均值与关联函数（velocity
correlation function）。这些量在算符框架里对应算符期望值的计算，在
路径积分里则是对生成泛函求导。注意一个记号约定：$\epsilon$用剩余
时间$\tau=T-t=N\epsilon$表达，因此对真实时间$t$的导数换算成对剩余
时间的导数时多一个负号。

\paragraph{模型设定与记号}
路径泛函$\mathcal{B}$的期望值定义为
\begin{equation*}
\langle\mathcal{B}\rangle=\frac{1}{Z_{BS}}\int_{BS}\mathcal{D}X\,
e^{S_{BS}}\,\mathcal{B}\ ;\quad
Z_{BS}=\int_{BS}\mathcal{D}X\,e^{S_{BS}}
\tag{5.17}
\end{equation*}
因$\langle 1\rangle=1$，除以$Z_{BS}$才能把$e^{S_{BS}}/Z_{BS}$归一成
概率测度。

\paragraph{关键公式}
在式(5.16)的积分里加一项与$\zeta_{i}$耦合的外源$j_{i}$，得矩生成
配分函数：
\begin{align*}
Z_{BS}[j]&=\int_{BS}\mathcal{D}X\,e^{S_{BS}}\,e^{\sum_{i=1}^{N}j_{i}\zeta_{i}}\\
&=e^{-r\tau}\Bigl[\frac{1}{2\pi\epsilon\sigma^{2}}\Bigr]^{N/2}
\prod_{i=1}^{N}\int_{-\infty}^{+\infty}d\zeta_{i}\,
e^{-\frac{1}{2\epsilon\sigma^{2}}\sum_{i=1}^{N}\zeta_{i}^{2}}\,
\delta\Bigl(m-\sum_{i=1}^{N}\zeta_{i}\Bigr)\,
e^{\sum_{i=1}^{N}j_{i}\zeta_{i}}
\end{align*}
与式(5.16)同理做高斯积分（$\mathcal{N}$为归一常数），得
\begin{align*}
Z_{BS}[j]&=\mathcal{N}\,e^{F[j]}\\
F[j]&=\frac{1}{2}\epsilon\sigma^{2}\sum_{i=1}^{N}j_{i}^{2}
-\frac{m^{2}}{2\epsilon\sigma^{2}N}
+\frac{m}{N}\sum_{i=1}^{N}j_{i}
-\frac{\epsilon}{2N}\sigma^{2}\Bigl(\sum_{i=1}^{N}j_{i}\Bigr)^{2}
\end{align*}
其中$m=x-x'+\tau(r-\sigma^{2}/2)$见式(5.15)。对源求导即得各阶矩，
并按剩余时间约定翻译回速度：
\begin{align*}
\frac{1}{\epsilon}\langle\zeta_{n}\rangle
&=\frac{1}{\epsilon}\frac{\partial F[j]}{\partial j_{n}}\Big|_{j=0}
=\frac{m}{\epsilon N}\\
\Rightarrow\ \langle\frac{dx}{dt}\rangle
&=-\frac{x-x'}{\tau}=\frac{x'-x}{T-t}
\tag{5.18}
\end{align*}
进一步取两个时刻$t=n\epsilon$、$t'=i\epsilon$的速度关联：
\begin{align*}
\frac{1}{\epsilon^{2}}\langle\zeta_{i}\zeta_{n}\rangle
&=\frac{1}{\epsilon^{2}}\frac{\partial^{2}F[j]}{\partial j_{i}\,\partial j_{n}}\Big|_{j=0}
=\frac{1}{\epsilon}\sigma^{2}\delta_{n-i}-\frac{\sigma^{2}}{\epsilon N}
+\Bigl(\frac{m}{\epsilon N}\Bigr)^{2}\\
\Rightarrow\ \langle\frac{dx(t)}{dt}\frac{dx(t')}{dt}\rangle
&=\sigma^{2}\delta(t-t')-\frac{\sigma^{2}}{T-t}
+\Bigl(\frac{x'-x}{T-t}\Bigr)^{2}
\tag{5.19}
\end{align*}
速度对自身平均值的偏离（涨落，fluctuation）为
\begin{equation*}
\langle\Bigl[\frac{dx(t)}{dt}-\langle\frac{dx(t)}{dt}\rangle\Bigr]
\Bigl[\frac{dx(t')}{dt}-\langle\frac{dx(t)}{dt}\rangle\Bigr]\rangle
=\sigma^{2}\delta(t-t')-\frac{\sigma^{2}}{T-t}
\end{equation*}

\paragraph{解读与联系}
平均速度(5.18)不含$\sigma$——令$\sigma\to0$，随机速度退化为确定性
的"经典速度"$(x-x')/(T-t)$，这是"经典极限"在金融语境的直接对应。
速度涨落则完全由$\sigma\neq0$贡献：把式(5.19)扣除平均后的涨落项
$\sigma^{2}\delta(t-t')-\sigma^{2}/(T-t)$与布朗桥（Brownian bridge）
的协方差对照很有启发——固定两端$x(0)=x'$、$x(\tau)=x$的布朗运动，
桥内涨落正是白噪声减去一个$1/\tau$型负修正，两套语言严丝合缝。
\footnote{熟悉Shreve或Hull的读者可对照桥过程
$W_{s}-\frac{s}{\tau}W_{\tau}$的协方差$\min(s,s')-\frac{ss'}{\tau}$：
对时间求导后恰为$\sigma^{2}$倍的$\delta(t-t')-\frac{1}{\tau}$结构。}
生成泛函技术（耦合外源、求导取矩）是第8章以后计算远期利率场论各阶
矩时反复使用的标准武器，此处是其首次亮相。

\subsection*{连续极限}

离散路径积分(5.12)--(5.13)能安全取$N\to\infty$极限，原因是其测度
本质上就是平直空间$\mathbb{R}^{N}$的测度。由式(5.10)与(5.12)取
$\epsilon\to0$，作用量变成熟悉的时间积分：
\begin{equation*}
S_{BS}=\int_{0}^{\tau}dt\,L_{BS}
=-\frac{1}{2\sigma^{2}}\int_{0}^{\tau}dt\,
\Bigl(\frac{dx}{dt}+r-\frac{1}{2}\sigma^{2}\Bigr)^{2}
\end{equation*}
边界条件为$x(0)=x'$、$x(\tau)=x$。常数波动率的Black--Scholes情形
由此显出它的物理身份：一个质量为$1/\sigma^{2}$的自由量子粒子的演化。
略去无关常数后，连续时间的Black--Scholes定价核路径积分为
\begin{align*}
p_{BS}(x,\tau\vert x')&=\langle x\vert e^{-\tau H_{BS}}\vert x'\rangle
=\int_{BS}\mathcal{D}X\,e^{S_{BS}}\\
\int_{BS}\mathcal{D}X&=\prod_{t=0}^{\tau}\int_{-\infty}^{+\infty}dx(t)
\tag{5.20}
\end{align*}

\origfig{5.1}{量子粒子同时走过所有虚路径（书页98；原图在附录5.9，属本
章后半）。从$(x_{i},t_{i})$到$(x_{f},t_{f})$画出一束连接两端的路径：
路径积分即对这束路径的振幅求和；$|\psi|^{2}$给出粒子出现在某点附近
的概率。股价语境下同一张图读作：股票价格从$x'$演化到$x$，定价核是
对其间所有可能路径的加权和。}

与量子力学情形一样，式(5.20)是对全部股价变量$x(t)$的连续泛函积分，
图形上即原书图5.1所示："股票价格在演化中同时取遍所有可能的取值"，
定价核是对这束路径的求和。

\section{路径依赖期权的路径积分}

本节解决的问题是：路径积分框架如何容纳收益函数依赖整条路径的期权。
路径依赖期权（path-dependent option）大体分两类：一类像障碍期权
（barrier option），要求股价在约定期内停留在预定区间内，否则期权
作废——这对应把路径积分的定义域收窄；另一类的收益函数依赖从初始
时刻到到期日的全部股价取值（如亚式期权）——这对应在收益函数里
塞进路径的时间积分。

\paragraph{模型设定与记号}
先看第一类。以附录4.14讨论过的双重敲出障碍期权（double knock-out
barrier option）为例，股价被限制在$a\leq x(t')\leq b$（$t\leq t'\leq T$，
$t$为当前时刻、$T$为到期时刻）。路径积分域从$\mathbb{R}^{N-1}$换成
区间$[a,b]^{N-1}$，与式(5.12)对照得
\begin{equation*}
\int_{DB}\mathcal{D}X=\Bigl[\frac{1}{2\pi\epsilon\sigma^{2}}\Bigr]^{N/2}
\prod_{i=1}^{N-1}\int_{a}^{b}dx_{i}
\end{equation*}
障碍期权的定价核即
\begin{equation*}
p_{DB}(x,\tau;x')=\int_{DB}\mathcal{D}X\,e^{S_{BS}}
\end{equation*}
再看第二类。一大类路径依赖期权的收益函数可写成（$\tau=T-t$）
\begin{equation*}
g[x,K]=\max\Bigl(K,\ \frac{1}{\tau}\int_{0}^{\tau}ds\,f(x(s))\Bigr)_{+}
\end{equation*}
亚式期权（Asian option）取$f(x(s))=e^{x(s)}$，即按股价的路径平均值
结算（收益函数按原书所印形式照录）。Matacz曾用路径积分表示系统处理
过一批这类期权。

\paragraph{关键公式}
现时$t$的期权价格由Black--Scholes路径积分配上路径收益函数给出：
\begin{align*}
C(x_{0};\tau;K)&=\int_{BS}\mathcal{D}X\,e^{S_{BS}}\,g(x,K)\\
&=\int_{-\infty}^{+\infty}d\xi\,g(\xi,K)\,
\Bigl[\int_{BS}\mathcal{D}X\,
\delta\Bigl(\xi-\frac{1}{\tau}\int_{0}^{\tau}ds\,f(x(s))\Bigr)\,e^{S_{BS}}\Bigr]\\
&\qquad\text{混合边界条件：}\ x(\tau)=x_{0}\ ;\ \frac{dx(0)}{dt}=0
\end{align*}
用$\delta$函数的傅里叶表示（附录A式(A.11)）改写：
\begin{align*}
C(x_{0},\tau;K)&=\int_{-\infty}^{+\infty}\frac{d\xi\,dp}{2\pi}\,
e^{ip\xi}\,Z(p)\,g(\xi,K)\\
\Rightarrow\ Z(p)&=\int_{BS}\mathcal{D}X\,e^{S_{\mathrm{eff}}(p)}\\
\text{其中}\quad S_{\mathrm{eff}}(p)&=S_{BS}-i\frac{p}{\tau}
\int_{0}^{\tau}ds\,f(x(s))
\tag{5.21}
\end{align*}
对亚式期权$f(x)=e^{x}$，特别地有
\begin{equation*}
L_{\text{Asian option}}=L_{BS}-i\frac{p}{\tau}\,e^{x}
\tag{5.22}
\end{equation*}

\paragraph{推导主线}
两条路径、两种技术。障碍期权一路：作用量原封不动仍是$S_{BS}$，变的
只是积分域$DB$——"路径依赖进入定义域而非作用量"。这个观察同时解释
了为什么第4章用哈密顿量（本征函数展开天然吸收边界条件）处理障碍期权
更高效，其结果即式(4.54)。路径平均一路则相反：没有障碍、作用量照旧，
收益函数里的路径平均$\frac{1}{\tau}\int_{0}^{\tau}ds\,f(x(s))$先用
$\delta$函数$\delta(\xi-\text{平均})$钉成一个新变量$\xi$，再把
$\delta$函数傅里叶展开——约束被共轭变量$p$取代后，路径积分的指数上
多出一项$-i\frac{p}{\tau}\int ds\,f(x(s))$，期权价格化成一个对$\xi$、
$p$的普通二重积分乘以修正后的"有效作用量"路径积分。

\paragraph{解读与联系}
式(5.22)的读法是：路径依赖被翻译成一个"虚势"——亚式期权的拉格朗日
量等于Black--Scholes拉格朗日量减去一个正比于$e^{x}$的线性项。原
$S_{BS}$是$x$的二次型，加上$-i(p/\tau)e^{x}$后指数仍属可积的高斯
家族，故路径依赖期权在路径积分框架里往往比在偏微分方程框架里更接近
可解；这也为第7章以后"把复杂结构塞进指数、尽量保持高斯可积"的处理
范式做了预告。混合边界条件（末端钉值$x(\tau)=x_{0}$、初端自由）值得
留意：对路径初值不施加固定值，等价于在该端点取自然（Neumann）边界
条件——这正是下一节简谐振子要正面处理的技术。

\section{期权定价哈密顿量的作用量}

本节解决的问题是：给定一个带任意势$V(x)$的期权定价哈密顿量，写出它
对应的作用量——把5.1--5.2的流程抽象成通用模板，使"任意的（路径
依赖）结构都可以表示成路径积分中的一项"。

\paragraph{关键公式}
带任意势$V(x)$的定价哈密顿量（第4章的结果）为
\begin{equation*}
H_{V}=-\frac{\sigma^{2}}{2}\frac{\partial^{2}}{\partial x^{2}}
+\Bigl(\frac{1}{2}\sigma^{2}-V(x)\Bigr)\frac{\partial}{\partial x}
+V(x)
\tag{5.23}
\end{equation*}
与Black--Scholes情形完全平行的分析给出定价核的路径积分表示：
\begin{align*}
p_{V}(x,\tau\vert x')&=\langle x\vert e^{-\tau H_{V}}\vert x'\rangle\\
&=\int_{BS}\mathcal{D}X\,e^{S_{V}}
\end{align*}
其中
\begin{equation*}
S_{V}=-\int_{0}^{\tau}dt\,
\Bigl[\frac{1}{2\sigma^{2}}\Bigl(\frac{dx}{dt}-V(x)+\frac{1}{2}\sigma^{2}\Bigr)^{2}
+V(x)\Bigr]
\tag{5.24}
\end{equation*}

\paragraph{解读与联系}
对照式(5.9)（相当于$V=0$的特例）可见势$V(x)$从两处进入作用量：一是
把"速度平移"从$r-\sigma^{2}/2$换成$-V(x)+\sigma^{2}/2$，等效于给
扩散过程一个位置相关的漂移$\frac{1}{2}\sigma^{2}-V(x)$；二是显式多出
一项$-V(x)$。用费曼--卡茨（Feynman--Kac）公式的话说：$-V$既修正
生成元的漂移、又是"杀势"，这正是抛物型方程与路径期望互相翻译的
标准结构。本节的连续作用量在附录5.10（见后半）将以离散形式重现，用于
推导金融版对易关系$[\hat{x},\hat{\dot{x}}]=\sigma^{2}$；结论是势对
对易子只有$O(\epsilon)$的效应——Black--Scholes拉格朗日量足以支撑
普适结论，这也说明5.2节的特殊情形并不特殊。

\section{简谐振子的路径积分}

本节转场到量子力学：计算简谐振子（simple harmonic oscillator）的路径
积分。它是量子力学里最简单、也最有用的可解模型之一，本质上是无穷维
高斯积分的一个特例。放在期权定价的教科书里，它绝非闲笔：本节算出的
传播子$D(t,t')$（propagator）——动能算符的逆——正是第7--10章远期
利率量子场论中传播子的原型，那里的场论传播子不过是它的"多变量版"。
本节的技术主线是把路径积分的泛函积分化成一串单变量高斯积分，手法是
傅里叶（简正模式）展开——这是路径积分三大技术之二。

\paragraph{模型设定与记号}
一个量子粒子，随机位置为$x$，在势$\hat{V}(x)=\frac{1}{2}\omega^{2}x^{2}$
中运动——即恢复力正比于离原点位移的胡克定律。粒子质量$m$、弹簧常数
$\omega$，受外力$j$。从初始时刻位置$t_{i},x_{i}$到最终时刻位置
$t_{f},x_{f}$，拉格朗日量与作用量为
\begin{align*}
S&=\int_{t_{i}}^{t_{f}}dt\,L\Bigl[x,\frac{dx}{dt}\Bigr]\\
L&=-\frac{1}{2}m\Bigl(\frac{dx}{dt}\Bigr)^{2}-\frac{1}{2}m\omega^{2}x^{2}+jx
\tag{5.25}
\end{align*}
（注意欧几里得符号约定：动能与势能两项都带负号，$L$是负定的，与式
(5.9)、(5.24)的约定一致。）转移振幅为
\begin{equation*}
p(x;t_{i};x_{f},t_{f})=\langle x_{f}\vert e^{-(t_{f}-t_{i})H}\vert x_{i}\rangle
=\int\mathcal{D}X\,e^{S}
\end{equation*}
（原书此式测度沿用了$\int_{BS}$的下标，系笔误，按5.1节记号应为
$\int\mathcal{D}X$。）当初末位置$x_{i},x_{f}$也给随机性、被积分掉时，
定义生成泛函\footnote{原书脚注3：因系统有无穷多个变量（每个时刻一个
$x(t)$），故称生成泛函（generating functional）而非生成函数
（generating function）。}
\begin{equation*}
Z(t_{i},t_{f};j)=\int_{-\infty}^{+\infty}dx_{i}\,dx_{f}\,
p(x_{i},t_{i};x_{f},t_{f})
\tag{5.26}
\end{equation*}
对初末位置的积分可以换一种方式实现——改边界条件。路径$x(t)$不再
钉死两端取值，而是取
\begin{equation*}
\frac{dx(t_{i})}{dt}=0=\frac{dx(t_{f})}{dt}
\quad\text{：Neumann边界条件}
\tag{5.27}
\end{equation*}
Neumann边界条件允许对式(5.25)的作用量分部积分（边界项因式(5.27)
消失），得到二次型：
\begin{equation*}
S=-\frac{1}{2}m\int_{t_{i}}^{t_{f}}dt\,x(t)
\Bigl[-\frac{d^{2}}{dt^{2}}+\omega^{2}\Bigr]x(t)
+\int_{t_{i}}^{t_{f}}dt\,j(t)x(t)
\tag{5.28}
\end{equation*}
生成泛函即Neumann边界条件下的路径积分：
\begin{equation*}
Z(t_{i},t_{f};j)=\int_{N}\mathcal{D}X\,e^{S}
\tag{5.29}
\end{equation*}
下标$N$表示对路径积分施加Neumann边界条件。套用附录A式(A.19)--(A.20)
的无穷维高斯积分公式，定义传播子
\begin{equation*}
m\Bigl[-\frac{d^{2}}{dt^{2}}+\omega^{2}\Bigr]
D(t,t';t_{i},t_{f})=\delta(t-t')
\quad\text{：Neumann边界条件}
\end{equation*}
便有
\begin{equation*}
Z(t_{i},t_{f};j)=e^{\frac{1}{2}\int_{t_{i}}^{t_{f}}dt\,dt'\,
j(t)\,D(t,t';t_{i},t_{f})\,j(t')}
\tag{5.30}
\end{equation*}
$D(t,t';t_{i},t_{f})$就是简谐振子的传播子——它同时是关联函数
$\langle x(t)x(t')\rangle$。

\subsection*{简谐振子路径积分：傅里叶展开}

式(5.29)迄今只是形式解；本小节把它真正积出来。路径积分
$\int\mathcal{D}X\,e^{S}$跑遍所有满足Neumann边界条件(5.27)的路径
（函数）$x(t)$，而这类函数恰可用傅里叶余弦级数展开——Neumann条件
正是余弦基的特征（端点导数为零）。

\paragraph{模型设定与记号}
记$\tau\equiv t_{f}-t_{i}$，展开
\begin{align*}
x(t)&=a_{0}+\sum_{n=1}^{\infty}a_{n}
\cos\Bigl[n\pi\frac{(t-t_{i})}{\tau}\Bigr]\\
\int_{N}\mathcal{D}X&=\mathcal{N}\prod_{n=0}^{\infty}
\int_{-\infty}^{+\infty}da_{n}
\quad\text{：无穷多重积分}
\end{align*}
$\mathcal{N}$为归一常数。泛函积分从此变成可操作的无穷维积分：对每个
傅里叶系数$a_{n}$各积一次。

\paragraph{关键公式}
余弦基的正交性（$m,n\geq1$）：
\begin{align*}
\int_{t_{i}}^{t_{f}}dt\,\cos\Bigl[n\pi\frac{(t-t_{i})}{\tau}\Bigr]
\cos\Bigl[m\pi\frac{(t-t_{i})}{\tau}\Bigr]
&=\int_{t_{i}}^{t_{f}}dt\,\sin\Bigl[n\pi\frac{(t-t_{i})}{\tau}\Bigr]
\sin\Bigl[m\pi\frac{(t-t_{i})}{\tau}\Bigr]\\
&=\frac{\tau}{2}\delta_{m-n}\ ;\quad m,n\geq1
\tag{5.31}
\end{align*}
把展开代入式(5.28)，作用量在傅里叶基下对角化：
\begin{align*}
S&=-\frac{1}{2}m\omega^{2}\tau\biggl\{a_{0}^{2}
+\frac{1}{2}\sum_{n=1}^{\infty}
\Bigl[1+\Bigl(\frac{n\pi}{\omega\tau}\Bigr)^{2}\Bigr]a_{n}^{2}\biggr\}\\
&\quad+\int_{t_{i}}^{t_{f}}dt\,j(t)\biggl\{a_{0}
+\sum_{n=1}^{\infty}a_{n}\cos\Bigl[n\pi\frac{(t-t_{i})}{\tau}\Bigr]\biggr\}\\
&=-\frac{1}{2}\sum_{n=0}^{\infty}\kappa_{n}a_{n}^{2}
+\sum_{n=0}^{\infty}j_{n}a_{n}
\tag{5.32}
\end{align*}
其中
\begin{align*}
&\kappa_{0}=m\omega^{2}\tau\ ;\quad
\kappa_{n}=\frac{1}{2}m\omega^{2}\tau
\Bigl[1+\Bigl(\frac{n\pi}{\omega\tau}\Bigr)^{2}\Bigr]\ ;\quad n\geq1\\
&j_{n}=\int_{t_{i}}^{t_{f}}dt\,j(t)\cos\Bigl[n\pi\frac{(t-t_{i})}{\tau}\Bigr]\ ;
\quad n=0,1,\ldots\infty
\end{align*}
各$a_{n}$的高斯积分在作用量(5.32)中完全解耦，路径积分化成无穷多个
单变量高斯积分的乘积，逐个套用附录A式(A.16)：
\begin{align*}
Z(t_{i},t_{f};j)&=\int\mathcal{D}X\,e^{S}\\
&=\mathcal{N}\prod_{n=0}^{\infty}
\Bigl[\int_{-\infty}^{+\infty}da_{n}\,
e^{-\frac{1}{2}\kappa_{n}a_{n}^{2}+j_{n}a_{n}}\Bigr]\\
&=e^{\frac{1}{2}\sum_{n=0}^{\infty}j_{n}\frac{1}{\kappa_{n}}j_{n}}
\tag{5.33}
\end{align*}
与式(5.30)对照读出传播子的模式展开：
\begin{equation*}
D(t,t';t_{i},t_{f})=\frac{1}{m\omega^{2}\tau}\biggl\{1
+2\sum_{n=1}^{\infty}\cos\Bigl[n\pi\frac{(t-t_{i})}{\tau}\Bigr]
\Bigl[\frac{1}{1+\bigl(\frac{n\pi}{\omega\tau}\bigr)^{2}}\Bigr]
\cos\Bigl[n\pi\frac{(t'-t_{i})}{\tau}\Bigr]\biggr\}
\tag{5.34}
\end{equation*}
令$\theta=t-t_{i}>0$、$\theta'=t'-t_{i}>0$，用积化和差把分子改写：
\begin{equation*}
2\sum_{n=1}^{\infty}
\frac{\cos(n\pi\theta/\tau)\cos(n\pi\theta'/\tau)}
{1+\bigl(\frac{n\pi}{\omega\tau}\bigr)^{2}}
=\Bigl(\frac{\omega\tau}{\pi}\Bigr)^{2}\sum_{n=1}^{\infty}
\frac{\cos(n\pi(\theta+\theta')/\tau)+\cos(n\pi(\theta-\theta')/\tau)}
{\bigl(\frac{\omega\tau}{\pi}\bigr)^{2}+n^{2}}
\tag{5.35}
\end{equation*}
对整数$n$求和用恒等式\footnote{原书脚注4：式(5.36)对任意复数$a$
成立，后文将在$a$为复数的场合使用；右端只需$a^{2}$，无须指定$a^{2}$
平方根的分支。}
\begin{equation*}
\sum_{n=1}^{\infty}\frac{\cos(n\theta)}{a^{2}+n^{2}}
=\frac{\pi}{2a}\,\frac{\cosh\bigl((\pi-|\theta|)a\bigr)}{\sinh\pi a}
-\frac{1}{2a^{2}}
\tag{5.36}
\end{equation*}
（原书印作$\cosh(\pi-|\theta|)a$，按恒等式读作
$\cosh\bigl((\pi-|\theta|)a\bigr)$。）得封闭形式：
\begin{equation*}
D(t,t';t_{i},t_{f})
=\frac{\cosh\omega\{\tau-|\theta-\theta'|\}
+\cosh\omega\{\tau-(\theta+\theta')\}}
{2m\omega\sinh\omega\tau}
\tag{5.37}
\end{equation*}
代回$\theta,\theta'$并注意$\tau=t_{f}-t_{i}$，传播子为
\begin{equation*}
D(t,t';t_{i},t_{f})=
\frac{\cosh\omega\{(t_{f}-t_{i})-|t-t'|\}
+\cosh\omega\{(t_{f}-t_{i})-(t+t'-2t_{i})\}}
{2m\omega\sinh\omega(t_{f}-t_{i})}
\tag{5.38}
\end{equation*}
取无穷时间极限$t_{i}\to-\infty$、$t_{f}\to+\infty$，双曲函数中来自
边界的第二项消失，只剩平移不变部分：
\begin{equation*}
D(t,t')=\frac{1}{2m\omega}\,e^{-\omega|t-t'|}
\tag{5.39}
\end{equation*}
此式在附录A用另一方法（式(A.3.3)）亦曾导出。

\paragraph{推导主线}
五步。第一，式(5.26)对初末位置的显式积分等价于给路径换Neumann边界
条件(5.27)——"积分掉端点"="放开端点"，这是场论里"对场构形积分
自动给出自然边界条件"的最简实例。第二，分部积分把作用量写成
$x\cdot(\text{动能算符})\cdot x$的二次型(5.28)：没有Neumann条件，
分部积分的边界项下不去。第三，傅里叶展开把泛函自由度离散成可数无穷
个系数$a_{n}$，正交性(5.31)保证二次型矩阵是对角的；动能算符
$-d^{2}/dt^{2}$在余弦基下的本征值是$(n\pi/\tau)^{2}$，与$\omega^{2}$
合并即式(5.32)中的$1+(n\pi/\omega\tau)^{2}$因子——"简正模式"的
频谱一目了然。第四，每个模式做单变量高斯积分(5.33)，指数上是
$j_{n}(1/\kappa_{n})j_{n}$：源的两两耦合经"质量的倒数"传播——
传播子就是二次型系数矩阵的逆。第五，模方级数(5.34)经积化和差(5.35)
与经典求和公式(5.36)化成双曲函数的封闭形式(5.37)--(5.38)。

\paragraph{解读与联系}
这节给出的不只是教科书习题。式(5.30)的结构$Z=\exp(\frac{1}{2}jDj)$
是所有高斯型（自由）理论的总公式：知道传播子$D$就知道一切关联函数。
式(5.38)这类"有限区间+边界条件"的传播子正是第7--10章远期利率场论
传播子$D(\theta,\theta')$的直接前身后——附录A.7将用本征函数法与
格林函数法对不同边界条件重推此式（含Dirichlet情形），供后文反复调用。
对金融读者不妨这样理解$D$：它量度"今天的外力扰动$j$对日后位置$x$
的影响残留"，$\omega$越大（回复越快）残留衰减越快，指数
$e^{-\omega|t-t'|}$即记忆的时间尺度。

\section{随机波动率股票价格的拉格朗日量}

至此路径积分只处理过线性理论——作用量是随机变量的二次型，高斯积分
包打天下。作用量含三次及以上幂的非线性系统一般出了名地难解，可精确
求解者屈指可数；本节处理的股价+随机波动率（stochastic volatility）
系统恰是这样的稀有物种：它能用路径积分精确求解。系统每个时刻有两个
随机变量（两个自由度）——股价$S(t)$与其波动率$V(t)$。这个例子既因
随机波动率本身而重要，也因其复杂度恰到好处：它是演示路径积分推导
技术的样板戏，此后遇到更难的系统（如远期利率场论）时，这里的每一招
都会重演。

\paragraph{模型设定与记号}
回到第3章的Merton--Garman框架。式(3.34)的哈密顿量对任意参数$\alpha$
成立；期权价格的实证数据显示结果对$\alpha$的具体取值相当不敏感，
$\frac{1}{2}\leq\alpha\leq1$的任何值拟合得一样好。本节只取$\alpha=1$
（计算大为简化，且给出高效的数值算法），并取波动率随机方程(3.27)中
$\lambda=0$。沿用第3章记号$x=\ln S$、$y=\ln V$，两个高斯白噪声
$Q$、$R$的相关系数为$\rho$，股价与波动率的基本随机方程为
\begin{equation*}
\frac{dV}{dt}=\mu V+\xi VQ\ ;\quad
\frac{dS}{dt}=\phi S+\sigma SR
\tag{5.40}
\end{equation*}
$\alpha=1$、$\lambda=0$的Merton--Garman哈密顿量（第3章式(3.34)
改写为算符形式）：
\begin{align*}
H_{MG}=-\frac{e^{y}}{2}\frac{\partial^{2}}{\partial x^{2}}
+\Bigl(\frac{1}{2}e^{y}-r\Bigr)\frac{\partial}{\partial x}
-\xi\rho e^{y/2}\frac{\partial^{2}}{\partial x\,\partial y}
-\frac{\xi^{2}}{2}\frac{\partial^{2}}{\partial y^{2}}
+\Bigl(\frac{1}{2}\xi^{2}-\mu\Bigr)\frac{\partial}{\partial y}
\tag{5.41}
\end{align*}
注意$x$方向与$y$方向通过交叉导数项$-\xi\rho e^{y/2}\partial_{x}\partial_{y}$
耦合，耦合强度由相关系数$\rho$控制；$e^{y}$取代了常数$\sigma^{2}$
的位置——波动率进入"动能"。

\subsection*{Merton--Garman拉格朗日量的推导}

本小节从$\mathcal{H}_{MG}$出发导出拉格朗日量$L_{MG}$（为简明，
推导中省去哈密顿量的下标$MG$）。困难在于$x$与$y$两个自由度的算符
都不对易于坐标，须借用第4章4.6.1节引入的动量基：把短时间矩阵元
先在动量表象中算出（那里哈密顿量是普通函数），再做二维高斯积分。

\paragraph{模型设定与记号}
由动量基完备性（第4章式(4.34)）：
\begin{align*}
\langle x,y\vert e^{-\epsilon H}\vert x',y'\rangle
&=\int_{-\infty}^{\infty}\frac{dp_{x}}{2\pi}\frac{dp_{y}}{2\pi}\,
\langle x,y\vert e^{-\epsilon H}\vert p_{x},p_{y}\rangle
\langle p_{x},p_{y}\vert x',y'\rangle\\
&=\int_{-\infty}^{\infty}\frac{dp_{x}}{2\pi}\frac{dp_{y}}{2\pi}\,
e^{ip_{x}(x-x')}e^{ip_{y}(y-y')}e^{-\epsilon H(x,y,p_{x},p_{y})}
\tag{5.42}
\end{align*}
（原书此式首行积分测度印作$dp_{y}\,dp_{y}$、指数印作
$H(x,y,p_{y},p_{y})$，均为$p_{x}$之误，已按上下文订正。）由式(5.41)，
哈密顿量的动量表象（$p\to$乘$ip$的代换）是
\begin{equation*}
H(x,y,p_{x},p_{y})=\frac{e^{y}}{2}p_{x}^{2}
+\Bigl(\frac{1}{2}e^{y}-r\Bigr)ip_{x}
+\xi\rho e^{y/2}p_{x}p_{y}
+\frac{\xi^{2}}{2}p_{y}^{2}
+\Bigl(\frac{1}{2}\xi^{2}-\mu\Bigr)ip_{y}
\end{equation*}
写成矩阵形式：
\begin{align*}
\langle x,y\vert e^{-\epsilon H}\vert x',y'\rangle
&=\int_{-\infty}^{\infty}\frac{dp_{x}}{2\pi}\frac{dp_{y}}{2\pi}
\exp\biggl\{-\frac{\epsilon}{2}
\begin{bmatrix}p_{x} & p_{y}\end{bmatrix}\mathcal{M}
\begin{bmatrix}p_{x}\\ p_{y}\end{bmatrix}
+i\begin{bmatrix}A & B\end{bmatrix}
\begin{bmatrix}p_{x}\\ p_{y}\end{bmatrix}\biggr\}
\tag{5.43}
\end{align*}
其中
\begin{equation*}
A=x-x'+\epsilon r-\frac{\epsilon}{2}e^{y},\qquad
B=y-y'+\epsilon\mu-\frac{\epsilon}{2}\xi^{2}
\end{equation*}
而"质量矩阵"为
\begin{equation*}
\mathcal{M}=\begin{bmatrix}
e^{y} & \xi\rho e^{y/2}\\
\xi\rho e^{y/2} & \xi^{2}
\end{bmatrix}
\end{equation*}
对$p_{x}$、$p_{y}$做高斯积分所需的行列式与逆矩阵为
\begin{equation*}
\det\mathcal{M}=\xi^{2}e^{y}(1-\rho^{2}),\qquad
\mathcal{M}^{-1}=\frac{1}{\xi^{2}(1-\rho^{2})}
\begin{bmatrix}
\xi^{2}e^{-y} & -\xi\rho e^{-y/2}\\
-\xi\rho e^{-y/2} & 1
\end{bmatrix}
\end{equation*}

\paragraph{关键公式}
用附录A式(A.18)完成式(5.43)中的二维高斯积分，并把5.1节的定义式(5.5)
推广到随机波动率情形（每步写出$\epsilon L$）：
\begin{equation*}
\langle x,y\vert e^{-\epsilon H}\vert x',y'\rangle
=\frac{1}{2\pi\epsilon\sqrt{\det\mathcal{M}}}\,e^{\epsilon L}
\tag{5.44}
\end{equation*}
其中
\begin{equation*}
L_{MG}=-\frac{1}{2\epsilon^{2}(1-\rho^{2})}
\Bigl(e^{-y}A^{2}+\frac{1}{\xi^{2}}B^{2}
-2\frac{\rho}{\xi}e^{-y/2}AB\Bigr)+O(\epsilon)
\tag{5.45}
\end{equation*}
化简式(5.45)，记$\delta x=x-x'$、$\delta y=y-y'$，得$\alpha=1$的
（负定的）Merton--Garman拉格朗日量：
\begin{align*}
L_{MG}&=-\frac{1}{2\xi^{2}}\Bigl(\frac{\delta y}{\epsilon}
+\mu-\frac{1}{2}\xi^{2}\Bigr)^{2}\\
&\quad-\frac{e^{-y}}{2(1-\rho^{2})}
\biggl[\frac{\delta x}{\epsilon}+r-\frac{1}{2}e^{y}
-\frac{\rho}{\xi}e^{y/2}\Bigl(\frac{\delta y}{\epsilon}
+\mu-\frac{1}{2}\xi^{2}\Bigr)\biggr]^{2}+O(\epsilon)\\
&\equiv L_{0}+L_{X}
\tag{5.46}
\end{align*}
$\epsilon\to0$时单步矩阵元趋于双$\delta$函数（归一因子恰好保证
极限正确）：
\begin{equation*}
\lim_{\epsilon\to0}\ \langle x_{i},y_{i}\vert e^{-\epsilon H}\vert
x_{i-1},y_{i-1}\rangle
=\delta(x_{i}-x_{i-1})\,\delta(y_{i}-y_{i-1})+O(\epsilon)
\tag{5.47}
\end{equation*}
作用量为
\begin{align*}
S&=\epsilon\sum_{i=1}^{N}L(i)+O(\epsilon)\\
&\equiv S_{0}+S_{X}
\tag{5.48}
\end{align*}
$S$对$x_{i}$是二次的、对$y_{i}$是非线性的（$e^{-y}$与$e^{y/2}$在场）
——故至少股价部分的路径积分原则上可以精确做出。汇集式(5.6)与
(5.48)，Merton--Garman定价核为
\begin{align*}
p(x,y,\tau;x')&=\int_{-\infty}^{\infty}dy'\,p(x,y,\tau;x',y')\\
&=\lim_{N\to\infty}\int\mathcal{D}X\,\mathcal{D}Y\,e^{S}
\tag{5.49}
\end{align*}
其中（$\epsilon=\tau/N$）股价测度带非线性权重：
\begin{equation*}
\int\mathcal{D}X=\frac{e^{-y_{N}/2}}{\sqrt{2\pi\epsilon(1-\rho^{2})}}
\prod_{i=1}^{N-1}\int_{-\infty}^{\infty}
\frac{dx_{i}\,e^{-y_{i}/2}}{\sqrt{2\pi\epsilon(1-\rho^{2})}}
\tag{5.50}
\end{equation*}
波动率测度则为平直测度，且因式(5.49)对初值$y'$积分而多出一重
$dy_{0}=dy'$：
\begin{equation*}
\int\mathcal{D}Y=\int_{-\infty}^{\infty}dy_{0}
\prod_{i=1}^{N-1}\int_{-\infty}^{\infty}
\frac{dy_{i}}{\sqrt{2\pi\epsilon\xi^{2}}}
\tag{5.51}
\end{equation*}

\paragraph{推导主线}
全部难度集中在一次二维高斯积分上。第一步(5.42)把算符问题翻译成
动量表象里的普通积分：在$\vert p_{x},p_{y}\rangle$基下，
$\partial_{x}\to ip_{x}$、$\partial_{y}\to ip_{y}$，哈密顿量退化成
$p_{x},p_{y}$的多项式。第二步把指数组织成式(5.43)的"二次型+线性型"
标准形——注意线性项的系数$A$、$B$正是两端坐标差加漂移修正，与
Black--Scholes情形的$m$（式5.15）角色相同；二次型矩阵$\mathcal{M}$
的两个本征方向就是"股价+波动率"的联合涨落方向，$\rho$藏在非对角元。
第三步是教科书式的矩阵高斯积分：积分值由$\det\mathcal{M}$与
$\mathcal{M}^{-1}$决定（附录A式(A.18)），指数上留下
$-\frac{1}{2\epsilon}$乘"坐标差经$\mathcal{M}^{-1}$收缩"的二次型，
即式(5.45)——三个平方项的组合正是逆矩阵元的排列。第四步把$\epsilon$
吸收回$\delta x/\epsilon$、$\delta y/\epsilon$，整理成式(5.46)。

\paragraph{解读与联系}
式(5.46)的结构值得细读。$L_{0}$是波动率自身的拉格朗日量：平直空间里
的自由粒子（质量$1/\xi^{2}$，速度平移$\mu-\xi^{2}/2$）——波动率方程
(5.40)本身是几何布朗运动，对数为线性，故这部分与Black--Scholes同构。
$L_{X}$是股价部分：先是常见的$e^{-y}$权重（波动率越高、股价路径的
"质量"越小），方括号内则是股价速度减去两项修正——$-\frac{1}{2}e^{y}$
是对数坐标的伊藤修正（现在是随机的），$-\frac{\rho}{\xi}e^{y/2}
(\cdots)$是$\rho\neq0$时波动率涨落通过交叉项泄露给股价的部分；
$\rho=0$时这一项消失，$L_{X}$退化为"参数随路径变化的Black--Scholes"。
另一关键在测度(5.50)：因子$e^{-y_{i}/2}$来自$\sqrt{\det\mathcal{M}}
\propto e^{y/2}$，它使$\int\mathcal{D}X$不再是平直测度——这正是
5.7节"必须先做$\mathcal{D}X$积分、再取连续极限"的原因。求
$S(t)$与$V(t')$这类交叉关联时也须回到路径积分(5.49)。任意$\alpha$的
一般情形，原书放到章末附录5.11--5.14（见后半）。

\section{随机波动率的定价核}

本节完成任务：把式(5.49)的路径积分真正积出来。策略是"先难后易"：
先精确做出$\int\mathcal{D}X$（对$x_{i}$的高斯积分），再对剩下的
$\int\mathcal{D}Y$取连续极限。解的形式是一个无穷级数，尚余若干残差
积分，但它们全都可以精确完成——因此这个问题可谓已被形式地彻底解决。
（连续极限$N\to\infty$的严格依据由后半附录5.11的式(5.81)给出。）

\paragraph{模型设定与记号}
精确做完$\int\mathcal{D}X$后，剩下的$\int\mathcal{D}Y$的测度正如
Black--Scholes情形一样本质上是平直空间$\mathbb{R}^{N}$的测度（这与
式(5.50)的非线性形式截然不同），故对$\int\mathcal{D}Y$与
$S_{0}+S_{1}$可以放心取$N\to\infty$极限。取$\epsilon\to0$：相邻差
$\delta y_{i}/\epsilon\to dy/dt$，黎曼和$\epsilon\sum_{i=1}^{N}e^{y_{i}}
\to\int_{0}^{\tau}dt\,e^{y(t)}\equiv\tau w$。于是连续时间作用量分解为
\begin{equation*}
S=S_{0}+S_{1}
\tag{5.52}
\end{equation*}
其中
\begin{align*}
S_{0}&=-\frac{1}{2\xi^{2}}\int_{0}^{\tau}dt\,
\Bigl(\frac{dy}{dt}+\mu-\frac{1}{2}\xi^{2}\Bigr)^{2}\\
S_{1}&=-\frac{1}{2(1-\rho^{2})w}\,\Bigl\{x-x'+r\tau
-\frac{1}{2}\int_{0}^{\tau}dt\,e^{y(t)}\\
&\quad+\frac{2\rho}{\xi}\bigl(e^{y(0)/2}-e^{y/2}\bigr)
-\frac{\rho}{\xi}\Bigl(\mu-\frac{\xi^{2}}{2}\Bigr)
\int_{0}^{\tau}dt\,e^{y(t)/2}\Bigr\}^{2}
\tag{5.53}
\end{align*}
边界值为$y(\tau)=y$。定价核是只余波动率路径的积分：
\begin{equation*}
p(x,y,\tau;x')=\int\mathcal{D}Y\,
\frac{e^{S}}{\sqrt{2\pi(1-\rho^{2})\tau w}}
\tag{5.54}
\end{equation*}

\paragraph{关键公式}
$\rho\neq0$时（$\rho=0$的情形见下文脚注），由式(5.54)可见定价核除
$w$外还依赖$u=\frac{1}{\tau}\int_{0}^{\tau}e^{y(t)/2}dt$与
$e^{y(0)/2}$。引入三个"路径矩"变量
\begin{equation*}
u=\frac{1}{\tau}\int_{0}^{\tau}e^{y(t)/2}dt\ ;\quad
v=e^{y(0)/2}\ ;\quad
w=\frac{1}{\tau}\int_{0}^{\tau}e^{y(t)}dt
\end{equation*}
把定价核写成对这些矩的加权平均：
\begin{equation*}
p(x,y,\tau;x')=\int_{0}^{\infty}dw\ du\ dv\,
\frac{e^{S_{1}(u,v,w)}}{\sqrt{2\pi(1-\rho^{2})\tau w}}\,g(u,v,w)
\tag{5.55}
\end{equation*}
其中
\begin{align*}
S_{1}(u,v,w)&=-\frac{1}{2(1-\rho^{2})\tau w}\,
\Bigl[x-x'+r\tau-\frac{\tau}{2}w
+\frac{2\rho}{\xi}(v-e^{y/2})\\
&\qquad\qquad-\frac{\rho}{\xi}\Bigl(\mu-\frac{\xi^{2}}{2}\Bigr)u\Bigr]^{2}
\end{align*}
而$g(u,v,w)$是矩$(u,v,w)$的概率分布，由带三个$\delta$约束的路径
积分给出：
\begin{align*}
g(u,v,w)&=\int\mathcal{D}Y\,e^{S_{0}}\,
\delta\bigl\{v-e^{y(0)/2}\bigr\}\,
\delta\Bigl\{u-\frac{1}{\tau}\int_{0}^{\tau}e^{y(t)/2}\,dt\Bigr\}\,
\delta\Bigl\{w-\frac{1}{\tau}\int_{0}^{\tau}e^{y(t)}\,dt\Bigr\}
\tag{5.56}
\end{align*}
由式(5.56)看，$g(u,v,w)$是$u,v,w$的概率密度，定价核(5.55)是
$e^{S_{1}(u,v,w)}/\sqrt{2\pi(1-\rho^{2})\tau w}$对分布$g$的加权平均。
$g$的路径积分是非线性的，无法精确做出；但它有一个非凡性质——与
$\rho$完全无关。沿用Hull--White的思路，把式(5.55)的被积函数按
$u,v,w$展成无穷幂级数，定价核的计算便化归为求$u,v,w$的各阶矩：
\begin{align*}
\langle u^{n}w^{m}v^{p}\rangle
&\equiv\int_{0}^{\infty}du\ dw\ u^{n}\,w^{m}\,v^{p}\ g(u,v,w)\\
&=\int\mathcal{D}Y\,
\Bigl[\frac{1}{\tau}\int_{0}^{\tau}e^{y(t)/2}\,dt\Bigr]^{n}
\Bigl[\frac{1}{\tau}\int_{0}^{\tau}e^{y(t)}\,dt\Bigr]^{m}
e^{py(0)/2}\,e^{S_{0}}
\tag{5.57}
\end{align*}
这个路径积分可以精确完成。把矩写成对插入时刻的多重积分：
\begin{equation*}
\langle u^{n}w^{m}v^{p}\rangle
=\frac{1}{\tau^{n+m}}\int_{0}^{\tau}dt_{1}\cdots dt_{n}\,
dt_{n+1}\cdots dt_{n+m}\,Z(j,y,p)
\tag{5.58}
\end{equation*}
其中生成泛函
\begin{equation*}
Z(j,y,p)=\int\mathcal{D}Y\,\exp\Bigl[\int_{0}^{\tau}dt\,j(t)y(t)\Bigr]\,e^{S_{0}}
\tag{5.59}
\end{equation*}
由式(5.57)--(5.59)反推，外源$j(t)$是一串$\delta$函数的集合：
\begin{align*}
j(t)&=\frac{1}{2}\sum_{i=1}^{n}\delta(t-t_{i})
+\sum_{i=n+1}^{n+m}\delta(t-t_{i})+\frac{p}{2}\,\delta(t)\\
&\equiv\sum_{i=1}^{n+m}a_{i}\,\delta(t-t_{i})+\frac{p}{2}\,\delta(t)
\tag{5.60}
\end{align*}
$Z(j,y,p)$在附录5.12（见后半；用到其式(5.93)、(5.94)）中被精确算出，
于是
\begin{equation*}
\langle u^{n}w^{m}v^{p}\rangle
=\sigma^{p}\,\frac{e^{y\sum_{i}a_{i}}}{\tau^{n+m}}
\Biggl[\prod_{i=1}^{n+m}\int_{0}^{\tau}dt_{i}\Biggr]\,e^{F}
\qquad\text{：}\alpha=1\ \text{精确解}
\tag{5.61}
\end{equation*}
（此处$\sigma=e^{y(0)/2}=v$即初始波动率；$n+m=0$时约定空积分为1、
空求和为0。）由式(5.60)与附录5.12的结果，经整理得
\begin{equation*}
F=\Bigl(\mu-\frac{1}{2}\xi^{2}\Bigr)\sum_{i=1}^{n+m}a_{i}t_{i}
+\xi^{2}\sum_{i,j=1}^{n+m}a_{i}a_{j}\,\Theta(t_{i}-t_{j})\,t_{j}+F'
\tag{5.62}
\end{equation*}
其中
\begin{equation*}
F'=\frac{1}{8}p^{2}\xi^{2}\tau
+\frac{1}{2}p\tau\Bigl(\mu-\frac{1}{2}\xi^{2}\Bigr)
+\frac{1}{2}p\xi^{2}\sum_{i=1}^{n+m}a_{i}t_{i}
\end{equation*}
（化简中用了恒等式$\int_{0}^{\tau}dt\,\delta(t-\tau)=\frac{1}{2}$，
它可由附录A式(A.9)得到。）式(5.61)构成$\alpha=1$情形的精确解：
指数$F$对每个$t_{i}$都是线性的，故所有$t_{i}$积分都能精确完成。
股价及其随机波动率的若干具体矩在附录5.13（见后半）中算出。

\paragraph{推导主线}
顺序不可颠倒：连续极限只能在离散的$\int\mathcal{D}X$积分完成之后取，
因为式(5.50)的非线性测度（权重$e^{-y_{i}/2}$随路径变化）没有有限的
连续极限；做完$x$积分后权重被吸收进闭式结果，剩下的$y$测度是平直的，
连续极限才存在。$\int\mathcal{D}X$本身如何积出：仿照5.2.1节的变量
替换思路，把式(5.46)方括号里"速度减漂移"的组合换成新变量（其边界
条件把$x_{0}=x'$、$x_{N}=x$一并吸收），二次型对角化后逐点高斯积分；
计算细节原书放在附录5.11（见后半），式(5.53)中的三个时间积分
$\int e^{y}$、$\int e^{y/2}$与端点差$e^{y(0)/2}-e^{y/2}$正是那次
积分的遗留物。注意$S_{0}$只含$y$、恰是5.2节自由粒子作用的翻版；
全部$\rho$依赖都集中在$S_{1}$。

\paragraph{解读与联系}
本节的架构值得从高处看一遍：定价核(5.55)=$S_{1}$部分（给定矩后的
"平均Black--Scholes"因子）$\times$矩分布$g$。$S_{1}(u,v,w)$正是把
常数波动率核里的$\sigma^{2}$换成了平均波动率$w$、再补上$\rho$修正项
的产物；而$g(u,v,w)$编码"波动率路径的平均平方与平均开方各是多少"
的全部统计。三个矩各司其职：$w$（平均方差）进入分母与二次型，
$u$（平均波动率）只通过$\rho$项进入，$v$（初值波动率）体现当前
信息。$\rho=0$时退化极干净：定价核只经由$w$依赖随机波动率——
这正是第3章Merton定理与Hull--White经典结论的路径积分重述。\footnote{%
原书脚注5：$\rho=0$时令$w=\sigma^{2}$（$\sigma=e^{y/2}$为$t=0$时刻
波动率）即还原Black--Scholes公式；此时定价核对随机波动率$y(t)$的
依赖仅通过组合$w=\frac{1}{\tau}\int_{0}^{\tau}e^{y(t)}dt$，与Merton
定理（附录3.9）及Hull--White一致。}求矩技术(5.57)--(5.62)是5.2.1节
外源生成泛函法的双变量升级版：$j(t)$里的$\delta$函数串把"时间平均的
幂"翻译成"外源对路径抽样"，$S_{0}$是二次型故$Z(j,y,p)$可精确积出，
$\Theta(t_{i}-t_{j})$因子记录了波动率路径的时间有序性（因果记忆）。
整个构造在第8--10章的远期利率场论中将以"场变量+外源+传播子"的
形式再次登场——本章的三个模型正分别是它的三个侧面。

%TAGS: 5.1,5.2,5.3,5.4,5.5,5.6,5.7,5.8,5.9,5.10,5.11,5.12,5.13,5.14,5.15,5.16,5.17,5.18,5.19,5.20,5.21,5.22,5.23,5.24,5.25,5.26,5.27,5.28,5.29,5.30,5.31,5.32,5.33,5.34,5.35,5.36,5.37,5.38,5.39,5.40,5.41,5.42,5.43,5.44,5.45,5.46,5.47,5.48,5.49,5.50,5.51,5.52,5.53,5.54,5.55,5.56,5.57,5.58,5.59,5.60,5.61,5.62（5.1--5.62连续无断档；5.7节末未编号公式F'位于书页96/p114，为保本节完整一并收入；5.63起属后半附录）


\input{chapters/ch06}
\input{chapters/ch07}
\input{chapters/ch08}
\input{chapters/ch09}
\input{chapters/ch10}
\input{chapters/ch11}

\appendix
\input{chapters/chA}
\input{chapters/chB_symbols}
\input{chapters/chC_glossary}

\end{document}
